\title{导数（2016--2025 高考真题）}
\author{}
\date{}
\maketitle
\noindent 本专题共 104 道题，按年份排列。每题标注原卷与原题号。

\section*{2016 年}

\begin{question}[points=5]
\textbf{来源：}2016 年 全国III卷-理科，第 15 题\par
已知 $f(x)$ 为偶函数, 当 $x<0$ 时, $f(x)=\ln(-x)+3x$, 则曲线 $y=f(x)$ 在点 $(1,-3)$ 处的切线方程是 .
\fillin[{$2x+y+1=0$}]
\begin{solution}
\textbf{答案：}$2x+y+1=0$

\textbf{分析：}由偶函数的定义, 可得 $f(-x)=f(x)$, 即有 $x>0$ 时, $f(x)=\ln x-3x$, 求出导数, 求得切线的斜率, 由点斜式方程可得切线的方程.

\textbf{详解：}$f(x)$ 为偶函数, 可得 $f(-x)=f(x)$, 当 $x<0$ 时, $f(x)=\ln(-x)+3x$, 即有 $x>0$ 时, $f(x)=\ln x-3x$, $f'(x)=\frac{1}{x}-3$. 可得 $f(1)=\ln1-3=-3$, $f'(1)=1-3=-2$, 则曲线 $y=f(x)$ 在点 $(1,-3)$ 处的切线方程为 $y-(-3)=-2(x-1)$, 即为 $2x+y+1=0$. 故答案为 $2x+y+1=0$.
\end{solution}
\end{question}

\begin{question}[points=5]
\textbf{来源：}2016 年 全国III卷-文科，第 16 题\par
已知 $f(x)$ 为偶函数, 当 $x\le 0$ 时, $f(x)=e^{-x-1}-x$, 则曲线 $y=f(x)$ 在点 $(1,2)$ 处的切线方程是 .
\fillin[{$y=2x$}]
\begin{solution}
\textbf{答案：}$y=2x$

\textbf{分析：}由已知函数的奇偶性结合 $x\le 0$ 时的解析式求出 $x>0$ 时的解析式, 求出导函数, 得到 $f'(1)$, 然后代入直线方程的点斜式得答案.

\textbf{详解：}$f(x)$ 为偶函数, 当 $x\le 0$ 时, $f(x)=e^{-x-1}-x$. 设 $x>0$, 则 $-x<0$, $f(x)=f(-x)=e^{x-1}+x$. 则 $f'(x)=e^{x-1}+1$, $f'(1)=e^0+1=2$. 所以曲线 $y=f(x)$ 在点 $(1,2)$ 处的切线方程为 $y-2=2(x-1)$, 即 $y=2x$. 故答案为 $y=2x$.
\end{solution}
\end{question}

\begin{question}[points=12]
\textbf{来源：}2016 年 全国III卷-文科，第 21 题\par
设函数 $f(x)=\ln x-x+1$.

(1) 讨论 $f(x)$ 的单调性;

(2) 证明当 $x\in(1,+\infty)$ 时, $1<\dfrac{x-1}{\ln x}<x$;

(3) 设 $c>1$, 证明当 $x\in(0,1)$ 时, $1+(c-1)x>c^x$.
\begin{solution}
\textbf{答案：}见解析

\textbf{分析：}(1) 求出导数, 由导数大于0可得增区间, 导数小于0可得减区间, 注意函数的定义域;
(2) 由题意可得即证 $\ln x<x-1<x\ln x$. 运用(1)的单调性可得 $\ln x<x-1$, 设 $F(x)=x\ln x-x+1$, $x>1$, 求出单调性, 即可得到 $x-1<x\ln x$ 成立;
(3) 设 $G(x)=1+(c-1)x-c^x$, 求 $G(x)$ 的二次导数, 判断 $G'(x)$ 的单调性, 进而证明原不等式.

\textbf{详解：}(1) $f(x)=\ln x-x+1$, $f'(x)=\frac{1}{x}-1=\frac{1-x}{x}$, $x>0$. 由 $f'(x)>0$ 得 $0<x<1$, 由 $f'(x)<0$ 得 $x>1$. 所以 $f(x)$ 的增区间为 $(0,1)$, 减区间为 $(1,+\infty)$.
(2) 当 $x>1$ 时, $1<\frac{x-1}{\ln x}<x$ 等价于 $\ln x<x-1<x\ln x$. 由(1)知 $f(x)=\ln x-x+1$ 在 $(1,+\infty)$ 递减, 所以 $f(x)<f(1)=0$, 即 $\ln x<x-1$. 设 $F(x)=x\ln x-x+1$, $x>1$, 则 $F'(x)=\ln x+1-1=\ln x>0$, 所以 $F(x)$ 在 $(1,+\infty)$ 递增, 所以 $F(x)>F(1)=0$, 即 $x\ln x>x-1$. 故原不等式成立.
(3) 设 $G(x)=1+(c-1)x-c^x$, $x\in(0,1)$, $c>1$. 则 $G'(x)=c-1-c^x\ln c$, $G''(x)=-(\ln c)^2 c^x<0$, 所以 $G'(x)$ 在 $(0,1)$ 单调递减. $G'(0)=c-1-\ln c$, $G'(1)=c-1-c\ln c$. 由(1)知 $f(x)=\ln x-x+1$, 当 $x=c$ 时, $f(c)=\ln c-c+1<0$, 所以 $c-1-\ln c>0$, 即 $G'(0)>0$. 由(2)知 $\frac{x-1}{\ln x}<x$, 取 $x=c$ 得 $\frac{c-1}{\ln c}<c$, 即 $c-1<c\ln c$, 所以 $G'(1)=c-1-c\ln c<0$. 故存在 $t\in(0,1)$ 使得 $G'(t)=0$. 所以 $x\in(0,t)$ 时 $G'(x)>0$, $x\in(t,1)$ 时 $G'(x)<0$, 即 $G(x)$ 在 $(0,t)$ 递增, 在 $(t,1)$ 递减. 又 $G(0)=1+0-1=0$, $G(1)=1+(c-1)-c=0$, 所以 $x\in(0,1)$ 时 $G(x)>0$ 成立, 即 $1+(c-1)x>c^x$.
\end{solution}
\end{question}

\begin{question}[points=5]
\textbf{来源：}2016 年 全国II卷-理科，第 16 题\par
若直线$y=kx+b$是曲线$y=\ln x+2$的切线，也是曲线$y=\ln(x+1)$的切线，则$b=$．
\fillin[{$1-\ln2$}]
\begin{solution}
\textbf{答案：}$1-\ln2$

\textbf{分析：}先设切点，然后利用切点来寻找切线斜率的联系，以及对应的函数值，综合联立求解即可．

\textbf{详解：}设直线与两曲线切点分别为$(x_1,\ln x_1+2)$和$(x_2,\ln(x_2+1))$，则$k=\frac{1}{x_1}=\frac{1}{x_2+1}$，得$x_1=x_2+1$，又$\ln x_1+2=\ln(x_2+1)+b$，解得$x_1=\frac12$，$k=2$，$b=1-\ln2$．
\end{solution}
\end{question}

\begin{question}[points=12]
\textbf{来源：}2016 年 全国II卷-理科，第 21 题\par
（Ⅰ）讨论函数$f(x)=\frac{x-2}{x+2}e^x$的单调性，并证明当$x>0$时，$(x-2)e^x+x+2>0$；

（Ⅱ）证明：当$a\in[0,1)$时，函数$g(x)=\frac{e^x-ax-a}{x^2}(x>0)$有最小值．设$g(x)$的最小值为$h(a)$，求函数$h(a)$的值域．
\begin{solution}
\textbf{答案：}（Ⅰ）$f(x)$在$(-\infty,-2)$和$(-2,+\infty)$上单调递增；证明略；（Ⅱ）$h(a)\in(\frac{1}{2},\frac{e^2}{4}]$

\textbf{分析：}从导数作为切入点探求函数的单调性，通过函数单调性来求得函数的值域，利用复合函数的求导公式进行求导，然后逐步分析即可．

\textbf{详解：}（Ⅰ）$f'(x)=\frac{x^2+4}{(x+2)^2}e^x>0$，故$f(x)$在$(-\infty,-2)$和$(-2,+\infty)$上单调递增．当$x>0$时，$f(x)>f(0)=-1$，即$\frac{x-2}{x+2}e^x>-1$，整理得$(x-2)e^x+x+2>0$．
（Ⅱ）$g'(x)=\frac{(e^x-a)x^2-2x(e^x-ax-a)}{x^4}=\frac{xe^x-2e^x+ax+2a}{x^3}=\frac{(x-2)e^x+a(x+2)}{x^3}$，由（Ⅰ）知$\frac{(x-2)e^x}{x+2}$在$(0,+\infty)$上单调递增，值域为$(-1,+\infty)$，故存在唯一$t>0$使得$\frac{(t-2)e^t}{t+2}=-a$，且$g(x)$在$(0,t)$递减，$(t,+\infty)$递增，$h(a)=g(t)=\frac{e^t-at-a}{t^2}=\frac{e^t+\frac{t-2}{t+2}e^t(t+1)}{t^2}=\frac{e^t}{t+2}$，$t\in(0,2]$（由$a\in[0,1)$得$\frac{(t-2)e^t}{t+2}\in(-1,0]$，解得$t\in(0,2]$），$h(a)=\frac{e^t}{t+2}$在$(0,2]$上递增，值域为$(\frac12,\frac{e^2}{4}]$．
\end{solution}
\end{question}

\begin{question}[points=12]
\textbf{来源：}2016 年 全国II卷-文科，第 20 题\par
已知函数$f(x)=(x+1)\ln x-a(x-1)$。

（Ⅰ）当$a=4$时，求曲线$y=f(x)$在$(1,f(1))$处的切线方程；

（Ⅱ）若当$x\in(1,+\infty)$时，$f(x)>0$，求$a$的取值范围。
\begin{solution}
\textbf{答案：}（Ⅰ）$y=-2x+2$；（Ⅱ）$a\le2$

\textbf{分析：}（Ⅰ）$f(1)=0$，$f'(x)=\ln x+\frac{x+1}{x}-4$，$f'(1)=-2$，切线$y=-2(x-1)=-2x+2$。
（Ⅱ）$f'(x)=1+\frac{1}{x}+\ln x-a$，$f''(x)=\frac{1}{x^2}+\frac{1}{x}>0$，$f'(x)$在$(1,+\infty)$递增，$f'(x)>f'(1)=2-a$。若$a\le2$，$f'(x)\ge0$，$f(x)$递增，$f(x)>f(1)=0$；若$a>2$，存在$x_0$使$f'(x_0)=0$，$f(x)$在$(1,x_0)$递减，$f(x_0)<0$，不合。故$a\le2$。
\end{solution}
\end{question}

\begin{question}[points=12]
\textbf{来源：}2016 年 全国I卷-理科，第 21 题\par
已知函数$f(x)=(x-2)e^x+a(x-1)^2$有两个零点．

（Ⅰ）求$a$的取值范围；

（Ⅱ）设$x_1$，$x_2$是$f(x)$的两个零点，证明：$x_1+x_2<2$．
\begin{solution}
\textbf{答案：}（Ⅰ）$(0,+\infty)$；（Ⅱ）证明见解析

\textbf{分析：}（Ⅰ）由函数$f(x)=(x-2)e^x+a(x-1)^2$可得：$f'(x)=(x-1)e^x+2a(x-1)=(x-1)(e^x+2a)$，对$a$进行分类讨论，综合讨论结果，可得答案．（Ⅱ）设$x_1$，$x_2$是$f(x)$的两个零点，则$-a=\frac{(x_1-2)e^{x_1}}{(x_1-1)^2}=\frac{(x_2-2)e^{x_2}}{(x_2-1)^2}$，令$g(x)=\frac{(x-2)e^x}{(x-1)^2}$，则$g(x_1)=g(x_2)=-a$，分析$g(x)$的单调性，令$m>0$，则$g(1+m)-g(1-m)=\frac{(m-1)e^{1+m}}{m^2}-\frac{(-m-1)e^{1-m}}{m^2}=\frac{e[ (m-1)e^m - (-m-1)e^{-m}]}{m^2}$，设$h(m)= (m-1)e^m + (m+1)e^{-m}$，$m>0$，利用导数法可得$h(m)>h(0)=0$恒成立，即$g(1+m)>g(1-m)$恒成立，令$m=1-x_1>0$，可得结论．

\textbf{详解：}解：（Ⅰ）因为函数$f(x)=(x-2)e^x+a(x-1)^2$，所以$f'(x)=(x-1)e^x+2a(x-1)=(x-1)(e^x+2a)$，①若$a=0$，那么$f(x)=0\Leftrightarrow(x-2)e^x=0\Leftrightarrow x=2$，函数$f(x)$只有唯一的零点$2$，不合题意；②若$a>0$，那么$e^x+2a>0$恒成立，当$x<1$时，$f'(x)<0$，此时函数为减函数；当$x>1$时，$f'(x)>0$，此时函数为增函数；此时当$x=1$时，函数$f(x)$取极小值$-e$，由$f(2)=a>0$，可得：函数$f(x)$在$x>1$存在一个零点；当$x<1$时，$e^x<e$，$x-2<-1<0$，所以$f(x)=(x-2)e^x+a(x-1)^2>(x-2)e+a(x-1)^2=a(x-1)^2+e(x-1)-e$，令$a(x-1)^2+e(x-1)-e=0$的两根为$t_1$，$t_2$，且$t_1<t_2$，则当$x<t_1$，或$x>t_2$时，$f(x)>a(x-1)^2+e(x-1)-e>0$，故函数$f(x)$在$x<1$存在一个零点；即函数$f(x)$在$R$是存在两个零点，满足题意；③若$-\frac{e}{2}<a<0$，则$\ln(-2a)<\ln e=1$，当$x<\ln(-2a)$时，$x-1<\ln(-2a)-1<\ln e-1=0$，$e^x+2a<e^{\ln(-2a)}+2a=0$，即$f'(x)=(x-1)(e^x+2a)>0$恒成立，故$f(x)$单调递增，当$\ln(-2a)<x<1$时，$x-1<0$，$e^x+2a>e^{\ln(-2a)}+2a=0$，即$f'(x)=(x-1)(e^x+2a)<0$恒成立，故$f(x)$单调递减，当$x>1$时，$x-1>0$，$e^x+2a>e^{\ln(-2a)}+2a=0$，即$f'(x)=(x-1)(e^x+2a)>0$恒成立，故$f(x)$单调递增，故当$x=\ln(-2a)$时，函数取极大值，由$f(\ln(-2a))=[\ln(-2a)-2](-2a)+a[\ln(-2a)-1]^2=a\{[\ln(-2a)-2]^2+1\}<0$得：函数$f(x)$在$R$上至多存在一个零点，不合题意；④若$a=-\frac{e}{2}$，则$\ln(-2a)=1$，当$x<1=\ln(-2a)$时，$x-1<0$，$e^x+2a<e^{\ln(-2a)}+2a=0$，即$f'(x)=(x-1)(e^x+2a)>0$恒成立，故$f(x)$单调递增，当$x>1$时，$x-1>0$，$e^x+2a>e^{\ln(-2a)}+2a=0$，即$f'(x)=(x-1)(e^x+2a)>0$恒成立，故$f(x)$单调递增，故函数$f(x)$在$R$上单调递增，函数$f(x)$在$R$上至多存在一个零点，不合题意；⑤若$a<-\frac{e}{2}$，则$\ln(-2a)>\ln e=1$，当$x<1$时，$x-1<0$，$e^x+2a<e^{\ln(-2a)}+2a=0$，即$f'(x)=(x-1)(e^x+2a)>0$恒成立，故$f(x)$单调递增，当$1<x<\ln(-2a)$时，$x-1>0$，$e^x+2a<e^{\ln(-2a)}+2a=0$，即$f'(x)=(x-1)(e^x+2a)<0$恒成立，故$f(x)$单调递减，当$x>\ln(-2a)$时，$x-1>0$，$e^x+2a>e^{\ln(-2a)}+2a=0$，即$f'(x)=(x-1)(e^x+2a)>0$恒成立，故$f(x)$单调递增，故当$x=1$时，函数取极大值，由$f(1)=-e<0$得：函数$f(x)$在$R$上至多存在一个零点，不合题意；综上所述，$a$的取值范围为$(0,+\infty)$证明：（Ⅱ）因为$x_1$，$x_2$是$f(x)$的两个零点，所以$f(x_1)=f(x_2)=0$，且$x_1\neq1$，且$x_2\neq1$，所以$-a=\frac{(x_1-2)e^{x_1}}{(x_1-1)^2}=\frac{(x_2-2)e^{x_2}}{(x_2-1)^2}$，令$g(x)=\frac{(x-2)e^x}{(x-1)^2}$，则$g(x_1)=g(x_2)=-a$，因为$g'(x)=\frac{[(x-1)-1]e^x(x-1)-2(x-2)e^x}{(x-1)^3}=\frac{e^x[(x-1)^2- (x-1)-2(x-2)]}{(x-1)^3}=\frac{e^x[(x-1)^2-3x+5]}{(x-1)^3}$，当$x<1$时，$g'(x)<0$，$g(x)$单调递减；当$x>1$时，$g'(x)>0$，$g(x)$单调递增；设$m>0$，则$g(1+m)-g(1-m)=\frac{(m-1)e^{1+m}}{m^2}-\frac{(-m-1)e^{1-m}}{m^2}=\frac{e[(m-1)e^m+(m+1)e^{-m}]}{m^2}$，设$h(m)= (m-1)e^m+(m+1)e^{-m}$，$m>0$，则$h'(m)= e^m+(m-1)e^m + e^{-m}-(m+1)e^{-m}= me^m - me^{-m}= m(e^m - e^{-m})>0$恒成立，即$h(m)$在$(0,+\infty)$上为增函数，$h(m)>h(0)=0$恒成立，即$g(1+m)>g(1-m)$恒成立，令$m=1-x_1>0$，则$g(1+1-x_1)>g(1-1+x_1)\Leftrightarrow g(2-x_1)>g(x_1)=g(x_2)\Leftrightarrow 2-x_1>x_2$，即$x_1+x_2<2$．
\end{solution}
\end{question}

\begin{question}[points=5]
\textbf{来源：}2016 年 全国I卷-文科，第 12 题\par
若函数$f(x)=x-\frac{1}{3}\sin2x+a\sin x$在$(-\infty,+\infty)$单调递增，则$a$的取值范围是
\begin{choices}
  \item $[-1,1]$
  \item $[-1,\frac{1}{3}]$
  \item $[-\frac{1}{3},\frac{1}{3}]$
  \item $[-1,-\frac{1}{3}]$
\end{choices}
\begin{solution}
\textbf{答案：}$[-\frac{1}{3},\frac{1}{3}]$

\textbf{分析：}求出$f(x)$的导数，由题意可得$f'(x)\ge0$恒成立，设$t=\cos x(-1\le t\le1)$，即有$5-4t^2+3at\ge0$，对$t$讨论，分$t=0$，$0<t\le1$，$-1\le t<0$，分离参数，运用函数的单调性可得最值，解不等式即可得到所求范围。

\textbf{详解：}函数$f(x)=x-\frac{1}{3}\sin2x+a\sin x$的导数为$f'(x)=1-\frac{2}{3}\cos2x+a\cos x$，由题意可得$f'(x)\ge0$恒成立，即为$1-\frac{2}{3}\cos2x+a\cos x\ge0$，即有$\frac{5}{3}-\frac{4}{3}\cos^2x+a\cos x\ge0$，设$t=\cos x(-1\le t\le1)$，即有$5-4t^2+3at\ge0$，当$t=0$时，不等式显然成立；当$0<t\le1$时，$3a\ge4t-\frac{5}{t}$，由$4t-\frac{5}{t}$在$(0,1]$递增，可得$t=1$时，取得最大值$-1$，可得$3a\ge-1$，即$a\ge-\frac{1}{3}$；当$-1\le t<0$时，$3a\le4t-\frac{5}{t}$，由$4t-\frac{5}{t}$在$[-1,0)$递增，可得$t=-1$时，取得最小值$1$，可得$3a\le1$，即$a\le\frac{1}{3}$。综上可得$a$的范围是$[-\frac{1}{3},\frac{1}{3}]$。另解：设$t=\cos x(-1\le t\le1)$，即有$5-4t^2+3at\ge0$，由题意可得$5-4+3a\ge0$，且$5-4-3a\ge0$，解得$a$的范围是$[-\frac{1}{3},\frac{1}{3}]$。故选：$C$。
\end{solution}
\end{question}

\begin{question}[points=12]
\textbf{来源：}2016 年 全国I卷-文科，第 21 题\par
已知函数$f(x)=(x-2)e^x+a(x-1)^2$。

（Ⅰ）讨论$f(x)$的单调性；

（Ⅱ）若$f(x)$有两个零点，求$a$的取值范围。
\begin{solution}
\textbf{答案：}（Ⅰ）当$a\ge0$时，$f(x)$在$(-\infty,1)$递减，在$(1,+\infty)$递增；当$a=-\frac{e}{2}$时，$f(x)$在$\mathbb{R}$上递增；当$a<-\frac{e}{2}$时，$f(x)$在$(-\infty,1)$，$(\ln(-2a),+\infty)$递增，在$(1,\ln(-2a))$递减；当$-\frac{e}{2}<a<0$时，$f(x)$在$(-\infty,\ln(-2a))$，$(1,+\infty)$递增，在$(\ln(-2a),1)$递减。（Ⅱ）$(0,+\infty)$。

\textbf{分析：}（Ⅰ）求出$f(x)$的导数，讨论当$a\ge0$时，$a<-\frac{e}{2}$时，$a=-\frac{e}{2}$时，$-\frac{e}{2}<a<0$，由导数大于0，可得增区间；由导数小于0，可得减区间；
（Ⅱ）由（Ⅰ）的单调区间，对$a$讨论，结合单调性和函数值的变化特点，即可得到所求范围。

\textbf{详解：}（Ⅰ）由$f(x)=(x-2)e^x+a(x-1)^2$，可得$f'(x)=(x-1)e^x+2a(x-1)=(x-1)(e^x+2a)$，
①当$a\ge0$时，由$f'(x)>0$，可得$x>1$；由$f'(x)<0$，可得$x<1$，即有$f(x)$在$(-\infty,1)$递减；在$(1,+\infty)$递增；
②当$a<0$时，若$a=-\frac{e}{2}$，则$f'(x)\ge0$恒成立，即有$f(x)$在$R$上递增；
若$a<-\frac{e}{2}$时，由$f'(x)>0$，可得$x<1$或$x>\ln(-2a)$；由$f'(x)<0$，可得$1<x<\ln(-2a)$。即有$f(x)$在$(-\infty,1)$，$(\ln(-2a),+\infty)$递增；在$(1,\ln(-2a))$递减；
若$-\frac{e}{2}<a<0$，由$f'(x)>0$，可得$x<\ln(-2a)$或$x>1$；由$f'(x)<0$，可得$\ln(-2a)<x<1$。即有$f(x)$在$(-\infty,\ln(-2a))$，$(1,+\infty)$递增；在$(\ln(-2a),1)$递减；
（Ⅱ）①由（Ⅰ）可得当$a>0$时，$f(x)$在$(-\infty,1)$递减；在$(1,+\infty)$递增，且$f(1)=-e<0$，$x\to+\infty$，$f(x)\to+\infty$；当$x\to-\infty$时$f(x)>0$或找到一个$x<1$使得$f(x)>0$对于$a>0$恒成立，$f(x)$有两个零点；
②当$a=0$时，$f(x)=(x-2)e^x$，所以$f(x)$只有一个零点$x=2$；
③当$a<0$时，若$a<-\frac{e}{2}$时，$f(x)$在$(1,\ln(-2a))$递减，在$(-\infty,1)$，$(\ln(-2a),+\infty)$递增，又当$x\le1$时，$f(x)<0$，所以$f(x)$不存在两个零点；当$a\ge-\frac{e}{2}$时，在$(-\infty,\ln(-2a))$单调增，在$(1,+\infty)$单调增，在$(\ln(-2a),1)$单调减，只有$f(\ln(-2a))$等于0才有两个零点，而当$x\le1$时，$f(x)<0$，所以只有一个零点不符题意。
综上可得，$f(x)$有两个零点时，$a$的取值范围为$(0,+\infty)$。
\end{solution}
\end{question}

\section*{2017 年}

\begin{question}[points=5]
\textbf{来源：}2017 年 全国III卷-理科，第 11 题\par
已知函数$f(x)=x^2-2x+a(e^{x-1}+e^{-x+1})$有唯一零点，则$a=$（ ）
\begin{choices}
  \item $-\frac{1}{2}$
  \item $\frac{1}{3}$
  \item $\frac{1}{2}$
  \item $1$
\end{choices}
\begin{solution}
\textbf{答案：}$C$

\textbf{分析：}通过转化可知问题等价于函数$y=1-(x-1)^2$的图象与$y=a(e^{x-1}+e^{-x+1})$的图象只有一个交点求$a$的值。分$a=0$、$a<0$、$a>0$三种情况，结合函数的单调性分析可得结论。

\textbf{详解：}因为$f(x)=x^2-2x+a(e^{x-1}+e^{-x+1})=-1+(x-1)^2+a(e^{x-1}+e^{-x+1})=0$，所以函数$f(x)$有唯一零点等价于方程$1-(x-1)^2=a(e^{x-1}+e^{-x+1})$有唯一解，等价于函数$y=1-(x-1)^2$的图象与$y=a(e^{x-1}+e^{-x+1})$的图象只有一个交点。①当$a=0$时，$f(x)=x^2-2x\ge -1$，此时有两个零点，矛盾；②当$a<0$时，由于$y=1-(x-1)^2$在$(-\infty,1)$上递增、在$(1,+\infty)$上递减，且$y=a(e^{x-1}+e^{-x+1})$在$(-\infty,1)$上递增、在$(1,+\infty)$上递减，所以函数$y=1-(x-1)^2$的图象的最高点为$A(1,1)$，$y=a(e^{x-1}+e^{-x+1})$的图象的最高点为$B(1,2a)$，由于$2a<0<1$，此时函数$y=1-(x-1)^2$的图象与$y=a(e^{x-1}+e^{-x+1})$的图象有两个交点，矛盾；③当$a>0$时，由于$y=1-(x-1)^2$在$(-\infty,1)$上递增、在$(1,+\infty)$上递减，且$y=a(e^{x-1}+e^{-x+1})$在$(-\infty,1)$上递减、在$(1,+\infty)$上递增，所以函数$y=1-(x-1)^2$的图象的最高点为$A(1,1)$，$y=a(e^{x-1}+e^{-x+1})$的图象的最低点为$B(1,2a)$，由题可知点$A$与点$B$重合时满足条件，即$2a=1$，即$a=\frac{1}{2}$，符合条件；综上所述，$a=\frac{1}{2}$。故选C。
\end{solution}
\end{question}

\begin{question}[points=12]
\textbf{来源：}2017 年 全国III卷-理科，第 21 题\par
已知函数$f(x)=x-1-a\ln x$．

（1）若$f(x)\ge 0$，求$a$的值；

（2）设$m$为整数，且对于任意正整数$n$，$(1+\frac{1}{2})(1+\frac{1}{2^2})\cdots(1+\frac{1}{2^n})<m$，求$m$的最小值．
\begin{solution}
\textbf{答案：}（1）$a=1$；（2）$m$的最小值为$3$

\textbf{分析：}（1）通过对函数$f(x)=x-1-a\ln x(x>0)$求导，分$a\le 0$、$a>0$两种情况考虑导函数$f'(x)$与$0$的大小关系可得结论；
（2）通过（1）可知$\ln x\le x-1$，进而取特殊值可知$\ln(1+\frac{1}{2^k})<\frac{1}{2^k}$，$k\in N^*$．一方面利用等比数列的求和公式放缩可知$(1+\frac{1}{2})(1+\frac{1}{2^2})\cdots(1+\frac{1}{2^n})<e$，另一方面可知$(1+\frac{1}{2})(1+\frac{1}{2^2})\cdots(1+\frac{1}{2^n})>2$，从而当$n\ge 3$时，$(1+\frac{1}{2})(1+\frac{1}{2^2})\cdots(1+\frac{1}{2^n})\in(2,e)$，比较可得结论。

\textbf{详解：}解：（1）因为函数$f(x)=x-1-a\ln x$，$x>0$，所以$f'(x)=1-\frac{a}{x}=\frac{x-a}{x}$，且$f(1)=0$．所以当$a\le 0$时$f'(x)>0$恒成立，此时$y=f(x)$在$(0,+\infty)$上单调递增，这与$f(x)\ge 0$矛盾；当$a>0$时令$f'(x)=0$，解得$x=a$，所以$y=f(x)$在$(0,a)$上单调递减，在$(a,+\infty)$上单调递增，即$f(x)_{\min}=f(a)$，若$a\neq 1$，则$f(a)<f(1)=0$，从而与$f(x)\ge 0$矛盾；所以$a=1$；
（2）由（1）可知当$a=1$时$f(x)=x-1-\ln x\ge 0$，即$\ln x\le x-1$，所以$\ln(x+1)\le x$当且仅当$x=0$时取等号，所以$\ln(1+\frac{1}{2^k})<\frac{1}{2^k}$，$k\in N^*$．$\ln(1+\frac{1}{2})+\ln(1+\frac{1}{2^2})+\cdots+\ln(1+\frac{1}{2^n})<\frac{1}{2}+\frac{1}{2^2}+\cdots+\frac{1}{2^n}=1-\frac{1}{2^n}<1$，即$(1+\frac{1}{2})(1+\frac{1}{2^2})\cdots(1+\frac{1}{2^n})<e$；因为$m$为整数，且对于任意正整数$n$，$(1+\frac{1}{2})(1+\frac{1}{2^2})\cdots(1+\frac{1}{2^n})<m$成立，当$n=3$时，不等式左边大于$2$，所以$m$的最小值为$3$。
\end{solution}
\end{question}

\begin{question}[points=12]
\textbf{来源：}2017 年 全国III卷-文科，第 21 题\par
已知函数$f(x)=\ln x+ax^2+(2a+1)x$。

（1）讨论$f(x)$的单调性；

（2）当$a<0$时，证明$f(x)\le -\frac{1}{4a}-2$。
\begin{solution}
\textbf{答案：}（1）当$a\ge 0$时，$f(x)$在$(0,+\infty)$上单调递增；当$a<0$时，$f(x)$在$(0,-\frac{1}{2a})$上单调递增，在$(-\frac{1}{2a},+\infty)$上单调递减；（2）证明略

\textbf{分析：}（1）求导$f'(x)=\frac{2ax+1}{x}+2a+1=\frac{(2ax+1)(x+1)}{x}$，分类讨论即可。（2）利用（1）中最大值，构造函数证明。

\textbf{详解：}（1）$f'(x)=\frac{1}{x}+2ax+2a+1=\frac{(2ax+1)(x+1)}{x}\ (x>0)$。当$a\ge 0$时，$f'(x)>0$恒成立，$f(x)$在$(0,+\infty)$递增；当$a<0$时，令$f'(x)=0$得$x=-\frac{1}{2a}$，在$(0,-\frac{1}{2a})$上$f'(x)>0$，在$(-\frac{1}{2a},+\infty)$上$f'(x)<0$。
（2）由（1）知$f(x)_{\max}=f(-\frac{1}{2a})=\ln(-\frac{1}{2a})+\frac{1}{4a}-1$。要证$f(x)\le -\frac{1}{4a}-2$，即证$\ln(-\frac{1}{2a})+\frac{1}{2a}+1\le 0$。令$t=-\frac{1}{2a}>0$，即证$\ln t-t+1\le 0$。设$g(t)=\ln t-t+1$，$g'(t)=\frac{1}{t}-1$，$g(t)$在$(0,1)$递增，在$(1,+\infty)$递减，$g(t)_{\max}=g(1)=0$，故原不等式成立。
\end{solution}
\end{question}

\begin{question}[points=5]
\textbf{来源：}2017 年 全国II卷-理科，第 11 题\par
若 $x=-2$ 是函数 $f(x)=(x^2+ax-1)e^{x-1}$ 的极值点，则 $f(x)$ 的极小值为（   ）
\begin{choices}
  \item $-1$
  \item $-2e^{-3}$
  \item $5e^{-3}$
  \item $1$
\end{choices}
\begin{solution}
\textbf{答案：}$A$

\textbf{分析：}求出函数的导数，利用极值点，求出 $a$，然后判断函数的单调性，求解函数的极小值即可。

\textbf{详解：}函数 $f(x)=(x^2+ax-1)e^{x-1}$，可得 $f'(x)=(2x+a)e^{x-1}+(x^2+ax-1)e^{x-1}$，$x=-2$ 是函数 $f(x)=(x^2+ax-1)e^{x-1}$ 的极值点，可得：$f'(-2)=(-4+a)e^{-3}+(4-2a-1)e^{-3}=0$，即 $-4+a+(3-2a)=0$．解得 $a=-1$．可得 $f'(x)=(2x-1)e^{x-1}+(x^2-x-1)e^{x-1}=(x^2+x-2)e^{x-1}$，函数的极值点为：$x=-2$，$x=1$，当 $x<-2$ 或 $x>1$ 时，$f'(x)>0$ 函数是增函数，$x\in(-2,1)$ 时，函数是减函数，$x=1$ 时，函数取得极小值：$f(1)=(1^2-1-1)e^{1-1}=-1$．故选：$A$。
\end{solution}
\end{question}

\begin{question}[points=12]
\textbf{来源：}2017 年 全国II卷-理科，第 21 题\par
已知函数 $f(x)=ax^2-ax-x\ln x$，且 $f(x)\ge0$．

（1）求 $a$；

（2）证明：$f(x)$ 存在唯一的极大值点 $x_0$，且 $e^{-2}<f(x_0)<2^{-2}$．
\begin{solution}
\textbf{答案：}（1）$a=1$；（2）略

\textbf{分析：}（1）通过分析可知 $f(x)\ge0$ 等价于 $h(x)=ax-a-\ln x\ge0$，进而利用 $h'(x)=a-\frac{1}{x}$ 可得 $h(x)_{\min}=h(\frac{1}{a})$，从而可得结论；（2）通过（1）可知 $f(x)=x^2-x-x\ln x$，记 $t(x)=f'(x)=2x-2-\ln x$，解不等式可知 $t(x)_{\min}=t(\frac{1}{2})=\ln2-1<0$，从而可知 $f'(x)=0$ 存在两根 $x_0$，$x_2$，利用 $f(x)$ 必存在唯一极大值点 $x_0$ 及 $x_0<\frac{1}{2}$ 可知 $f(x_0)<\frac{1}{4}$，另一方面可知 $f(x_0)>f(\frac{1}{e})=\frac{1}{e^2}$。

\textbf{详解：}（1）解：因为 $f(x)=ax^2-ax-x\ln x=x(ax-a-\ln x)$（$x>0$），则 $f(x)\ge0$ 等价于 $h(x)=ax-a-\ln x\ge0$，求导可知 $h'(x)=a-\frac{1}{x}$．则当 $a\le0$ 时 $h'(x)<0$，即 $y=h(x)$ 在 $(0,+\infty)$ 上单调递减，所以当 $x_0>1$ 时，$h(x_0)<h(1)=0$，矛盾，故 $a>0$．因为当 $0<x<\frac{1}{a}$ 时 $h'(x)<0$、当 $x>\frac{1}{a}$ 时 $h'(x)>0$，所以 $h(x)_{\min}=h(\frac{1}{a})$，又因为 $h(1)=a-a-\ln1=0$，所以 $\frac{1}{a}=1$，解得 $a=1$；
（2）证明：由（1）可知 $f(x)=x^2-x-x\ln x$，$f'(x)=2x-2-\ln x$，令 $f'(x)=0$，可得 $2x-2-\ln x=0$，记 $t(x)=2x-2-\ln x$，则 $t'(x)=2-\frac{1}{x}$，令 $t'(x)=0$，解得：$x=\frac{1}{2}$，所以 $t(x)$ 在区间 $(0,\frac{1}{2})$ 上单调递减，在 $(\frac{1}{2},+\infty)$ 上单调递增，所以 $t(x)_{\min}=t(\frac{1}{2})=\ln2-1<0$，从而 $t(x)=0$ 有解，即 $f'(x)=0$ 存在两根 $x_0$，$x_2$，且不妨设 $f'(x)$ 在 $(0,x_0)$ 上为正、在 $(x_0,x_2)$ 上为负、在 $(x_2,+\infty)$ 上为正，所以 $f(x)$ 必存在唯一极大值点 $x_0$，且 $2x_0-2-\ln x_0=0$，所以 $f(x_0)=x_0^2-x_0-x_0\ln x_0=x_0^2-x_0+2x_0-2x_0^2=x_0-x_0^2$，由 $x_0<\frac{1}{2}$ 可知 $f(x_0)<(x_0-x_0^2)_{\max}=-\frac{1}{4}+\frac{1}{2}=\frac{1}{4}$；由 $f'(\frac{1}{e})<0$ 可知 $x_0<\frac{1}{e}<\frac{1}{2}$，所以 $f(x)$ 在 $(0,x_0)$ 上单调递增，在 $(x_0,\frac{1}{e})$ 上单调递减，所以 $f(x_0)>f(\frac{1}{e})=\frac{1}{e^2}$；综上所述，$f(x)$ 存在唯一的极大值点 $x_0$，且 $e^{-2}<f(x_0)<2^{-2}$。
\end{solution}
\end{question}

\begin{question}[points=12]
\textbf{来源：}2017 年 全国II卷-文科，第 21 题\par
设函数 $f(x)=(1-x^2)e^x$. \\ (1) 讨论 $f(x)$ 的单调性; \\ (2) 当 $x\ge0$ 时, $f(x)\le ax+1$, 求 $a$ 的取值范围.
\begin{solution}

\textbf{分析：}(1) 求出函数的导数, 求出极值点, 利用导函数的符号, 判断函数的单调性即可. (2) 化简 $f(x)=(1-x)(1+x)e^x$. 下面对 $a$ 的范围进行讨论: ① 当 $a\ge1$ 时, ② 当 $0<a<1$ 时, ③ 当 $a\le0$ 时, 推出结论, 然后得到 $a$ 的取值范围.

\textbf{详解：}(1) 因为 $f(x)=(1-x^2)e^x$, $x\in\mathbb{R}$, 所以 $f'(x)=(1-2x-x^2)e^x$, 令 $f'(x)=0$ 可知 $x=-1\pm\sqrt{2}$, 当 $x<-1-\sqrt{2}$ 或 $x>-1+\sqrt{2}$ 时 $f'(x)<0$, 当 $-1-\sqrt{2}<x<-1+\sqrt{2}$ 时 $f'(x)>0$, 所以 $f(x)$ 在 $(-\infty,-1-\sqrt{2})$ 和 $(-1+\sqrt{2},+\infty)$ 上单调递减, 在 $(-1-\sqrt{2},-1+\sqrt{2})$ 上单调递增; \\ (2) 由题可知 $f(x)=(1-x)(1+x)e^x$. 下面对 $a$ 的范围进行讨论: ① 当 $a\ge1$ 时, 设函数 $h(x)=(1-x)e^x$, 则 $h'(x)=-xe^x<0$ $(x>0)$, 因此 $h(x)$ 在 $[0,+\infty)$ 上单调递减, 又因为 $h(0)=1$, 所以 $h(x)\le1$, 所以 $f(x)=(1+x)h(x)\le x+1\le ax+1$; ② 当 $0<a<1$ 时, 设函数 $g(x)=e^x-x-1$, 则 $g'(x)=e^x-1>0$ $(x>0)$, 所以 $g(x)$ 在 $[0,+\infty)$ 上单调递增, 又 $g(0)=1-0-1=0$, 所以 $e^x\ge x+1$. 因为当 $0<x<1$ 时 $f(x)>(1-x)(1+x)^2$, 所以 $(1-x)(1+x)^2-ax-1=x(1-a-x-x^2)$, 取 $x_0=\frac{\sqrt{5-4a}-1}{2}\in(0,1)$, 则 $(1-x_0)(1+x_0)^2-ax_0-1=0$, 所以 $f(x_0)>ax_0+1$, 矛盾; ③ 当 $a\le0$ 时, 取 $x_0=\frac{\sqrt{5}-1}{2}\in(0,1)$, 则 $f(x_0)>(1-x_0)(1+x_0)^2=1\ge ax_0+1$, 矛盾; 综上所述, $a$ 的取值范围是 $[1,+\infty)$.
\end{solution}
\end{question}

\begin{question}[points=12]
\textbf{来源：}2017 年 全国I卷-理科，第 21 题\par
已知函数 $f(x)=ae^{2x}+(a-2)e^x-x$．

（1）讨论 $f(x)$ 的单调性；

（2）若 $f(x)$ 有两个零点，求 $a$ 的取值范围．
\begin{solution}
\textbf{答案：}（1）见解析；（2）$(0,1)$

\textbf{分析：}（1）求导，根据导数与函数单调性的关系，分类讨论，即可求得 $f(x)$ 单调性；（2）由（1）可知：当 $a>0$ 时才有两个零点，根据函数的单调性求得 $f(x)$ 最小值，由 $f(x)_{\min}<0$，$g(a)=a\ln a+a-1$，$a>0$，求导，由 $g(a)_{\min}=g(e^{-2})=e^{-2}\ln e^{-2}+e^{-2}-1=-\frac{1}{e^2}-1$，$g(1)=0$，即可求得 $a$ 的取值范围．

\textbf{详解：}解：（1）由 $f(x)=ae^{2x}+(a-2)e^x-x$，求导 $f'(x)=2ae^{2x}+(a-2)e^x-1$，当 $a=0$ 时，$f'(x)=-2e^x-1<0$，所以当 $x\in\mathbb{R}$，$f(x)$ 单调递减，当 $a>0$ 时，$f'(x)=(2e^x+1)(ae^x-1)=2a(e^x+\frac{1}{2})(e^x-\frac{1}{a})$，令 $f'(x)=0$，解得：$x=\ln\frac{1}{a}$，当 $f'(x)>0$，解得：$x>\ln\frac{1}{a}$，当 $f'(x)<0$，解得：$x<\ln\frac{1}{a}$，所以$x\in(-\infty,\ln\frac{1}{a})$ 时，$f(x)$ 单调递减，$x\in(\ln\frac{1}{a},+\infty)$ 单调递增；当 $a<0$ 时，$f'(x)=2a(e^x+\frac{1}{2})(e^x-\frac{1}{a})<0$，恒成立，所以当 $x\in\mathbb{R}$，$f(x)$ 单调递减，综上可知：当 $a\leq0$ 时，$f(x)$ 在 $\mathbb{R}$ 单调减函数，当 $a>0$ 时，$f(x)$ 在 $(-\infty,\ln\frac{1}{a})$ 是减函数，在 $(\ln\frac{1}{a},+\infty)$ 是增函数．
（2）①若 $a\leq0$ 时，由（1）可知：$f(x)$ 最多有一个零点，当 $a>0$ 时，$f(x)=ae^{2x}+(a-2)e^x-x$，当 $x\to-\infty$ 时，$e^{2x}\to0$，$e^x\to0$，所以当 $x\to-\infty$ 时，$f(x)\to+\infty$，当 $x\to\infty$，$e^{2x}\to+\infty$，且远远大于 $e^x$ 和 $x$，所以当 $x\to\infty$，$f(x)\to+\infty$，所以函数有两个零点，$f(x)$ 的最小值小于 0 即可，由 $f(x)$ 在 $(-\infty,\ln\frac{1}{a})$ 是减函数，在 $(\ln\frac{1}{a},+\infty)$ 是增函数，所以$f(x)_{\min}=f(\ln\frac{1}{a})=a\times(\frac{1}{a})^2+(a-2)\times\frac{1}{a}-\ln\frac{1}{a}=1-\frac{2}{a}-\ln\frac{1}{a}<0$，所以$1-\frac{2}{a}+\ln a<0$，即 $\ln a+\frac{1}{a}-1>0$，设 $t=\frac{1}{a}$，则 $g(t)=\ln t+t-1$，($t>0$)，求导 $g'(t)=\frac{1}{t}+1$，由 $g(1)=0$，所以$t=\frac{1}{a}>1$，解得：$0<a<1$，所以$a$ 的取值范围 $(0,1)$．
\end{solution}
\end{question}

\begin{question}[points=5]
\textbf{来源：}2017 年 全国I卷-文科，第 14 题\par
曲线$y=x^2+\frac{1}{x}$在点$(1,2)$处的切线方程为．
\fillin[{$x-y+1=0$}]
\begin{solution}
\textbf{答案：}$x-y+1=0$

\textbf{分析：}求出函数的导数，求出切线的斜率，利用点斜式求解切线方程即可。

\textbf{详解：}解：曲线$y=x^2+\frac{1}{x}$，可得$y'=2x-\frac{1}{x^2}$，切线的斜率为：$k=2-1=1$。切线方程为：$y-2=x-1$，即：$x-y+1=0$。故答案为：$x-y+1=0$。
\end{solution}
\end{question}

\begin{question}[points=12]
\textbf{来源：}2017 年 全国I卷-文科，第 21 题\par
已知函数$f(x)=e^x(e^x-a)-a^2x$．

（1）讨论$f(x)$的单调性；

（2）若$f(x)\ge0$，求$a$的取值范围．
\begin{solution}
\textbf{答案：}（1）当$a=0$时，$f(x)$在$\mathbb{R}$上单调递增；当$a>0$时，$f(x)$在$(-\infty,\ln a)$上单调递减，在$(\ln a,+\infty)$上单调递增；当$a<0$时，$f(x)$在$(-\infty,\ln(-\frac{a}{2}))$上单调递减，在$(\ln(-\frac{a}{2}),+\infty)$上单调递增；（2）$[-2\sqrt{e},1]$。

\textbf{分析：}（1）先求导，再分类讨论，根据导数和函数的单调性即可判断；（2）根据（1）的结论，分别求出函数的最小值，即可求出$a$的范围。

\textbf{详解：}解：（1）$f(x)=e^x(e^x-a)-a^2x=e^{2x}-ae^x-a^2x$，$\therefore f'(x)=2e^{2x}-ae^x-a^2=(2e^x+a)(e^x-a)$，①当$a=0$时，$f'(x)>0$恒成立，$\therefore f(x)$在$\mathbb{R}$上单调递增，②当$a>0$时，$2e^x+a>0$，令$f'(x)=0$，解得$x=\ln a$，当$x<\ln a$时，$f'(x)<0$，函数$f(x)$单调递减，当$x>\ln a$时，$f'(x)>0$，函数$f(x)$单调递增，③当$a<0$时，$e^x-a>0$，令$f'(x)=0$，解得$x=\ln(-\frac{a}{2})$，当$x<\ln(-\frac{a}{2})$时，$f'(x)<0$，函数$f(x)$单调递减，当$x>\ln(-\frac{a}{2})$时，$f'(x)>0$，函数$f(x)$单调递增，综上所述，当$a=0$时，$f(x)$在$\mathbb{R}$上单调递增，当$a>0$时，$f(x)$在$(-\infty,\ln a)$上单调递减，在$(\ln a,+\infty)$上单调递增，当$a<0$时，$f(x)$在$(-\infty,\ln(-\frac{a}{2}))$上单调递减，在$(\ln(-\frac{a}{2}),+\infty)$上单调递增，
（2）①当$a=0$时，$f(x)=e^{2x}>0$恒成立，②当$a>0$时，由（1）可得$f(x)_{\min}=f(\ln a)=-a^2\ln a\ge0$，$\therefore\ln a\le0$，$\therefore0<a\le1$，③当$a<0$时，由（1）可得：$f(x)_{\min}=f(\ln(-\frac{a}{2}))=\frac{a^2}{4}-a^2\ln(-\frac{a}{2})\ge0$，$\therefore\ln(-\frac{a}{2})\le\frac{1}{4}$，$\therefore -2\sqrt{e}\le a<0$，综上所述$a$的取值范围为$[-2\sqrt{e},1]$．
\end{solution}
\end{question}

\section*{2018 年}

\begin{question}[points=5]
\textbf{来源：}2018 年 全国III卷-理科，第 14 题\par
曲线$y=(ax+1)e^x$在点$(0,1)$处的切线的斜率为$-2$，则$a=$．
\fillin[{$-3$}]
\begin{solution}
\textbf{答案：}$-3$

\textbf{分析：}求导，利用导数的几何意义列方程求解．

\textbf{详解：}$y'=ae^x+(ax+1)e^x=(ax+a+1)e^x$，则$y'|_{x=0}=a+1=-2$，解得$a=-3$．
\end{solution}
\end{question}

\begin{question}[points=12]
\textbf{来源：}2018 年 全国III卷-理科，第 21 题\par
已知函数$f(x)=(2+x+ax^2)\ln(1+x)-2x$．

（1）若$a=0$，证明：当$-1<x<0$时，$f(x)<0$；当$x>0$时，$f(x)>0$；

（2）若$x=0$是$f(x)$的极大值点，求$a$．
\begin{solution}

\textbf{分析：}（1）对函数$f(x)$两次求导数，分别判断$f'(x)$和$f(x)$的单调性，结合$f(0)=0$即可得出结论；（2）令$h(x)$为$f'(x)$的分子，令$h''(0)$计算$a$，讨论$a$的范围，得出$f(x)$的单调性，从而得出$a$的值．

\textbf{详解：}（1）当$a=0$时，$f(x)=(2+x)\ln(1+x)-2x$，$f'(x)=\ln(1+x)+\frac{2+x}{1+x}-2=\ln(1+x)-\frac{x}{1+x}$，$f''(x)=\frac{1}{1+x}-\frac{1}{(1+x)^2}=\frac{x}{(1+x)^2}$，当$-1<x<0$时，$f''(x)<0$，$f'(x)$递减；当$x>0$时，$f''(x)>0$，$f'(x)$递增，所以$f'(x)\geq f'(0)=0$，故$f(x)$在$(-1,+\infty)$上单调递增，又$f(0)=0$，所以当$-1<x<0$时$f(x)<0$，当$x>0$时$f(x)>0$．
（2）$f'(x)=(1+2ax)\ln(1+x)+\frac{2+x+ax^2}{1+x}-2=\frac{(1+2ax)(1+x)\ln(1+x)+ax^2-x}{1+x}$，令$h(x)=(1+2ax)(1+x)\ln(1+x)+ax^2-x$，则$h'(x)=4ax+(4ax+2a+1)\ln(1+x)$，$h''(x)=8a+4a\ln(1+x)+\frac{4ax+2a+1}{1+x}$．由题意$x=0$是极大值点，则$h'(0)=0$，$h''(0)\leq0$，$h''(0)=8a+2a+1=10a+1$，令$10a+1=0$得$a=-\frac{1}{10}$，下面验证：当$a=-\frac{1}{10}$时，$h''(x)=-\frac{4}{5}-\frac{2}{5}\ln(1+x)+\frac{-\frac{2}{5}x-\frac{1}{5}}{1+x}$，$h''(0)=0$，$h'''(x)=-\frac{2}{5(1+x)}+\frac{2}{5(1+x)^2}<0$，故$h''(x)$递减，当$x<0$时$h''(x)>0$，$h'(x)$递增，$h'(x)<h'(0)=0$，$h(x)$递减，$h(x)>h(0)=0$，$f'(x)>0$；当$x>0$时$h''(x)<0$，$h'(x)$递减，$h'(x)<h'(0)=0$，$h(x)$递减，$h(x)<h(0)=0$，$f'(x)<0$，所以$x=0$是极大值点，符合．若$a>-\frac{1}{10}$时，$h''(0)>0$，存在$x>0$使$h'(x)>0$，与极大值矛盾；若$a<-\frac{1}{10}$时，$h''(0)<0$，存在$x<0$使$h'(x)>0$，也与极大值矛盾．综上$a=-\frac{1}{10}$．
\end{solution}
\end{question}

\begin{question}[points=12]
\textbf{来源：}2018 年 全国III卷-文科，第 21 题\par
已知函数 $f(x)=\frac{ax^2+x-1}{e^x}$．

(1) 求曲线 $y=f(x)$ 在点 $(0,-1)$ 处的切线方程；

(2) 证明：当 $a\ge 1$ 时，$f(x)+e\ge 0$．
\begin{solution}

\textbf{分析：}(1) 由 $f'(x)=\frac{-ax^2+2ax-x+2}{e^x}$，$f'(0)=2$，可得切线斜率 $k=2$，即可得到切线方程．
(2) 可得 $f'(x)=\frac{-ax^2+(2a-1)x+2}{e^x}=\frac{-(ax+1)(x-2)}{e^x}$．令 $f'(x)=0$，可得 $x=-\frac{1}{a}$ 或 $x=2$．当 $x<-\frac{1}{a}$ 时，$f'(x)<0$；当 $-\frac{1}{a}<x<2$ 时，$f'(x)>0$；当 $x>2$ 时，$f'(x)<0$．$\therefore f(x)$ 在 $(-\infty,-\frac{1}{a})$，$(2,+\infty)$ 递减，在 $(-\frac{1}{a},2)$ 递增．注意到 $a\ge 1$ 时，函数 $g(x)=ax^2+x-1$ 在 $(2,+\infty)$ 单调递增，且 $g(2)=4a+1>0$，只需 $f(x)_{\min}\ge -e$，即可．

\textbf{详解：}(1) $f'(x)=\frac{(2ax+1)e^x-(ax^2+x-1)e^x}{e^{2x}}=\frac{-ax^2+2ax-x+2}{e^x}$，$\therefore f'(0)=2$，即曲线 $y=f(x)$ 在点 $(0,-1)$ 处的切线斜率 $k=2$，$\therefore$ 切线方程为 $y-(-1)=2x$，即 $2x-y-1=0$．
(2) 证明：函数 $f(x)$ 的定义域为 $\mathbb{R}$，$f'(x)=\frac{-ax^2+(2a-1)x+2}{e^x}=\frac{-(ax+1)(x-2)}{e^x}$．令 $f'(x)=0$，可得 $x=-\frac{1}{a}$ 或 $x=2$．当 $x<-\frac{1}{a}$ 时，$f'(x)<0$；当 $-\frac{1}{a}<x<2$ 时，$f'(x)>0$；当 $x>2$ 时，$f'(x)<0$．$\therefore f(x)$ 在 $(-\infty,-\frac{1}{a})$，$(2,+\infty)$ 递减，在 $(-\frac{1}{a},2)$ 递增．注意到 $a\ge 1$ 时，函数 $g(x)=ax^2+x-1$ 在 $(2,+\infty)$ 单调递增，且 $g(2)=4a+1>0$．$\because a\ge 1$，$\therefore -\frac{1}{a}\in[-1,0)$，则 $f(x)_{\min}=\min\{f(-\frac{1}{a}),f(2)\}$，$f(2)=\frac{4a+2-1}{e^2}=\frac{4a+1}{e^2}>0$，$f(-\frac{1}{a})=\frac{a\cdot\frac{1}{a^2}-\frac{1}{a}-1}{e^{-1/a}}=\frac{\frac{1}{a}-\frac{1}{a}-1}{e^{-1/a}}=-e^{1/a}\ge -e$，$\therefore f(x)\ge -e$，$\therefore$ 当 $a\ge 1$ 时，$f(x)+e\ge 0$．
\end{solution}
\end{question}

\begin{question}[points=5]
\textbf{来源：}2018 年 全国II卷-理科，第 13 题\par
曲线$y=2\ln(x+1)$在点$(0,0)$处的切线方程为                   ．
\fillin[{$y=2x$}]
\begin{solution}
\textbf{答案：}$y=2x$

\textbf{分析：}欲求出切线方程，只须求出其斜率即可，故先利用导数求出在$x=0$处的导函数值，再结合导数的几何意义即可求出切线的斜率．

\textbf{详解：}$y=2\ln(x+1)$，$y'=\frac{2}{x+1}$，当$x=0$时，$y'=2$，$\therefore$切线方程为$y=2x$．
\end{solution}
\end{question}

\begin{question}[points=12]
\textbf{来源：}2018 年 全国II卷-理科，第 21 题\par
已知函数$f(x)=e^x-ax^2$．

（1）若$a=1$，证明：当$x\ge 0$时，$f(x)\ge 1$；

（2）若$f(x)$在$(0,+\infty)$只有一个零点，求$a$．
\begin{solution}

\textbf{分析：}（1）通过两次求导，利用导数研究函数的单调性极值与最值即可证明，
（2）方法一、分离参数可得$a=\frac{e^x}{x^2}$在$(0,+\infty)$只有一个根，即函数$y=a$与$G(x)=\frac{e^x}{x^2}$的图象在$(0,+\infty)$只有一个交点．结合图象即可求得$a$．
方法二、：①当$a\le 0$时，$f(x)=e^x-ax^2>0$，$f(x)$在$(0,+\infty)$没有零点．②当$a>0$时，设函数$h(x)=1-ax^2e^{-x}$．$f(x)$在$(0,+\infty)$只有一个零点$\Leftrightarrow h(x)$在$(0,+\infty)$只有一个零点．利用$h'(x)=x(x-2)e^{-x}$，可得$h(x)$在$(0,2)$递减，在$(2,+\infty)$递增，结合函数$h(x)$图象即可求得$a$．

\textbf{详解：}（1）证明：当$a=1$时，函数$f(x)=e^x-x^2$．则$f'(x)=e^x-2x$，令$g(x)=e^x-2x$，则$g'(x)=e^x-2$，令$g'(x)=0$，得$x=\ln2$．当$x\in(0,\ln2)$时，$g'(x)<0$，当$x\in(\ln2,+\infty)$时，$g'(x)>0$，$\therefore g(x)\ge g(\ln2)=e^{\ln2}-2\cdot\ln2=2-2\ln2>0$，$\therefore f(x)$在$[0,+\infty)$单调递增，$\therefore f(x)\ge f(0)=1$；
（2）方法一、解：$f(x)$在$(0,+\infty)$只有一个零点$\Leftrightarrow$方程$e^x-ax^2=0$在$(0,+\infty)$只有一个根，$\Leftrightarrow a=\frac{e^x}{x^2}$在$(0,+\infty)$只有一个根，即函数$y=a$与$G(x)=\frac{e^x}{x^2}$的图象在$(0,+\infty)$只有一个交点．$G'(x)=\frac{e^x(x-2)}{x^3}$，当$x\in(0,2)$时，$G'(x)<0$，当$x\in(2,+\infty)$时，$G'(x)>0$，$\therefore G(x)$在$(0,2)$递减，在$(2,+\infty)$递增，当$x\to0^+$时，$G(x)\to+\infty$，当$x\to+\infty$时，$G(x)\to+\infty$，$\therefore f(x)$在$(0,+\infty)$只有一个零点时，$a=G(2)=\frac{e^2}{4}$．
方法二：①当$a\le 0$时，$f(x)=e^x-ax^2>0$，$f(x)$在$(0,+\infty)$没有零点．②当$a>0$时，设函数$h(x)=1-ax^2e^{-x}$．$f(x)$在$(0,+\infty)$只有一个零点$\Leftrightarrow h(x)$在$(0,+\infty)$只有一个零点．$h'(x)=x(x-2)e^{-x}$，当$x\in(0,2)$时，$h'(x)<0$，当$x\in(2,+\infty)$时，$h'(x)>0$，$\therefore h(x)$在$(0,2)$递减，在$(2,+\infty)$递增，$\therefore h(x)\ge h(2)=1-4ae^{-2}$（$x\ge0$）．当$h(2)<0$时，即$a>\frac{e^2}{4}$，由于$h(0)=1$，当$x>0$时，$e^x>x^2$，可得$h(4a)=1-\frac{16a^3}{e^{4a}}=1-\frac{16a^3}{(e^a)^4}$，因为$e^a>a^2$（$a>0$），所以$h(4a)>1-\frac{16a^3}{a^8}=1-\frac{16}{a^5}$，当$a$足够大时，$h(4a)>0$，$h(x)$在$(0,+\infty)$有$2$个零点；当$h(2)>0$时，即$a<\frac{e^2}{4}$，$h(x)$在$(0,+\infty)$没有零点；当$h(2)=0$时，即$a=\frac{e^2}{4}$，$h(x)$在$(0,+\infty)$只有一个零点．综上，$f(x)$在$(0,+\infty)$只有一个零点时，$a=\frac{e^2}{4}$．
\end{solution}
\end{question}

\begin{question}[points=5]
\textbf{来源：}2018 年 全国II卷-文科，第 13 题\par
曲线 $y=2\ln x$ 在点 $(1,0)$ 处的切线方程为
\fillin[{$y=2x-2$}]
\begin{solution}
\textbf{答案：}$y=2x-2$

\textbf{分析：}利用导数求出在 x=1 的导函数值，再结合导数的几何意义求出切线的斜率。

\textbf{详解：}因为$y=2\ln x$，所以$y'=\frac{2}{x}$，当 $x=1$ 时，$y'=2$，所以切线方程为 $y=2(x-1)$，即 $y=2x-2$。故答案为：$y=2x-2$。
\end{solution}
\end{question}

\begin{question}[points=12]
\textbf{来源：}2018 年 全国II卷-文科，第 21 题\par
已知函数 $f(x)=\frac{1}{3}x^3-a(x^2+x+1)$。

（1）若 $a=3$，求 $f(x)$ 的单调区间；

（2）证明：$f(x)$ 只有一个零点。
\begin{solution}

\textbf{分析：}（1）利用导数，求出极值点，判断导函数的符号，即可得到结果。
（2）分离参数后求导，先找点确定零点的存在性，再利用单调性确定唯一性。

\textbf{详解：}（1）当 $a=3$ 时，$f(x)=\frac{1}{3}x^3-3(x^2+x+1)$，所以 $f'(x)=x^2-6x-3$，令 $f'(x)=0$ 解得 $x=3\pm2\sqrt{3}$，当 $x\in(-\infty,3-2\sqrt{3})$ 和 $(3+2\sqrt{3},+\infty)$ 时，$f'(x)>0$，函数是增函数；当 $x\in(3-2\sqrt{3},3+2\sqrt{3})$ 时，$f'(x)<0$，函数是单调递减，综上，$f(x)$ 在 $(-\infty,3-2\sqrt{3})$ 和 $(3+2\sqrt{3},+\infty)$ 上是增函数，在 $(3-2\sqrt{3},3+2\sqrt{3})$ 上递减。
（2）证明：因为 $x^2+x+1=(x+\frac{1}{2})^2+\frac{3}{4}>0$，所以 $f(x)=0$ 等价于 $\frac{1}{3}x^3=a(x^2+x+1)$，令 $g(x)=\frac{x^3}{3(x^2+x+1)}$，则 $g'(x)=\frac{x^2(x^2+2x+3)}{3(x^2+x+1)^2}$，仅当 $x=0$ 时，$g'(x)=0$，所以 $g(x)$ 在 $R$ 上是增函数；$g(x)$ 至多有一个零点，从而 $f(x)$ 至多有一个零点。又因为 $f(3a-1)=-6a^2+2a-\frac{1}{3}=-6(a-\frac{1}{6})^2-\frac{1}{6}<0$，$f(3a+1)=\frac{1}{3}>0$，故 $f(x)$ 有一个零点，综上，$f(x)$ 只有一个零点。
\end{solution}
\end{question}

\begin{question}[points=5]
\textbf{来源：}2018 年 全国I卷-理科，第 5 题\par
设函数$f(x)=x^3+(a-1)x^2+ax$。若$f(x)$为奇函数，则曲线$y=f(x)$在点$(0,0)$处的切线方程为（    ）
\begin{choices}
  \item $y=-2x$
  \item $y=-x$
  \item $y=2x$
  \item $y=x$
\end{choices}
\begin{solution}
\textbf{答案：}D

\textbf{分析：}利用函数的奇偶性求出$a$，求出函数的导数，然后求解切线方程。

\textbf{详解：}由$f(x)$为奇函数，得$a=1$，$f(x)=x^3+x$，$f'(x)=3x^2+1$，$f'(0)=1$，切线方程为$y=x$。故选D。
\end{solution}
\end{question}

\begin{question}[points=12]
\textbf{来源：}2018 年 全国I卷-理科，第 21 题\par
已知函数$f(x)=\frac{1}{x}-x+a\ln x$。

（1）讨论$f(x)$的单调性；

（2）若$f(x)$存在两个极值点$x_1$，$x_2$，证明：$\frac{f(x_1)-f(x_2)}{x_1-x_2}<a-2$。
\begin{solution}

\textbf{分析：}（1）求出函数的定义域和导数，利用函数单调性和导数之间的关系进行求解即可。（2）将不等式进行等价转化，构造新函数，研究函数的单调性和最值即可得到结论。

\textbf{详解：}（1）定义域为$(0,+\infty)$，$f'(x)=-\frac{1}{x^2}-1+\frac{a}{x}=-\frac{x^2-ax+1}{x^2}$。设$g(x)=x^2-ax+1$。当$a\leq0$时，$g(x)>0$恒成立，$f'(x)<0$，$f(x)$在$(0,+\infty)$上单调递减。当$a>0$时，$\Delta=a^2-4$。若$0<a\leq2$，则$\Delta\leq0$，$g(x)\geq0$，$f'(x)\leq0$，$f(x)$在$(0,+\infty)$上单调递减。若$a>2$，则$g(x)=0$有两正根$x_1=\frac{a-\sqrt{a^2-4}}{2}$，$x_2=\frac{a+\sqrt{a^2-4}}{2}$，且$0<x_1<x_2$。当$x\in(0,x_1)\cup(x_2,+\infty)$时，$f'(x)<0$；当$x\in(x_1,x_2)$时，$f'(x)>0$。所以$f(x)$在$(0,x_1)$和$(x_2,+\infty)$上单调递减，在$(x_1,x_2)$上单调递增。
（2）由（1）知$a>2$，且$x_1x_2=1$，$x_1+x_2=a$。$\frac{f(x_1)-f(x_2)}{x_1-x_2}=-\frac{1}{x_1x_2}-1+a\frac{\ln x_1-\ln x_2}{x_1-x_2}=-2+a\frac{\ln x_1-\ln x_2}{x_1-x_2}$。要证$\frac{f(x_1)-f(x_2)}{x_1-x_2}<a-2$，即证$a\frac{\ln x_1-\ln x_2}{x_1-x_2}<a$，即$\frac{\ln x_1-\ln x_2}{x_1-x_2}<1$。不妨设$x_1<x_2$，则$x_1<1<x_2$，只需证$\ln x_1-\ln x_2>x_1-x_2$，即$\ln\frac{x_1}{x_2}>x_1-x_2$。令$t=\frac{x_1}{x_2}\in(0,1)$，则$x_1=tx_2$，由$x_1x_2=1$得$x_2=\frac{1}{\sqrt{t}}$，$x_1=\sqrt{t}$，则不等式化为$\ln t^2>\sqrt{t}-\frac{1}{\sqrt{t}}$，即$2\ln t>\sqrt{t}-\frac{1}{\sqrt{t}}$。构造函数$h(t)=2\ln t-\sqrt{t}+\frac{1}{\sqrt{t}}$，$t\in(0,1)$，$h'(t)=\frac{2}{t}-\frac{1}{2\sqrt{t}}-\frac{1}{2t\sqrt{t}}=\frac{4\sqrt{t}-t-1}{2t\sqrt{t}}=\frac{-(\sqrt{t}-1)^2}{2t\sqrt{t}}<0$，所以$h(t)$在$(0,1)$上单调递减，$h(t)>h(1)=0$，即$2\ln t-\sqrt{t}+\frac{1}{\sqrt{t}}>0$，故$2\ln t>\sqrt{t}-\frac{1}{\sqrt{t}}$成立，原不等式得证。
\end{solution}
\end{question}

\begin{question}[points=5]
\textbf{来源：}2018 年 全国I卷-文科，第 6 题\par
设函数 $f(x)=x^3+(a-1)x^2+ax$。若 $f(x)$ 为奇函数，则曲线 $y=f(x)$ 在点 $(0,0)$ 处的切线方程为
\begin{choices}
  \item $y=-2x$
  \item $y=-x$
  \item $y=2x$
  \item $y=x$
\end{choices}
\begin{solution}
\textbf{答案：}$y=x$

\textbf{分析：}利用函数的奇偶性求出 $a$，求出函数的导数，求出切线的斜率然后求解切线方程。

\textbf{详解：}函数 $f(x)=x^3+(a-1)x^2+ax$，若 $f(x)$ 为奇函数，可得 $a=1$，所以函数 $f(x)=x^3+x$，可得 $f'(x)=3x^2+1$，曲线 $y=f(x)$ 在点 $(0,0)$ 处的切线的斜率为 $1$，则曲线 $y=f(x)$ 在点 $(0,0)$ 处的切线方程为 $y=x$。故选：$D$。
\end{solution}
\end{question}

\begin{question}[points=12]
\textbf{来源：}2018 年 全国I卷-文科，第 21 题\par
已知函数 $f(x)=ae^x-\ln x-1$。

（1）设 $x=2$ 是 $f(x)$ 的极值点，求 $a$，并求 $f(x)$ 的单调区间；

（2）证明：当 $a\ge\frac{1}{e}$ 时，$f(x)\ge0$。
\begin{solution}

\textbf{分析：}（1）推导出 $x>0$，$f'(x)=ae^x-\frac{1}{x}$，由 $x=2$ 是极值点，解得 $a=\frac{1}{2e^2}$，从而 $f(x)=\frac{1}{2e^2}e^x-\ln x-1$，$f'(x)=\frac{1}{2}e^{x-2}-\frac{1}{x}$，由此能求出单调区间。（2）当 $a\ge\frac{1}{e}$ 时，$f(x)\ge\frac{1}{e}e^x-\ln x-1$，设 $g(x)=\frac{1}{e}e^x-\ln x-1$，则 $g'(x)=\frac{1}{e}e^x-\frac{1}{x}$，利用导数性质能证明。

\textbf{详解：}（1）$\because f(x)=ae^x-\ln x-1$，$x>0$，$f'(x)=ae^x-\frac{1}{x}$，$\because x=2$ 是 $f(x)$ 的极值点，$\therefore f'(2)=ae^2-\frac{1}{2}=0$，解得 $a=\frac{1}{2e^2}$，$\therefore f(x)=\frac{1}{2e^2}e^x-\ln x-1=\frac{1}{2}e^{x-2}-\ln x-1$，$f'(x)=\frac{1}{2}e^{x-2}-\frac{1}{x}$，当 $0<x<2$ 时，$f'(x)<0$，当 $x>2$ 时，$f'(x)>0$，$\therefore f(x)$ 在 $(0,2)$ 单调递减，在 $(2,+\infty)$ 单调递增。
（2）证明：当 $a\ge\frac{1}{e}$ 时，$f(x)\ge\frac{1}{e}e^x-\ln x-1=e^{x-1}-\ln x-1$，设 $g(x)=e^{x-1}-\ln x-1$，则 $g'(x)=e^{x-1}-\frac{1}{x}$，当 $0<x<1$ 时，$g'(x)<0$，当 $x>1$ 时，$g'(x)>0$，$\therefore x=1$ 是 $g(x)$ 的最小值点，故当 $x>0$ 时，$g(x)\ge g(1)=0$，$\therefore$ 当 $a\ge\frac{1}{e}$ 时，$f(x)\ge0$。
\end{solution}
\end{question}

\section*{2019 年}

\begin{question}[points=5]
\textbf{来源：}2019 年 全国III卷-理科，第 6 题\par
已知曲线 $y = ae^x + x\ln x$ 在点 $(1,ae)$ 处的切线方程为 $y=2x+b$，则
\begin{choices}
  \item $a=e, b=-1$
  \item $a=e, b=1$
  \item $a=e^{-1}, b=1$
  \item $a=e^{-1}, b=-1$
\end{choices}
\begin{solution}
\textbf{答案：}D

\textbf{分析：}通过求导数，确定切线斜率，求得 $a$，将点的坐标代入直线方程，求得 $b$。

\textbf{详解：}$y' = ae^x + \ln x + 1$，
$k = y'|_{x=1} = ae + 1 = 2$，所以 $a = e^{-1}$。
将点 $(1,1)$ 代入 $y=2x+b$ 得 $2+b=1$，$b=-1$。故选 D。
\end{solution}
\end{question}

\begin{question}[points=12]
\textbf{来源：}2019 年 全国III卷-理科，第 20 题\par
已知函数 $f(x)=2x^3-ax^2+b$。

（1）讨论 $f(x)$ 的单调性；

（2）是否存在 $a,b$，使得 $f(x)$ 在区间 $[0,1]$ 的最小值为 $-1$ 且最大值为 $1$？若存在，求出 $a,b$ 的所有值；若不存在，说明理由。
\begin{solution}
\textbf{答案：}（1）见详解；（2）$\begin{cases} a=0 \\ b=-1 \end{cases}$ 或 $\begin{cases} a=4 \\ b=1 \end{cases}$。

\textbf{分析：}(1)先求 $f(x)$ 的导数，再根据 $a$ 的范围分情况讨论函数单调性；(2)根据 $a$ 的各种范围，利用函数单调性进行最大值和最小值的判断，最终得出 $a,b$ 的值。

\textbf{详解：}（1）$f'(x)=6x^2-2ax=2x(3x-a)$。令 $f'(x)=0$，得 $x=0$ 或 $x=\dfrac{a}{3}$。
若 $a>0$，则当 $x\in(-\infty,0)\cup\left(\dfrac{a}{3},+\infty\right)$ 时，$f'(x)>0$；当 $x\in\left(0,\dfrac{a}{3}\right)$ 时，$f'(x)<0$。故 $f(x)$ 在 $(-\infty,0)$，$\left(\dfrac{a}{3},+\infty\right)$ 单调递增，在 $\left(0,\dfrac{a}{3}\right)$ 单调递减。
若 $a=0$，$f(x)$ 在 $(-\infty,+\infty)$ 单调递增。
若 $a<0$，则当 $x\in\left(-\infty,\dfrac{a}{3}\right)\cup(0,+\infty)$ 时，$f'(x)>0$；当 $x\in\left(\dfrac{a}{3},0\right)$ 时，$f'(x)<0$。故 $f(x)$ 在 $\left(-\infty,\dfrac{a}{3}\right)$，$(0,+\infty)$ 单调递增，在 $\left(\dfrac{a}{3},0\right)$ 单调递减。
（2）满足题设条件的 $a,b$ 存在。
（i）当 $a\le 0$ 时，由（1）知，$f(x)$ 在 $[0,1]$ 单调递增，所以 $f(x)$ 在区间 $[0,1]$ 的最小值为 $f(0)=b$，最大值为 $f(1)=2-a+b$。此时 $a,b$ 满足题设条件当且仅当 $b=-1$，$2-a+b=1$，即 $a=0$，$b=-1$。
（ii）当 $a\ge 3$ 时，由（1）知，$f(x)$ 在 $[0,1]$ 单调递减，所以 $f(x)$ 在区间 $[0,1]$ 的最大值为 $f(0)=b$，最小值为 $f(1)=2-a+b$。此时 $a,b$ 满足题设条件当且仅当 $2-a+b=-1$，$b=1$，即 $a=4$，$b=1$。
（iii）当 $0<a<3$ 时，由（1）知，$f(x)$ 在 $[0,1]$ 的最小值为 $f\left(\dfrac{a}{3}\right)=-\dfrac{a^3}{27}+b$，最大值为 $b$ 或 $2-a+b$。
若 $-\dfrac{a^3}{27}+b=-1$，$b=1$，则 $a=3\sqrt[3]{2}$，与 $0<a<3$ 矛盾。
若 $-\dfrac{a^3}{27}+b=-1$，$2-a+b=1$，则 $a=3\sqrt{3}$ 或 $a=-3\sqrt{3}$ 或 $a=0$，与 $0<a<3$ 矛盾。
综上，当且仅当 $a=0$，$b=-1$ 或 $a=4$，$b=1$ 时，$f(x)$ 在 $[0,1]$ 的最小值为 $-1$，最大值为 $1$。
\end{solution}
\end{question}

\begin{question}[points=5]
\textbf{来源：}2019 年 全国III卷-文科，第 7 题\par
已知曲线$y=ae^x+x\ln x$在点$(1,ae)$处的切线方程为$y=2x+b$，则
\begin{choices}
  \item $a=e,b=-1$
  \item $a=e,b=1$
  \item $a=e^{-1},b=1$
  \item $a=e^{-1},b=-1$
\end{choices}
\begin{solution}
\textbf{答案：}D

\textbf{分析：}通过求导数，确定得到切线斜率的表达式，求得$a$，将点的坐标代入直线方程，求得$b$。

\textbf{详解：}$y'=ae^x+\ln x+1$，$k=y'|_{x=1}=ae+1=2$，$\therefore a=e^{-1}$。将$(1,1)$代入$y=2x+b$得$2+b=1,b=-1$，故选D。
\end{solution}
\end{question}

\begin{question}[points=12]
\textbf{来源：}2019 年 全国III卷-文科，第 20 题\par
已知函数$f(x)=2x^3-ax^2+2$。

(1)讨论$f(x)$的单调性；

(2)当$0<a<3$时，记$f(x)$在区间$[0,1]$的最大值为M，最小值为m，求$M-m$的取值范围。
\begin{solution}
\textbf{答案：}(1)见详解; (2) $[\frac{8}{27},2)$

\textbf{分析：}(1)先求$f(x)$的导数，再根据a的范围分情况讨论函数单调性；(2)讨论a的范围，利用函数单调性进行最大值和最小值的判断，最终求得$M-m$的取值范围。

\textbf{详解：}(1)对$f(x)=2x^3-ax^2+2$求导得$f'(x)=6x^2-2ax=6x(x-\frac{a}{3})$。所以有：
当$a<0$时，$(-\infty,\frac{a}{3})$区间上单调递增，$(\frac{a}{3},0)$区间上单调递减，$(0,+\infty)$区间上单调递增；
当$a=0$时，$(-\infty,+\infty)$区间上单调递增；
当$a>0$时，$(-\infty,0)$区间上单调递增，$(0,\frac{a}{3})$区间上单调递减，$(\frac{a}{3},+\infty)$区间上单调递增。
(2)若$0<a\le2$，$f(x)$在区间$(0,\frac{a}{3})$单调递减，在区间$(\frac{a}{3},1)$单调递增，所以区间$[0,1]$上最小值为$f(\frac{a}{3})$。而$f(0)=2,f(1)=4-a\ge f(0)$，故区间$[0,1]$上最大值为$f(1)$。所以$M-m=f(1)-f(\frac{a}{3})=(4-a)-[2(\frac{a}{3})^3-a(\frac{a}{3})^2+2]=-a+\frac{a^3}{27}+2$。设$g(x)=\frac{x^3}{27}-x+2$，求导$g'(x)=\frac{x^2}{9}-1$，当$0<x\le2$时$g'(x)<0$从而$g(x)$单调递减。而$0<a\le2$，所以$\frac{8}{27}\le -a+\frac{a^3}{27}+2<2$，即$M-m$的取值范围是$[\frac{8}{27},2)$。
若$2<a<3$，$f(x)$在区间$(0,\frac{a}{3})$单调递减，在区间$(\frac{a}{3},1)$单调递增，所以区间$[0,1]$上最小值为$f(\frac{a}{3})$，而$f(0)=2,f(1)=4-a\le f(0)$，故区间$[0,1]$上最大值为$f(0)$。所以$M-m=f(0)-f(\frac{a}{3})=2-[2(\frac{a}{3})^3-a(\frac{a}{3})^2+2]=\frac{a^3}{27}$，而$2<a<3$，所以$\frac{8}{27}<\frac{a^3}{27}<1$，即$M-m$的取值范围是$(\frac{8}{27},1)$。
综上得$M-m$的取值范围是$[\frac{8}{27},2)$。
\end{solution}
\end{question}

\begin{question}[points=12]
\textbf{来源：}2019 年 全国II卷-理科，第 20 题\par
已知函数 $f(x)=\ln x-\frac{x+1}{x-1}$．

（1）讨论 $f(x)$ 的单调性，并证明 $f(x)$ 有且仅有两个零点；

（2）设 $x_0$ 是 $f(x)$ 的一个零点，证明曲线 $y=\ln x$ 在点 $A(x_0,\ln x_0)$ 处的切线也是曲线 $y=e^x$ 的切线．
\begin{solution}
\textbf{答案：}（1）函数 $f(x)$ 在 $(0,1)$ 和 $(1,+\infty)$ 上是单调增函数，证明见解析；（2）证明见解析．

\textbf{分析：}（1）对函数 $f(x)$ 求导，结合定义域，判断函数的单调性；（2）先求出曲线 $y=\ln x$ 在 $A(x_0,\ln x_0)$ 处的切线 $l$，然后求出当曲线 $y=e^x$ 切线的斜率与 $l$ 斜率相等时，证明曲线 $y=e^x$ 切线 $l'$ 在纵轴上的截距与 $l$ 在纵轴的截距相等即可．

\textbf{详解：}（1）$f(x)$ 的定义域为 $(0,1)\cup(1,+\infty)$，$f'(x)=\frac{x^2+1}{x(x-1)^2}>0$，所以 $f(x)$ 在 $(0,1)$ 和 $(1,+\infty)$ 上是单调增函数．因为 $f(e)=1-\frac{e+1}{e-1}=\frac{-2}{e-1}<0$，$f(e^2)=2-\frac{e^2+1}{e^2-1}=\frac{e^2-3}{e^2-1}>0$，所以 $f(x)$ 在 $(1,+\infty)$ 有唯一零点 $x_1$，即 $f(x_1)=0$．又 $0<\frac{1}{x_1}<1$，$f(\frac{1}{x_1})=-\ln x_1+\frac{x_1+1}{x_1-1}=-f(x_1)=0$，故 $f(x)$ 在 $(0,1)$ 有唯一零点 $\frac{1}{x_1}$．综上，$f(x)$ 有且仅有两个零点．
（2）因为 $\frac{1}{x_0}=e^{-\ln x_0}$，故点 $B(-\ln x_0,\frac{1}{x_0})$ 在曲线 $y=e^x$ 上．由题设知 $f(x_0)=0$，即 $\ln x_0=\frac{x_0+1}{x_0-1}$，故直线 $AB$ 的斜率 $k=\frac{\frac{1}{x_0}-\ln x_0}{-\ln x_0-x_0}=\frac{\frac{1}{x_0}-\frac{x_0+1}{x_0-1}}{-\frac{x_0+1}{x_0-1}-x_0}=\frac{1}{x_0}$．曲线 $y=e^x$ 在点 $B(-\ln x_0,\frac{1}{x_0})$ 处切线的斜率是 $\frac{1}{x_0}$，曲线 $y=\ln x$ 在点 $A(x_0,\ln x_0)$ 处切线的斜率也是 $\frac{1}{x_0}$，所以曲线 $y=\ln x$ 在点 $A(x_0,\ln x_0)$ 处的切线也是曲线 $y=e^x$ 的切线．
\end{solution}
\end{question}

\begin{question}[points=5]
\textbf{来源：}2019 年 全国II卷-文科，第 10 题\par
曲线 $y=2\sin x+\cos x$ 在点 $(\pi,-1)$ 处的切线方程为
\begin{choices}
  \item $x-y-\pi-1=0$
  \item $2x-y-2\pi-1=0$
  \item $2x+y-2\pi+1=0$
  \item $x+y-\pi+1=0$
\end{choices}
\begin{solution}
\textbf{答案：}$C$

\textbf{分析：}先判定点 $(\pi,-1)$ 是否为切点，再利用导数的几何意义求解．

\textbf{详解：}当 $x=\pi$ 时，$y=2\sin\pi+\cos\pi=-1$，即点 $(\pi,-1)$ 在曲线 $y=2\sin x+\cos x$ 上．$y'=2\cos x-\sin x$，$y'|_{x=\pi}=2\cos\pi-\sin\pi=-2$，则 $y=2\sin x+\cos x$ 在点 $(\pi,-1)$ 处的切线方程为 $y-(-1)=-2(x-\pi)$，即 $2x+y-2\pi+1=0$．故选 $C$．
\end{solution}
\end{question}

\begin{question}[points=12]
\textbf{来源：}2019 年 全国II卷-文科，第 21 题\par
已知函数 $f(x)=(x-1)\ln x-x-1$．证明：

（1）$f(x)$ 存在唯一的极值点；

（2）$f(x)=0$ 有且仅有两个实根，且两个实根互为倒数．
\begin{solution}
\textbf{答案：}（1）见详解；（2）见详解

\textbf{分析：}（1）先对函数 $f(x)$ 求导，根据导函数的单调性，得到存在唯一 $x_0$，使得 $f'(x_0)=0$，进而判断函数 $f(x)$ 的单调性，即可确定其极值点个数，证明出结论成立；
（2）先由（1）的结果，得到 $f(x_0)<f(1)=-2<0$，$f(e^2)=e^2-3>0$，得到 $f(x)=0$ 在 $(x_0,+\infty)$ 内存在唯一实根，记作 $x=\alpha$，再求出 $f\left(\dfrac{1}{\alpha}\right)=0$，即可结合题意，说明结论成立．

\textbf{详解：}（1）由题意可得，$f(x)$ 的定义域为 $(0,+\infty)$，由 $f(x)=(x-1)\ln x-x-1$，得 $f'(x)=\ln x+\dfrac{x-1}{x}-1=\ln x-\dfrac{1}{x}$，显然 $f'(x)=\ln x-\dfrac{1}{x}$ 单调递增；又 $f'(1)=-1<0$，$f'(2)=\ln2-\dfrac{1}{2}=\dfrac{\ln4-1}{2}>0$，故存在唯一 $x_0$，使得 $f'(x_0)=0$；又当 $x>x_0$ 时，$f'(x)>0$，函数 $f(x)$ 单调递增；当 $0<x<x_0$ 时，$f'(x)<0$，函数 $f(x)$ 单调递减；因此，$f(x)$ 存在唯一的极值点．
（2）由（1）知，$f(x_0)<f(1)=-2$，又 $f(e^2)=e^2-3>0$，所以 $f(x)=0$ 在 $(x_0,+\infty)$ 内存在唯一实根，记作 $x=\alpha$．由 $1<x_0<\alpha$ 得 $\dfrac{1}{\alpha}<1<x_0$，又 $f\left(\dfrac{1}{\alpha}\right)=\left(\dfrac{1}{\alpha}-1\right)\ln\dfrac{1}{\alpha}-\dfrac{1}{\alpha}-1=\dfrac{f(\alpha)}{\alpha}=0$，故 $\dfrac{1}{\alpha}$ 是方程 $f(x)=0$ 在 $(0,x_0)$ 内的唯一实根；综上，$f(x)=0$ 有且仅有两个实根，且两个实根互为倒数．
\end{solution}
\end{question}

\begin{question}[points=5]
\textbf{来源：}2019 年 全国I卷-理科，第 13 题\par
曲线 $y = 3(x^2 + x)e^x$ 在点 $(0, 0)$ 处的切线方程为．
\fillin[{$3x - y = 0$}]
\begin{solution}
\textbf{答案：}$3x - y = 0$

\textbf{分析：}本题根据导数的几何意义，通过求导数，确定得到切线的斜率，利用直线方程的点斜式求得切线方程．

\textbf{详解：}$y' = 3(2x + 1)e^x + 3(x^2 + x)e^x = 3(x^2 + 3x + 1)e^x$，所以，$k = y'|_{x=0} = 3$，所以，曲线 $y = 3(x^2 + x)e^x$ 在点 $(0, 0)$ 处的切线方程为 $y = 3x$，即 $3x - y = 0$．
\end{solution}
\end{question}

\begin{question}[points=12]
\textbf{来源：}2019 年 全国I卷-理科，第 20 题\par
已知函数 $f(x) = \sin x - \ln(1 + x)$, $f'(x)$ 为 $f(x)$ 的导数．证明：

(1) $f'(x)$ 在区间 $\left(-1, \frac{\pi}{2}\right)$ 存在唯一极大值点；

(2) $f(x)$ 有且仅有 $2$ 个零点．
\begin{solution}
\textbf{答案：}见解析

\textbf{分析：}(1) 求得导函数后，可判断出导函数在 $\left(-1, \frac{\pi}{2}\right)$ 上单调递减，根据零点存在定理可判断出 $\exists x_0 \in \left(0, \frac{\pi}{2}\right)$，使得 $g'(x_0) = 0$，进而得到导函数在 $\left(-1, \frac{\pi}{2}\right)$ 上的单调性，从而可证得结论；(2) 由（1）的结论可知 $x = 0$ 为 $f(x)$ 在 $(-1, 0]$ 上的唯一零点；当 $x \in \left(0, \frac{\pi}{2}\right]$ 时，首先可判断出在 $(0, x_0)$ 上无零点，再利用零点存在定理得到 $f(x)$ 在 $\left(x_0, \frac{\pi}{2}\right)$ 上的单调性，可知 $f(x) > 0$，不存在零点；当 $x \in \left[\frac{\pi}{2}, \pi\right]$ 时，利用零点存在定理和 $f(x)$ 单调性可判断出存在唯一一个零点；当 $x \in (\pi, +\infty)$，可证得 $f(x) < 0$；综合上述情况可证得结论．

\textbf{详解：}(1) 由题意知：$f(x)$ 定义域为 $(-1, +\infty)$ 且 $f'(x) = \cos x - \frac{1}{x+1}$．令 $g(x) = \cos x - \frac{1}{x+1}$, $x \in \left(-1, \frac{\pi}{2}\right)$，$\therefore g'(x) = -\sin x + \frac{1}{(x+1)^2}$, $x \in \left(-1, \frac{\pi}{2}\right)$．$\because \frac{1}{(x+1)^2}$ 在 $\left(-1, \frac{\pi}{2}\right)$ 上单调递减，$-\sin x$ 在 $\left(-1, \frac{\pi}{2}\right)$ 上单调递减，$\therefore g'(x)$ 在 $\left(-1, \frac{\pi}{2}\right)$ 上单调递减．又 $g'(0) = -\sin 0 + 1 = 1 > 0$, $g'\left(\frac{\pi}{2}\right) = -\sin\frac{\pi}{2} + \frac{4}{(\pi+2)^2} = -1 + \frac{4}{(\pi+2)^2} < 0$，$\therefore \exists x_0 \in \left(0, \frac{\pi}{2}\right)$，使得 $g'(x_0) = 0$．$\therefore$ 当 $x \in (-1, x_0)$ 时，$g'(x) > 0$；当 $x \in \left(x_0, \frac{\pi}{2}\right)$ 时，$g'(x) < 0$．即 $g(x)$ 在 $(-1, x_0)$ 上单调递增；在 $\left(x_0, \frac{\pi}{2}\right)$ 上单调递减，则 $x = x_0$ 为 $g(x)$ 唯一的极大值点，即 $f'(x)$ 在区间 $\left(-1, \frac{\pi}{2}\right)$ 上存在唯一的极大值点 $x_0$．
(2) 由（1）知：$f'(x) = \cos x - \frac{1}{x+1}$, $x \in (-1, +\infty)$．
①当 $x \in (-1, 0]$ 时，由（1）可知 $f'(x)$ 在 $(-1, 0]$ 上单调递增，$\therefore f'(x) \le f'(0) = 0$，$\therefore f(x)$ 在 $(-1, 0]$ 上单调递减，又 $f(0) = 0$，$\therefore x = 0$ 为 $f(x)$ 在 $(-1, 0]$ 上的唯一零点．
②当 $x \in \left(0, \frac{\pi}{2}\right]$ 时，$f'(x)$ 在 $(0, x_0)$ 上单调递增，在 $\left(x_0, \frac{\pi}{2}\right)$ 上单调递减，又 $f'(0) = 0$，$\therefore f'(x_0) > 0$，$\therefore f(x)$ 在 $(0, x_0)$ 上单调递增，此时 $f(x) > f(0) = 0$，不存在零点．又 $f'\left(\frac{\pi}{2}\right) = \cos\frac{\pi}{2} - \frac{2}{\pi+2} = -\frac{2}{\pi+2} < 0$，$\therefore \exists x_1 \in \left(x_0, \frac{\pi}{2}\right)$，使得 $f'(x_1) = 0$，$\therefore f(x)$ 在 $(x_0, x_1)$ 上单调递增，在 $\left(x_1, \frac{\pi}{2}\right)$ 上单调递减．又 $f(x_0) > f(0) = 0$, $f\left(\frac{\pi}{2}\right) = \sin\frac{\pi}{2} - \ln\left(1 + \frac{\pi}{2}\right) = \ln\frac{2e}{\pi+2} > \ln 1 = 0$，$\therefore f(x) > 0$ 在 $\left(x_0, \frac{\pi}{2}\right)$ 上恒成立，此时不存在零点．
③当 $x \in \left[\frac{\pi}{2}, \pi\right]$ 时，$\sin x$ 单调递减，$-\ln(x+1)$ 单调递减，$\therefore f(x)$ 在 $\left[\frac{\pi}{2}, \pi\right]$ 上单调递减．又 $f\left(\frac{\pi}{2}\right) > 0$, $f(\pi) = \sin\pi - \ln(\pi+1) = -\ln(\pi+1) < 0$，即 $f(\pi) \cdot f\left(\frac{\pi}{2}\right) < 0$，又 $f(x)$ 在 $\left[\frac{\pi}{2}, \pi\right]$ 上单调递减，$\therefore f(x)$ 在 $\left[\frac{\pi}{2}, \pi\right]$ 上存在唯一零点．
④当 $x \in (\pi, +\infty)$ 时，$\sin x \in [-1, 1]$, $\ln(x+1) > \ln(\pi+1) > \ln e = 1$，$\therefore \sin x - \ln(x+1) < 0$，即 $f(x)$ 在 $(\pi, +\infty)$ 上不存在零点．
综上所述：$f(x)$ 有且仅有 $2$ 个零点．
\end{solution}
\end{question}

\begin{question}[points=5]
\textbf{来源：}2019 年 全国I卷-文科，第 13 题\par
曲线 $y = 3(x^2+x)e^x$ 在点 $(0,0)$ 处的切线方程为．
\fillin[{$3x - y = 0$}]
\begin{solution}
\textbf{答案：}$3x - y = 0$

\textbf{分析：}本题根据导数的几何意义，通过求导数，确定得到切线的斜率，利用直线方程的点斜式求得切线方程。

\textbf{详解：}$y' = 3(2x+1)e^x + 3(x^2+x)e^x = 3(x^2+3x+1)e^x$，所以 $k = y'|_{x=0} = 3$，所以曲线 $y = 3(x^2+x)e^x$ 在点 $(0,0)$ 处的切线方程为 $y = 3x$，即 $3x - y = 0$。
\end{solution}
\end{question}

\begin{question}[points=12]
\textbf{来源：}2019 年 全国I卷-文科，第 20 题\par
已知函数 $f(x) = 2\sin x - x\cos x - x$，$f'(x)$ 为 $f(x)$ 的导数．

（1）证明：$f'(x)$ 在区间 $(0,\pi)$ 存在唯一零点；

（2）若 $x \in [0,\pi]$ 时，$f(x) \ge ax$，求 $a$ 的取值范围．
\begin{solution}
\textbf{答案：}（1）见解析；（2）$a \in (-\infty, 0]$。

\textbf{分析：}（1）求导得到导函数后，设为 $g(x)$ 进行再次求导，可判断出当 $x \in (0,\frac{\pi}{2})$ 时 $g'(x)>0$，当 $x \in (\frac{\pi}{2},\pi)$ 时 $g'(x)<0$，从而得到 $g(x)$ 单调性，由零点存在定理可判断出唯一零点；（2）构造函数 $h(x)=f(x)-ax$，通过二次求导可判断出 $h'(x)$ 的最值，分别在 $a \le -2$，$-2<a\le0$，$0<a<\frac{\pi-2}{2}$ 和 $a \ge \frac{\pi-2}{2}$ 的情况下根据导函数的符号判断 $h(x)$ 单调性，从而确定 $h(x)\ge0$ 恒成立时 $a$ 的取值范围。

\textbf{详解：}（1）$f'(x)=2\cos x - \cos x + x\sin x -1 = \cos x + x\sin x -1$，令 $g(x)=\cos x + x\sin x -1$，则 $g'(x)=-\sin x + \sin x + x\cos x = x\cos x$．当 $x \in (0,\pi)$ 时，令 $g'(x)=0$，解得 $x=\frac{\pi}{2}$，当 $x\in(0,\frac{\pi}{2})$ 时 $g'(x)>0$；当 $x\in(\frac{\pi}{2},\pi)$ 时 $g'(x)<0$，所以 $g(x)$ 在 $(0,\frac{\pi}{2})$ 上单调递增，在 $(\frac{\pi}{2},\pi)$ 上单调递减．又 $g(0)=1-1=0$，$g(\frac{\pi}{2})=\frac{\pi}{2}-1>0$，$g(\pi)=-1-1=-2$，所以 $g(x)$ 在 $(0,\frac{\pi}{2})$ 无零点，在 $(\frac{\pi}{2},\pi)$ 存在唯一零点，即 $f'(x)$ 在 $(0,\pi)$ 存在唯一零点。
（2）若 $x\in[0,\pi]$ 时，$f(x)\ge ax$，即 $f(x)-ax\ge0$ 恒成立，令 $h(x)=f(x)-ax=2\sin x - x\cos x - (a+1)x$，则 $h'(x)=\cos x + x\sin x -1 - a$，$h''(x)=x\cos x = g'(x)$．由（1）可知，$h'(x)$ 在 $(0,\frac{\pi}{2})$ 上单调递增，在 $(\frac{\pi}{2},\pi)$ 上单调递减，且 $h'(0)=-a$，$h'(\frac{\pi}{2})=\frac{\pi-2}{2}-a$，$h'(\pi)=-2-a$，所以 $h'(x)_{\min}=h'(\pi)=-2-a$，$h'(x)_{\max}=h'(\frac{\pi}{2})=\frac{\pi-2}{2}-a$．
①当 $a\le -2$ 时，$h'(x)_{\min}=-2-a\ge0$，即 $h'(x)\ge0$ 在 $[0,\pi]$ 上恒成立，$h(x)$ 单调递增，$h(x)\ge h(0)=0$，恒成立；
②当 $-2<a\le0$ 时，$h'(0)\ge0$，$h'(\frac{\pi}{2})>0$，$h'(\pi)<0$，$\exists x_1\in(\frac{\pi}{2},\pi)$ 使得 $h'(x_1)=0$，$h(x)$ 在 $[0,x_1)$ 单调递增，在 $(x_1,\pi]$ 单调递减，又 $h(0)=0$，$h(\pi)=2\sin \pi - \pi\cos \pi - (a+1)\pi = -a\pi \ge 0$，所以恒成立；
③当 $0<a<\frac{\pi-2}{2}$ 时，$h'(0)<0$，$h'(\frac{\pi}{2})>0$，$\exists x_2\in(0,\frac{\pi}{2})$ 使得 $h'(x_2)=0$，$h(x)$ 在 $[0,x_2)$ 单调递减，在 $(x_2,\frac{\pi}{2})$ 单调递增，则 $x\in(0,x_2)$ 时 $h(x)<h(0)=0$，不恒成立；
④当 $a\ge \frac{\pi-2}{2}$ 时，$h'(x)_{\max}=h'(\frac{\pi}{2})=\frac{\pi-2}{2}-a\le0$，$h(x)$ 在 $(0,\frac{\pi}{2})$ 单调递减，$h(x)<h(0)=0$，不恒成立．
综上所述：$a\in(-\infty,0]$。
\end{solution}
\end{question}

\section*{2020 年}

\begin{question}[points=12]
\textbf{来源：}2020 年 全国III卷-理科，第 21 题\par
设函数 $f(x)=x^3+bx+c$，曲线 $y=f(x)$ 在点 $\left(\dfrac12,f\left(\dfrac12\right)\right)$ 处的切线与 $y$ 轴垂直。

(1) 求 $b$；

(2) 若 $f(x)$ 有一个绝对值不大于 $1$ 的零点，证明：$f(x)$ 所有零点的绝对值都不大于 $1$。
\begin{solution}
\textbf{答案：}见解析

\textbf{分析：}(1) 利用导数几何意义求 $b$；(2) 利用导数研究函数单调性，结合零点存在性定理反证。

\textbf{详解：}(1) $f'(x)=3x^2+b$，由题意 $f'\left(\frac12\right)=0$，即 $3\times\frac14+b=0$，得 $b=-\frac34$。
(2) 由 (1) $f(x)=x^3-\frac34x+c$，$f'(x)=3x^2-\frac34=3\left(x+\frac12\right)\left(x-\frac12\right)$。$f(x)$ 在 $\left(-\frac12,\frac12\right)$ 递减，在 $\left(-\infty,-\frac12\right)$ 和 $\left(\frac12,+\infty\right)$ 递增。$f(-1)=c-\frac14$，$f\left(-\frac12\right)=c+\frac14$，$f\left(\frac12\right)=c-\frac14$，$f(1)=c+\frac14$。假设存在一个零点 $x_0$ 绝对值大于 $1$，则 $x_0>1$ 或 $x_0<-1$。若 $x_0>1$，则 $f(1)=c+\frac14<0$，$c<-\frac14$，则 $f\left(-\frac12\right)=c+\frac14<0$，$f\left(\frac12\right)=c-\frac14<0$，$f(-1)=c-\frac14<0$。又 $f(-4c)=64c^3+3c+c=4c(1-16c^2)>0$（因为 $c<-\frac14$ 时 $c<0$，$1-16c^2<0$，乘积为正），由零点存在性定理知 $f(x)$ 在 $(1,-4c)$ 上有唯一零点，在 $(-\infty,1)$ 上无零点，此时 $f(x)$ 不存在绝对值不大于 $1$ 的零点，矛盾。同理 $x_0<-1$ 也矛盾。所以所有零点的绝对值都不大于 $1$。
\end{solution}
\end{question}

\begin{question}[points=5]
\textbf{来源：}2020 年 全国III卷-文科，第 15 题\par
设函数 $f(x)=\frac{e^x}{x+a}$，若 $f'(1)=\frac{e}{4}$，则 $a=$
\fillin[{$1$}]
\begin{solution}
\textbf{答案：}$1$

\textbf{分析：}由题意首先求得导函数的解析式，然后得到关于实数 $a$ 的方程，解方程即可确定实数 $a$ 的值.

\textbf{详解：}由函数的解析式可得 $f'(x)=\frac{e^x(x+a)-e^x}{(x+a)^2}=\frac{e^x(x+a-1)}{(x+a)^2}$，则 $f'(1)=\frac{e^1\times(1+a-1)}{(1+a)^2}=\frac{ae}{(a+1)^2}$，据此可得 $\frac{ae}{(a+1)^2}=\frac{e}{4}$，整理得 $a^2-2a+1=0$，解得 $a=1$.
\end{solution}
\end{question}

\begin{question}[points=12]
\textbf{来源：}2020 年 全国III卷-文科，第 20 题\par
已知函数 $f(x)=x^3-kx+k^2$．

(1) 讨论 $f(x)$ 的单调性；

(2) 若 $f(x)$ 有三个零点，求 $k$ 的取值范围．
\begin{solution}
\textbf{答案：}(1) 当 $k\leq0$ 时，$f(x)$ 在 $(-\infty,+\infty)$ 上单调递增；当 $k>0$ 时，$f(x)$ 在 $(-\sqrt{\frac{k}{3}},\sqrt{\frac{k}{3}})$ 上单调递减，在 $(-\infty,-\sqrt{\frac{k}{3}})$ 和 $(\sqrt{\frac{k}{3}},+\infty)$ 上单调递增. (2) $(0,\frac{4}{27})$

\textbf{分析：}(1) $f'(x)=3x^2-k$，对 $k$ 分 $k\leq0$ 和 $k>0$ 两种情况讨论即可；(2) $f(x)$ 有三个零点，由(1)知 $k>0$，且 $\begin{cases} f(-\sqrt{\frac{k}{3}})>0 \\ f(\sqrt{\frac{k}{3}})<0 \end{cases}$，解不等式组得到 $k$ 的范围，再利用零点存在性定理加以说明即可.

\textbf{详解：}(1) 由题，$f'(x)=3x^2-k$. 当 $k\leq0$ 时，$f'(x)\geq0$ 恒成立，所以 $f(x)$ 在 $(-\infty,+\infty)$ 上单调递增；当 $k>0$ 时，令 $f'(x)=0$，得 $x=\pm\sqrt{\frac{k}{3}}$，令 $f'(x)<0$，得 $-\sqrt{\frac{k}{3}}<x<\sqrt{\frac{k}{3}}$，令 $f'(x)>0$，得 $x<-\sqrt{\frac{k}{3}}$ 或 $x>\sqrt{\frac{k}{3}}$，所以 $f(x)$ 在 $(-\sqrt{\frac{k}{3}},\sqrt{\frac{k}{3}})$ 上单调递减，在 $(-\infty,-\sqrt{\frac{k}{3}})$ 和 $(\sqrt{\frac{k}{3}},+\infty)$ 上单调递增.
(2) 由(1)知，$f(x)$ 有三个零点，则 $k>0$，且 $\begin{cases} f(-\sqrt{\frac{k}{3}})>0 \\ f(\sqrt{\frac{k}{3}})<0 \end{cases}$，即 $\begin{cases} k^2+\frac{2}{3}k\sqrt{\frac{k}{3}}>0 \\ k^2-\frac{2}{3}k\sqrt{\frac{k}{3}}<0 \end{cases}$，解得 $0<k<\frac{4}{27}$. 当 $0<k<\frac{4}{27}$ 时，$k>\sqrt{\frac{k}{3}}$，且 $f(k)=k^2>0$，所以 $f(x)$ 在 $(\sqrt{\frac{k}{3}},k)$ 上有唯一一个零点；同理 $-k-1<-\sqrt{\frac{k}{3}}$，$f(-k-1)=-k^3-(k+1)^2<0$，所以 $f(x)$ 在 $(-k-1,-\sqrt{\frac{k}{3}})$ 上有唯一一个零点；又 $f(x)$ 在 $(-\sqrt{\frac{k}{3}},\sqrt{\frac{k}{3}})$ 上有唯一一个零点，所以 $f(x)$ 有三个零点. 综上可知 $k$ 的取值范围为 $(0,\frac{4}{27})$.
\end{solution}
\end{question}

\begin{question}[points=12]
\textbf{来源：}2020 年 全国II卷-理科，第 21 题；次专题：三角函数\par
已知函数$f(x)=\sin^2x\sin2x$.

（1）讨论$f(x)$在区间$(0,\pi)$的单调性；

（2）证明：$|f(x)|\leq\dfrac{3\sqrt3}{8}$；

（3）设$n\in\mathbb{N}^*$，证明：$\sin^2x\sin^22x\sin^24x\cdots\sin^22^nx\leq\dfrac{3^n}{4^n}$.
\begin{solution}
\textbf{答案：}（1）当$x\in(0,\frac\pi3)$时，$f'(x)>0,f(x)$单调递增，当$x\in(\frac\pi3,\frac{2\pi}3)$时，$f'(x)<0,f(x)$单调递减，当$x\in(\frac{2\pi}3,\pi)$时，$f'(x)>0,f(x)$单调递增.（2）证明见解析；（3）证明见解析.

\textbf{分析：}(1)首先求得导函数的解析式，然后由导函数的零点确定其在各个区间上的符号，最后确定原函数的单调性即可；(2)首先确定函数的周期性，然后结合(1)中的结论确定函数在一个周期内的最大值和最小值即可证得题中的不等式；(3)对所给的不等式左侧进行恒等变形，然后结合(2)的结论和三角函数的有界性进行放缩即可证得题中的不等式.

\textbf{详解：}（1）$f(x)=\sin^2x\sin2x=2\sin^3x\cos x$，则$f'(x)=2(3\sin^2x\cos^2x-\sin^4x)=2\sin^2x(3\cos^2x-\sin^2x)=2\sin^2x(4\cos^2x-1)=2\sin^2x(2\cos x+1)(2\cos x-1)$，$f'(x)=0$在$x\in(0,\pi)$上的根为$x_1=\frac\pi3,x_2=\frac{2\pi}3$，当$x\in(0,\frac\pi3)$时，$f'(x)>0,f(x)$单调递增，当$x\in(\frac\pi3,\frac{2\pi}3)$时，$f'(x)<0,f(x)$单调递减，当$x\in(\frac{2\pi}3,\pi)$时，$f'(x)>0,f(x)$单调递增.
（2）注意到$f(x+\pi)=\sin^2(x+\pi)\sin[2(x+\pi)]=\sin^2x\sin2x=f(x)$，故函数$f(x)$是周期为$\pi$的函数，结合(1)的结论，计算可得：$f(0)=f(\pi)=0$，$f(\frac\pi3)=(\frac{\sqrt3}{2})^2\times\frac{\sqrt3}{2}=\frac{3\sqrt3}{8}$，$f(\frac{2\pi}3)=(\frac{\sqrt3}{2})^2\times(-\frac{\sqrt3}{2})=-\frac{3\sqrt3}{8}$，据此可得：$[f(x)]_{\max}=\frac{3\sqrt3}{8}$，$[f(x)]_{\min}=-\frac{3\sqrt3}{8}$，即$|f(x)|\leq\frac{3\sqrt3}{8}$.
（3）$\sin^2x\sin^22x\sin^24x\cdots\sin^22^nx=\left[\sin^3x\sin^32x\sin^34x\cdots\sin^32^nx\right]^{\frac23}=\left[\sin x(\sin^2x\sin2x)(\sin^22x\sin4x)\cdots(\sin^22^{n-1}x\sin2^nx)\sin^22^nx\right]^{\frac23}\leq\left[|\sin x|\times\frac{3\sqrt3}{8}\times\frac{3\sqrt3}{8}\times\cdots\times\frac{3\sqrt3}{8}\times|\sin2^nx|\right]^{\frac23}=\left[\left(\frac{3\sqrt3}{8}\right)^n\right]^{\frac23}=\left(\frac{3}{4}\right)^n$.
\end{solution}
\end{question}

\begin{question}[points=12]
\textbf{来源：}2020 年 全国II卷-文科，第 21 题\par
已知函数 $f(x)=2\ln x+1$。

（1）若 $f(x)\le 2x+c$，求 $c$ 的取值范围；

（2）设 $a>0$ 时，讨论函数 $g(x)=\frac{f(x)-f(a)}{x-a}$ 的单调性。
\begin{solution}
\textbf{答案：}（1）$c\ge -1$；（2）$g(x)$ 在区间 $(0,a)$ 和 $(a,+\infty)$ 上单调递减，没有递增区间。

\textbf{分析：}（1）不等式 $f(x)\le 2x+c$ 转化为 $f(x)-2x-c\le 0$，构造新函数，利用导数求出新函数的最大值，进而进行求解；
（2）对函数 $g(x)$ 求导，把导函数 $g'(x)$ 的分子构成一个新函数 $m(x)$，再求导得到 $m'(x)$，根据 $m'(x)$ 的正负，判断 $m(x)$ 的单调性，进而确定 $g'(x)$ 的正负性，最后求出函数 $g(x)$ 的单调性。

\textbf{详解：}（1）$f(x)\le 2x+c \Leftrightarrow f(x)-2x-c\le 0 \Leftrightarrow 2\ln x+1-2x-c\le 0$。
设 $h(x)=2\ln x+1-2x-c(x>0)$，则 $h'(x)=\frac{2}{x}-2=\frac{2(1-x)}{x}$。
当 $x>1$ 时，$h'(x)<0$，$h(x)$ 单调递减；当 $0<x<1$ 时，$h'(x)>0$，$h(x)$ 单调递增。
所以当 $x=1$ 时，$h(x)$ 有最大值 $h(1)=2\ln1+1-2\times1-c=-1-c$。
要想不等式恒成立，只需 $h(1)\le 0$，即 $-1-c\le 0$，解得 $c\ge -1$。
（2）$g(x)=\frac{2\ln x+1-(2\ln a+1)}{x-a}=\frac{2(\ln x-\ln a)}{x-a} \quad (x>0$ 且 $x\neq a)$。
$g'(x)=2\cdot\frac{\frac{1}{x}(x-a)-(\ln x-\ln a)}{(x-a)^2}=2\cdot\frac{(x-a)-x(\ln x-\ln a)}{x(x-a)^2}$。
设 $m(x)= (x-a)-x(\ln x-\ln a)$，则 $m'(x)=1-(\ln x-\ln a)-1 = \ln a - \ln x$。
当 $x>a$ 时，$\ln x>\ln a$，$m'(x)<0$，$m(x)$ 单调递减，所以 $m(x)<m(a)=0$，从而 $g'(x)<0$，$g(x)$ 单调递减；
当 $0<x<a$ 时，$\ln x<\ln a$，$m'(x)>0$，$m(x)$ 单调递增，所以 $m(x)<m(a)=0$，从而 $g'(x)<0$，$g(x)$ 单调递减。
所以函数 $g(x)$ 在区间 $(0,a)$ 和 $(a,+\infty)$ 上单调递减，没有递增区间。
\end{solution}
\end{question}

\begin{question}[points=5]
\textbf{来源：}2020 年 全国I卷-理科，第 6 题\par
函数$f(x)=x^4-2x^3$的图像在点$(1,f(1))$处的切线方程为
\begin{choices}
  \item $y=-2x-1$
  \item $y=-2x+1$
  \item $y=2x-3$
  \item $y=2x+1$
\end{choices}
\begin{solution}
\textbf{答案：}B

\textbf{详解：}因为$f(x)=x^4-2x^3$，所以$f'(x)=4x^3-6x^2$，$f(1)=-1$，$f'(1)=-2$，因此所求切线方程为$y+1=-2(x-1)$，即$y=-2x+1$。
\end{solution}
\end{question}

\begin{question}[points=12]
\textbf{来源：}2020 年 全国I卷-理科，第 21 题\par
已知函数$f(x)=e^x+ax^2-x$。

(1)当$a=1$时，讨论$f(x)$的单调性；

(2)当$x\geq 0$时，$f(x)\geq\frac{1}{2}x^3+1$，求$a$的取值范围。
\begin{solution}
\textbf{答案：}(1)当$x\in(-\infty,0)$时，$f(x)$单调递减；当$x\in(0,+\infty)$时，$f(x)$单调递增。(2)$a\in\left[\frac{7-e^2}{4},+\infty\right)$

\textbf{详解：}(1)当$a=1$时，$f(x)=e^x+x^2-x$，$f'(x)=e^x+2x-1$，$f''(x)=e^x+2>0$，故$f'(x)$单调递增。又$f'(0)=0$，所以当$x<0$时，$f'(x)<0$，$f(x)$单调递减；当$x>0$时，$f'(x)>0$，$f(x)$单调递增。
(2)由$f(x)\geq\frac{1}{2}x^3+1$得$e^x+ax^2-x\geq\frac{1}{2}x^3+1$，即$e^x-\frac{1}{2}x^3-x-1\geq -ax^2$。当$x=0$时，$1\geq1$成立。当$x>0$时，分离参数得$a\geq-\frac{e^x-\frac{1}{2}x^3-x-1}{x^2}$。令$g(x)=-\frac{e^x-\frac{1}{2}x^3-x-1}{x^2}$，则$g'(x)=-\frac{(x-2)\left(e^x-\frac{1}{2}x^2-x-1\right)}{x^3}$。令$h(x)=e^x-\frac{1}{2}x^2-x-1$，则$h'(x)=e^x-x-1$，$h''(x)=e^x-1\geq0$，故$h'(x)$单调递增，$h'(x)\geq h'(0)=0$，$h(x)$单调递增，$h(x)\geq h(0)=0$。所以当$x\in(0,2)$时，$g'(x)>0$，$g(x)$单调递增；当$x\in(2,+\infty)$时，$g'(x)<0$，$g(x)$单调递减。故$g(x)_{\max}=g(2)=\frac{7-e^2}{4}$，因此$a\geq\frac{7-e^2}{4}$。综上，$a$的取值范围是$\left[\frac{7-e^2}{4},+\infty\right)$。
\end{solution}
\end{question}

\begin{question}[points=5]
\textbf{来源：}2020 年 全国I卷-文科，第 15 题\par
曲线 $y = \ln x + x + 1$ 的一条切线的斜率为 $2$, 则该切线的方程为
\fillin[{$y = 2x$}]
\begin{solution}
\textbf{答案：}$y = 2x$

\textbf{分析：}设切线的切点坐标为 $(x_0, y_0)$, 对函数求导，利用 $y'|_{x_0} = 2$, 求出 $x_0$, 代入曲线方程求出 $y_0$, 得到切线的点斜式方程，化简即可.

\textbf{详解：}设切点坐标为 $(x_0, y_0)$, $y' = \frac{1}{x} + 1$, 则 $y'|_{x=x_0} = \frac{1}{x_0} + 1 = 2$, 解得 $x_0 = 1$, 代入得 $y_0 = \ln 1 + 1 + 1 = 2$, 所以切线方程为 $y - 2 = 2(x - 1)$, 即 $y = 2x$.
\end{solution}
\end{question}

\begin{question}[points=12]
\textbf{来源：}2020 年 全国I卷-文科，第 20 题\par
已知函数 $f(x) = e^x - a(x+2)$.



(1) 当 $a=1$ 时，讨论 $f(x)$ 的单调性;



(2) 若 $f(x)$ 有两个零点，求 $a$ 的取值范围.
\begin{solution}
\textbf{答案：}(1) 减区间为 $(-\infty, 0)$, 增区间为 $(0, +\infty)$; (2) $\left(\frac{1}{e}, +\infty\right)$.

\textbf{分析：}(1) 将 $a=1$ 代入函数解析式，对函数求导，分别令导数大于零和小于零，求得函数的单调增区间和减区间； (2) 若 $f(x)$ 有两个零点，即 $e^x - a(x+2)=0$ 有两个解，将其转化为 $a = \frac{e^x}{x+2}$ 有两个解，令 $h(x)=\frac{e^x}{x+2} (x\neq -2)$, 求导研究函数图象的走向，从而求得结果.

\textbf{详解：}(1) 当 $a=1$ 时，$f(x)=e^x-(x+2)$, $f'(x)=e^x-1$, 令 $f'(x)<0$ 解得 $x<0$, 令 $f'(x)>0$ 解得 $x>0$, 所以 $f(x)$ 的减区间为 $(-\infty,0)$, 增区间为 $(0,+\infty)$.

(2) 若 $f(x)$ 有两个零点，即 $e^x - a(x+2)=0$ 有两个解，从方程可知 $x=-2$ 不成立，即 $a = \frac{e^x}{x+2}$ 有两个解. 令 $h(x)=\frac{e^x}{x+2} (x\neq -2)$, 则 $h'(x)=\frac{e^x(x+2)-e^x}{(x+2)^2}=\frac{e^x(x+1)}{(x+2)^2}$. 令 $h'(x)>0$ 解得 $x>-1$, 令 $h'(x)<0$ 解得 $x<-2$ 或 $-2<x<-1$, 所以 $h(x)$ 在 $(-\infty,-2)$ 和 $(-2,-1)$ 上单调递减，在 $(-1,+\infty)$ 上单调递增，且当 $x<-2$ 时 $h(x)<0$, 当 $x\to -2^+$ 时 $h(x)\to +\infty$, 当 $x\to +\infty$ 时 $h(x)\to +\infty$. 所以当 $a = \frac{e^x}{x+2}$ 有两个解时，有 $a > h(-1) = \frac{1}{e}$, 所以 $a$ 的取值范围是 $\left(\frac{1}{e}, +\infty\right)$.
\end{solution}
\end{question}

\section*{2021 年}

\begin{question}[points=15]
\textbf{来源：}2021 年 北京卷，第 19 题\par
已知函数 $f(x)=\frac{3-2x}{x^2+a}$.

(1) 若 $a=0$, 求 $y=f(x)$ 在 $(1,f(1))$ 处的切线方程;

(2) 若函数 $f(x)$ 在 $x=-1$ 处取得极值, 求 $f(x)$ 的单调区间, 以及最大值和最小值.
\begin{solution}
\textbf{答案：}(1) $4x+y-5=0$; (2) 增区间为 $(-\infty,-1)$ 和 $(4,+\infty)$, 减区间为 $(-1,4)$, 最大值为 $1$, 最小值为 $-\frac{1}{4}$.

\textbf{分析：}(1) 求出 $f(1),f'(1)$ 的值, 利用点斜式写出切线方程; (2) 由 $f'(-1)=0$ 求 $a$, 然后利用导数分析单调性与极值.

\textbf{详解：}(1) 当 $a=0$ 时, $f(x)=\frac{3-2x}{x^2}$, $f'(x)=\frac{2(x-3)}{x^3}$, $f(1)=1$, $f'(1)=-4$, 切线方程为 $y-1=-4(x-1)$, 即 $4x+y-5=0$.
(2) $f(x)=\frac{3-2x}{x^2+a}$, $f'(x)=\frac{-2(x^2+a)-2x(3-2x)}{(x^2+a)^2}=\frac{2(x^2-3x-a)}{(x^2+a)^2}$, 由 $f'(-1)=0$ 得 $\frac{2(1+3-a)}{(1+a)^2}=0$, 解得 $a=4$. 故 $f(x)=\frac{3-2x}{x^2+4}$, $f'(x)=\frac{2(x+1)(x-4)}{(x^2+4)^2}$. 列表:
| $x$ | $(-\infty,-1)$ | $-1$ | $(-1,4)$ | $4$ | $(4,+\infty)$ |
|-----|----------------|------|----------|-----|---------------|
| $f'(x)$ | $+$ | $0$ | $-$ | $0$ | $+$ |
| $f(x)$ | 增 | 极大值 | 减 | 极小值 | 增 |
所以增区间为 $(-\infty,-1)$ 和 $(4,+\infty)$, 减区间为 $(-1,4)$. 当 $x<\frac{3}{2}$ 时 $f(x)>0$, $x>\frac{3}{2}$ 时 $f(x)<0$, 所以 $f(x)_{\max}=f(-1)=1$, $f(x)_{\min}=f(4)=-\frac{1}{4}$.
\end{solution}
\end{question}

\begin{question}[points=5]
\textbf{来源：}2021 年 全国甲卷-理科，第 13 题\par
曲线 $y = \frac{2x-1}{x+2}$ 在点 $(-1, -3)$ 处的切线方程为．
\fillin[{$5x - y + 2 = 0$}]
\begin{solution}
\textbf{答案：}$5x - y + 2 = 0$

\textbf{分析：}先验证点在曲线上，再求导，代入切线方程公式即可．

\textbf{详解：}$y' = \frac{2(x+2)- (2x-1)}{(x+2)^2} = \frac{5}{(x+2)^2}$，则 $y'|_{x=-1} = 5$，所以切线方程为 $y+3=5(x+1)$，即 $5x - y + 2 = 0$．
\end{solution}
\end{question}

\begin{question}[points=12]
\textbf{来源：}2021 年 全国甲卷-理科，第 21 题\par
已知 $a>0$ 且 $a\neq 1$，函数 $f(x) = \frac{x^a}{a^x}$（$x>0$）．（1）当 $a=2$ 时，求 $f(x)$ 的单调区间；（2）若曲线 $y=f(x)$ 与直线 $y=1$ 有且仅有两个交点，求 $a$ 的取值范围．
\begin{solution}
\textbf{答案：}（1）单调递增区间 $\left(0, \frac{2}{\ln 2}\right]$，单调递减区间 $\left[\frac{2}{\ln 2}, +\infty\right)$；（2）$(1,e) \cup (e, +\infty)$

\textbf{分析：}（1）求得函数的导函数，利用导函数的正负与函数的单调性的关系即可得到函数的单调性；（2）利用指数对数的运算法则，可以将曲线 $y=f(x)$ 与直线 $y=1$ 有且仅有两个交点等价转化为方程 $\frac{\ln x}{x} = \frac{\ln a}{a}$ 有两个不同的实数根，即曲线 $y=g(x)$ 与直线 $y=\frac{\ln a}{a}$ 有两个交点，利用导函数研究 $g(x)$ 的单调性，并结合 $g(x)$ 的正负，零点和极限值分析图象，进而得到 $0<\frac{\ln a}{a}<\frac{1}{e}$，即 $a>1$ 且 $a\neq e$．

\textbf{详解：}（1）$a=2$ 时，$f(x)=\frac{x^2}{2^x}$，$f'(x)=\frac{x\cdot2^x(2-x\ln2)}{4^x}=\frac{x(2-x\ln2)}{2^x}$，由 $f'(x)=0$ 得 $x=\frac{2}{\ln2}$，所以 $f(x)$ 在 $\left(0,\frac{2}{\ln2}\right]$ 单调递增，在 $\left[\frac{2}{\ln2}, +\infty\right)$ 单调递减．（2）$f(x)=1$ 等价于 $a^x = x^a$，即 $x\ln a = a\ln x$，即 $\frac{\ln x}{x} = \frac{\ln a}{a}$．设 $g(x)=\frac{\ln x}{x}$，则 $g'(x)=\frac{1-\ln x}{x^2}$，$g(x)$ 在 $(0,e]$ 单调递增，在 $[e,+\infty)$ 单调递减，最大值为 $g(e)=\frac{1}{e}$．曲线 $y=f(x)$ 与 $y=1$ 有两个交点等价于 $y=g(x)$ 与 $y=\frac{\ln a}{a}$ 有两个交点，故 $0<\frac{\ln a}{a}<\frac{1}{e}$．函数 $g(a)$ 在 $(0,e)$ 单调递增，$(e,+\infty)$ 单调递减，且 $g(1)=0$，$g(e)=1/e$，结合图象得 $a\in(1,e)\cup(e,+\infty)$．
\end{solution}
\end{question}

\begin{question}[points=12]
\textbf{来源：}2021 年 全国甲卷-文科，第 20 题\par
设函数 $f(x) = a^2 x^2 + ax - 3\ln x + 1$, 其中 $a>0$.



（1）讨论 $f(x)$ 的单调性；



（2）若 $y=f(x)$ 的图象与 $x$ 轴没有公共点，求 $a$ 的取值范围.
\begin{solution}
\textbf{答案：}（1）$f(x)$ 的减区间为 $\left(0, \frac{1}{a}\right)$, 增区间为 $\left(\frac{1}{a}, +\infty\right)$；（2）$a > \frac{1}{e}$.

\textbf{分析：}（1）求出函数的导数，讨论其符号后可得函数的单调性.（2）根据 $f(1)>0$ 及（1）的单调性性可得 $f(x)_{\min} > 0$, 从而可求 $a$ 的取值范围.

\textbf{详解：}（1）函数的定义域为 $(0, +\infty)$, 又 $f'(x) = \frac{(2ax+3)(ax-1)}{x}$. 因为 $a>0$, $x>0$, 故 $2ax+3>0$. 当 $0 < x < \frac{1}{a}$ 时, $f'(x) < 0$; 当 $x > \frac{1}{a}$ 时, $f'(x) > 0$. 所以 $f(x)$ 的减区间为 $\left(0, \frac{1}{a}\right)$, 增区间为 $\left(\frac{1}{a}, +\infty\right)$.

（2）因为 $f(1)=a^2+a+1>0$ 且 $y=f(x)$ 的图与 $x$ 轴没有公共点, 所以 $y=f(x)$ 的图象在 $x$ 轴的上方. 由（1）中函数的单调性可得 $f(x)_{\min} = f\left(\frac{1}{a}\right) = 3 - 3\ln\frac{1}{a} = 3 + 3\ln a$, 故 $3+3\ln a > 0$ 即 $a > \frac{1}{e}$.
\end{solution}
\end{question}

\begin{question}[points=5]
\textbf{来源：}2021 年 全国乙卷-理科，第 10 题\par
设$a\neq0$，若$x=a$为函数$f(x)=a(x-a)^2(x-b)$的极大值点，则（  ）
\begin{choices}
  \item $a<b$
  \item $a>b$
  \item $ab<a^2$
  \item $ab>a^2$
\end{choices}
\begin{solution}
\textbf{答案：}$ab<a^2$

\textbf{详解：}若$a>0$，其图像如图（1），此时$0<a<b$；若$a<0$，时图像如图（2），此时$b<a<0$。综上，$ab<a^2$。故选D。
\end{solution}
\end{question}

\begin{question}[points=12]
\textbf{来源：}2021 年 全国乙卷-理科，第 20 题\par
设函数$f(x)=\ln(a-x)$，已知$x=0$是函数$y=xf(x)$的极值点。

（1）求$a$；

（2）设函数$g(x)=\frac{x+f(x)}{xf(x)}$，证明：$g(x)<1$。
\begin{solution}
\textbf{答案：}（1）$a=1$；（2）证明见解析。

\textbf{详解：}（1）令$h(x)=xf(x)=x\ln(a-x)$，则$h'(x)=\ln(a-x)-\frac{x}{a-x}$。因为$x=0$是函数$y=xf(x)$的极值点，所以$h'(0)=0$，解得$a=1$。
（2）由（1）可知：$f(x)=\ln(1-x)$。$g(x)=\frac{x+f(x)}{xf(x)}=\frac{1}{f(x)}+\frac{1}{x}$。要证$g(x)<1$，即证$\frac{1}{\ln(1-x)}+\frac{1}{x}<1\Leftrightarrow\frac{1}{\ln(1-x)}+\frac{1-x}{x}<0$（$x<1$且$x\neq0$）$\Leftrightarrow\frac{x+(1-x)\ln(1-x)}{x\ln(1-x)}<0$。因为当$x<0$时，$x\ln(1-x)<0$；当$0<x<1$时，$x\ln(1-x)<0$。所以只需证明$x+(1-x)\ln(1-x)>0$。令$H(x)=x+(1-x)\ln(1-x)$，易知$H(0)=0$。则$H'(x)=1-\ln(1-x)-1=-\ln(1-x)$。当$x<0$时，$H'(x)<0$，$H(x)$在$(-\infty,0)$上单调递减，因为$H(0)=0$，所以$H(x)>0$；当$0<x<1$时，$H'(x)>0$，$H(x)$在$(0,1)$上单调递增，因为$H(0)=0$，所以$H(x)>0$。综上证得$g(x)<1$。
\end{solution}
\end{question}

\begin{question}[points=5]
\textbf{来源：}2021 年 全国乙卷-文科，第 12 题\par
设 $a \neq 0$ ，若 $x = a$ 为函数 $f(x) = a(x-a)^2(x-b)$ 的极大值点，则（ ）
\begin{choices}
  \item $a < b$
  \item $a > b$
  \item $ab < a^2$
  \item $ab > a^2$
\end{choices}
\begin{solution}
\textbf{答案：}D

\textbf{分析：}$f'(x) = 2a(x-a)(x-b) + a(x-a)^2 = a(x-a)(3x-2b-a)$.
当 $a > 0$ 时，原函数先增再减后增.
原函数在 $f'(x)=0$ 的较小零点时取得极大值.
即 $a < \frac{a+2b}{3}$ ，即 $a < b$ ，$\therefore a^2 < ab$.
当 $a < 0$ 时，原函数先减再增后减.
原函数在 $f'(x)=0$ 的较大零点时取得极大值.
即 $a > \frac{a+2b}{3}$ ，$a > b$ ，$a^2 < ab$ ，故选 D.
\end{solution}
\end{question}

\begin{question}[points=12]
\textbf{来源：}2021 年 全国乙卷-文科，第 21 题\par
已知函数 $f(x) = x^3 - x^2 + ax + 1$ .

（1）讨论 $f(x)$ 的单调性；

（2）求曲线 $y = f(x)$ 过坐标原点的切线与曲线 $y = f(x)$ 的公共点的坐标.
\begin{solution}
\textbf{答案：}见解析

\textbf{详解：}（1）$f'(x) = 3x^2 - 2x + a$ .
（i）当 $\Delta = 4 - 12a \le 0$ ，即 $a \ge \frac{1}{3}$ 时，$f'(x) \ge 0$ 恒成立，即 $f(x)$ 在 $\mathbb{R}$ 上单调递增.
（ii）当 $\Delta = 4 - 12a > 0$ ，即 $a < \frac{1}{3}$ 时，$f'(x)=0$ 解得 $x_1 = \frac{1-\sqrt{1-3a}}{3}$ ， $x_2 = \frac{1+\sqrt{1-3a}}{3}$ .
所以 $f(x)$ 在 $(-\infty, \frac{1-\sqrt{1-3a}}{3})$ 和 $(\frac{1+\sqrt{1-3a}}{3}, +\infty)$ 单调递增，在 $(\frac{1-\sqrt{1-3a}}{3}, \frac{1+\sqrt{1-3a}}{3})$ 单调递减.
综上所述：当 $a \ge \frac{1}{3}$ 时，$f(x)$ 在 $\mathbb{R}$ 上单调递增；当 $a < \frac{1}{3}$ 时，$f(x)$ 在 $(\frac{1-\sqrt{1-3a}}{3}, \frac{1+\sqrt{1-3a}}{3})$ 单调递减，在其余区间单调递增.
（2）设过原点切线的切点为 $(t, t^3 - t^2 + at + 1)$ ，切线斜率 $k = f'(t) = 3t^2 - 2t + a$ .
又 $k = \frac{t^3 - t^2 + at + 1}{t}$ ，可得 $\frac{t^3 - t^2 + at + 1}{t} = 3t^2 - 2t + a$ .
化简得 $(t-1)(2t^2 + t + 1) = 0$ ，即 $t = 1$ .
所以切点为 $(1, a+1)$ ，斜率 $k = a+1$ ，切线方程为 $y = (a+1)x$ .
将 $y = (a+1)x$ 与 $y = x^3 - x^2 + ax + 1$ 联立得 $x^3 - x^2 + ax + 1 = (a+1)x$ ，
化简得 $(x-1)^2(x+1) = 0$ ，解得 $x_1 = 1$ ， $x_2 = -1$ .
所以过原点的切线与 $y = f(x)$ 公共点坐标为 $(1, a+1)$ ， $(-1, -a-1)$ .
\end{solution}
\end{question}

\begin{question}[points=16]
\textbf{来源：}2021 年 天津卷，第 20 题\par
已知 $a > 0$，函数 $f(x) = ax - xe^x$．

（I）求曲线 $y = f(x)$ 在点 $(0, f(0))$ 处的切线方程；

（II）证明 $f(x)$ 存在唯一的极值点；

（III）若存在 $a$，使得 $f(x) \le a + b$ 对任意 $x \in \mathbb{R}$ 成立，求实数 $b$ 的取值范围．
\begin{solution}
\textbf{答案：}（I）$y = (a-1)x$；（II）证明见解析；（III）$[-e, +\infty)$

\textbf{分析：}（I）求出 $f(x)$ 在 $x=0$ 处的导数，即切线斜率，求出 $f(0)$，即可求出切线方程；
（II）令 $f'(x)=0$，可得 $a = (x+1)e^x$，则可化为证明 $y=a$ 与 $y=g(x)$ 仅有一个交点，利用导数求出 $g(x)$ 的变化情况，数形结合即可求解；
（III）令 $h(x) = (x^2 - x - 1)e^x$，$x > -1$，题目等价于存在 $x \in (-1, +\infty)$，使得 $h(x) \le b$，即 $b \ge h(x)_{\min}$，利用导数即可求出 $h(x)$ 的最小值.

\textbf{详解：}（I）$f'(x) = a - (x+1)e^x$，$f'(0) = a-1$，$f(0)=0$，切线方程为 $y = (a-1)x$；
（II）令 $f'(x)=0$ 得 $a = (x+1)e^x$，令 $g(x) = (x+1)e^x$，$g'(x) = (x+2)e^x$，当 $x \in (-\infty, -2)$ 时 $g'(x)<0$，$g(x)$ 递减；当 $x \in (-2, +\infty)$ 时 $g'(x)>0$，$g(x)$ 递增，$x \to -\infty$ 时 $g(x) \to 0^-$，$g(-1)=0$，$x \to +\infty$ 时 $g(x) \to +\infty$，故 $y=a$ 与 $y=g(x)$ 仅有一个交点，设交点为 $(m, a)$，则 $m > -1$，且 $f'(m)=0$，当 $x < m$ 时 $f'(x)>0$，$f(x)$ 递增；当 $x > m$ 时 $f'(x)<0$，$f(x)$ 递减，所以 $x=m$ 为极大值点，故 $f(x)$ 存在唯一的极值点；
（III）由（II）知 $f(x)_{\max} = f(m)$，且 $a = (m+1)e^m$，$m>-1$，则 $f(m)-a = (m^2 - m -1)e^m$，令 $h(x) = (x^2 - x - 1)e^x$，$x>-1$，则 $h'(x) = (x^2 + x -2)e^x = (x-1)(x+2)e^x$，当 $x \in (-1,1)$ 时 $h'(x)<0$，$h(x)$ 递减；当 $x \in (1, +\infty)$ 时 $h'(x)>0$，$h(x)$ 递增，$h(x)_{\min} = h(1) = -e$，存在 $a$ 使 $f(x) \le a+b$ 对任意 $x$ 成立等价于存在 $x>-1$ 使 $h(x) \le b$，即 $b \ge h(x)_{\min} = -e$，所以 $b$ 的取值范围是 $[-e, +\infty)$.
\end{solution}
\end{question}

\begin{question}[points=5]
\textbf{来源：}2021 年 新高考II卷，第 16 题\par
已知函数$f(x)=e^x-1$，$x_1<0,x_2>0$，函数$f(x)$的图象在点$A(x_1,f(x_1))$和点$B(x_2,f(x_2))$的两条切线互相垂直，且分别交$y$轴于$M$，$N$两点，则$\frac{|AM|}{|BN|}$取值范围是．
\fillin[{$(0,1)$}]
\begin{solution}
\textbf{答案：}$(0,1)$

\textbf{分析：}结合导数的几何意义可得$x_1+x_2=0$，结合直线方程及两点间距离公式可得$|AM|=\sqrt{1+e^{2x_1}}\cdot|x_1|$，$|BN|=\sqrt{1+e^{2x_2}}\cdot|x_2|$，化简即可得解.

\textbf{详解：}由题意，$f(x)=e^x-1=\begin{cases} 1-e^x, & x<0 \\ e^x-1, & x\ge0 \end{cases}$，则$f'(x)=\begin{cases} -e^x, & x<0 \\ e^x, & x>0 \end{cases}$，所以点$A(x_1,1-e^{x_1})$和点$B(x_2,e^{x_2}-1)$，$k_{AM}=-e^{x_1},k_{BN}=e^{x_2}$，所以$-e^{x_1}\cdot e^{x_2}=-1,x_1+x_2=0$，所以$AM: y-1+e^{x_1}=-e^{x_1}(x-x_1)$，$M(0,e^{x_1}x_1-e^{x_1}+1)$，所以$|AM|=\sqrt{x_1^2+(e^{x_1}x_1)^2}=\sqrt{1+e^{2x_1}}\cdot|x_1|$，同理$|BN|=\sqrt{1+e^{2x_2}}\cdot|x_2|$，所以$\frac{|AM|}{|BN|}=\frac{\sqrt{1+e^{2x_1}}\cdot|x_1|}{\sqrt{1+e^{2x_2}}\cdot|x_2|}=\frac{\sqrt{1+e^{2x_1}}}{\sqrt{1+e^{2x_2}}}=\frac{\sqrt{1+e^{2x_1}}}{\sqrt{1+e^{-2x_1}}}=e^{x_1}\in(0,1)$.故答案为：$(0,1)$.
\end{solution}
\end{question}

\begin{question}[points=12]
\textbf{来源：}2021 年 新高考II卷，第 22 题\par
已知函数$f(x)=(x-1)e^x-ax^2+b$．

（1）讨论$f(x)$的单调性；

（2）从下面两个条件中选一个，证明：$f(x)$有一个零点

① $\frac{1}{2}<a\leq\frac{e^2}{2}, b>2a$；

② $0<a<\frac{1}{2}, b\leq2a$．
\begin{solution}
\textbf{答案：}(1) 答案见解析；(2) 证明见解析.

\textbf{分析：}(1)首先求得导函数的解析式，然后分类讨论确定函数的单调性即可； (2)由题意结合(1)中函数的单调性和函数零点存在定理即可证得题中的结论.

\textbf{详解：}(1)由函数的解析式可得：$f'(x)=x(e^x-2a)$，当$a\leq0$时，若$x\in(-\infty,0)$，则$f'(x)<0,f(x)$单调递减，若$x\in(0,+\infty)$，则$f'(x)>0,f(x)$单调递增；当$0<a<\frac{1}{2}$时，若$x\in(-\infty,\ln(2a))$，则$f'(x)>0,f(x)$单调递增，若$x\in(\ln(2a),0)$，则$f'(x)<0,f(x)$单调递减，若$x\in(0,+\infty)$，则$f'(x)>0,f(x)$单调递增；当$a=\frac{1}{2}$时，$f'(x)\geq0,f(x)$在$\mathbb{R}$上单调递增；当$a>\frac{1}{2}$时，若$x\in(-\infty,0)$，则$f'(x)>0,f(x)$单调递增，若$x\in(0,\ln(2a))$，则$f'(x)<0,f(x)$单调递减，若$x\in(\ln(2a),+\infty)$，则$f'(x)>0,f(x)$单调递增； (2)若选择条件①：由于$\frac{1}{2}<a\leq\frac{e^2}{2}$，故$1<2a\leq e^2$，则$b>2a>1,f(0)=b-1>0$，而$f(-b)=(-1-b)e^{-b}-ab^2-b<0$，而函数在区间$(-\infty,0)$上单调递增，故函数在区间$(-\infty,0)$上有一个零点．$f(\ln(2a))=2a[\ln(2a)-1]-a[\ln(2a)]^2+b>2a[\ln(2a)-1]-a[\ln(2a)]^2+2a=2a\ln(2a)-a[\ln(2a)]^2=a\ln(2a)[2-\ln(2a)]$，由于$\frac{1}{2}<a\leq\frac{e^2}{2}$，$1<2a\leq e^2$，故$a\ln(2a)[2-\ln(2a)]\geq0$，结合函数的单调性可知函数在区间$(0,+\infty)$上没有零点．综上可得，题中的结论成立．若选择条件②：由于$0<a<\frac{1}{2}$，故$2a<1$，则$f(0)=b-1\leq2a-1<0$，当$b\geq0$时，$e^2>4,4a<2$，$f(2)=e^2-4a+b>0$，而函数在区间$(0,+\infty)$上单调递增，故函数在区间$(0,+\infty)$上有一个零点．当$b<0$时，构造函数$H(x)=e^x-x-1$，则$H'(x)=e^x-1$，当$x\in(-\infty,0)$时，$H'(x)<0,H(x)$单调递减，当$x\in(0,+\infty)$时，$H'(x)>0,H(x)$单调递增，注意到$H(0)=0$，故$H(x)\geq0$恒成立，从而有：$e^x\geq x+1$，此时：$f(x)=(x-1)e^x-ax^2+b\geq(x-1)(x+1)-ax^2+b=(1-a)x^2+(b-1)$，当$x>\sqrt{\frac{1-b}{1-a}}$时，$(1-a)x^2+(b-1)>0$，取$x_0=\sqrt{\frac{1-b}{1-a}}+1$，则$f(x_0)>0$，即：$f(0)<0,f(\sqrt{\frac{1-b}{1-a}}+1)>0$，而函数在区间$(0,+\infty)$上单调递增，故函数在区间$(0,+\infty)$上有一个零点．$f(\ln(2a))=2a[\ln(2a)-1]-a[\ln(2a)]^2+b\leq2a[\ln(2a)-1]-a[\ln(2a)]^2+2a=2a\ln(2a)-a[\ln(2a)]^2=a\ln(2a)[2-\ln(2a)]$，由于$0<a<\frac{1}{2}$，$0<2a<1$，故$a\ln(2a)[2-\ln(2a)]<0$，结合函数的单调性可知函数在区间$(-\infty,0)$上没有零点．综上可得，题中的结论成立．
\end{solution}
\end{question}

\begin{question}[points=5]
\textbf{来源：}2021 年 新高考I卷，第 7 题\par
若过点 $(a,b)$ 可以作曲线 $y = e^x$ 的两条切线, 则 ( )
\begin{choices}
  \item $e^b < a$
  \item $e^a < b$
  \item $0 < a < e^b$
  \item $0 < b < e^a$
\end{choices}
\begin{solution}
\textbf{答案：}$0 < b < e^a$

\textbf{分析：}根据导数几何意义求得切线方程, 再构造函数, 利用导数研究函数图象, 结合图形确定结果.

\textbf{详解：}在曲线 $y = e^x$ 上任取一点 $P(t, e^t)$, 对函数求导得 $y' = e^x$, 曲线在点 $P$ 处的切线方程为 $y - e^t = e^t(x - t)$, 即 $y = e^t x + (1 - t)e^t$. 由题意可知点 $(a,b)$ 在切线上, 得 $b = ae^t + (1 - t)e^t = (a+1-t)e^t$. 令 $f(t) = (a+1-t)e^t$, 则 $f'(t) = (a - t)e^t$. 当 $t < a$ 时 $f'(t) > 0$, $f(t)$ 递增; 当 $t > a$ 时 $f'(t) < 0$, $f(t)$ 递减, 所以 $f(t)_{\max} = f(a) = e^a$. 由题意, 直线 $y = b$ 与曲线 $y = f(t)$ 有两个交点, 则 $b < e^a$. 当 $t < a+1$ 时 $f(t) > 0$, 当 $t > a+1$ 时 $f(t) < 0$, 结合图象知 $0 < b < e^a$, 故选D.
\end{solution}
\end{question}

\begin{question}[points=12]
\textbf{来源：}2021 年 新高考I卷，第 22 题\par
已知函数 $f(x) = x(1 - \ln x)$.

(1) 讨论 $f(x)$ 的单调性;

(2) 设 $a$, $b$ 为两个不相等的正数, 且 $b\ln a - a\ln b = a - b$, 证明: $2 < \frac{1}{a} + \frac{1}{b} < e$.
\begin{solution}
\textbf{答案：}(1) $f(x)$ 的递增区间为 $(0,1)$, 递减区间为 $(1,+\infty)$; (2) 证明见解析.

\textbf{分析：}(1) 求导判断符号. (2) 变形为 $\frac{1}{a}(1-\ln\frac{1}{a}) = \frac{1}{b}(1-\ln\frac{1}{b})$, 设 $x_1 = \frac{1}{a}$, $x_2 = \frac{1}{b}$, 则 $f(x_1) = f(x_2)$, 利用极值点偏移证明 $x_1 + x_2 > 2$, 再构造函数证明 $x_1 + x_2 < e$.

\textbf{详解：}(1) $f'(x) = -\ln x$, 当 $x \in (0,1)$ 时 $f'(x) > 0$, 当 $x \in (1,+\infty)$ 时 $f'(x) < 0$, 故 $f(x)$ 在 $(0,1)$ 递增, 在 $(1,+\infty)$ 递减.
(2) 由 $b\ln a - a\ln b = a - b$ 得 $\frac{\ln a + 1}{a} = \frac{\ln b + 1}{b}$, 即 $\frac{1}{a}(1-\ln\frac{1}{a}) = \frac{1}{b}(1-\ln\frac{1}{b})$. 令 $x_1 = \frac{1}{a}$, $x_2 = \frac{1}{b}$, 则 $f(x_1) = f(x_2)$, 不妨设 $0 < x_1 < 1 < x_2$. 先证 $x_1 + x_2 > 2$: 若 $x_2 \geq 2$, 显然成立; 若 $1 < x_2 < 2$, 设 $g(x) = f(x) - f(2-x)$ $(1 < x < 2)$, $g'(x) = -\ln[x(2-x)] > 0$, $g(x)$ 递增, $g(x) > g(1) = 0$, 故 $f(x_2) > f(2-x_2)$, 即 $f(x_1) > f(2-x_2)$, 由单调性得 $x_1 < 2 - x_2$, 即 $x_1 + x_2 < 2$? 需要印证: 实际上 $f(x_1)=f(x_2)$, 由 $f(x_2) > f(2-x_2)$ 得 $f(x_1) > f(2-x_2)$, 因为 $0<x_1<1$, $0<2-x_2<1$, 且 $f$ 在 $(0,1)$ 递增, 所以 $x_1 > 2 - x_2$, 即 $x_1 + x_2 > 2$. 再证 $x_1 + x_2 < e$: 设 $x_2 = t x_1 (t>1)$, 由 $f(x_1)=f(x_2)$ 得 $\ln x_1 = \frac{t-1 - t\ln t}{t-1}$. 要证 $x_1 + x_2 < e$, 即 $\ln x_1 + \ln(t+1) < 1$, 代入得 $\frac{t-1 - t\ln t}{t-1} + \ln(t+1) < 1$, 即 $(t-1)\ln(t+1) - t\ln t < 0$. 令 $S(t) = (t-1)\ln(t+1) - t\ln t$, $S'(t) = \ln(1+\frac{1}{t}) - \frac{2}{t+1}$. 利用不等式 $\ln(1+x) \leq x$ 得 $\ln(1+\frac{1}{t}) < \frac{1}{t} < \frac{2}{t+1}$ (当 $t>1$), 所以 $S'(t) < 0$, $S(t)$ 递减, $S(t) < S(1) = 0$, 故 $x_1 + x_2 < e$ 成立. 综上, $2 < \frac{1}{a} + \frac{1}{b} < e$.
\end{solution}
\end{question}

\begin{question}[points=15]
\textbf{来源：}2021 年 浙江卷，第 22 题\par
设 $a,b$ 为实数, 且 $a>1$, 函数 $f(x)=a^x - bx + e^2$ $(x\in\mathbb{R})$.

(1) 求函数 $f(x)$ 的单调区间；

(2) 若对任意 $b>2e^2$, 函数 $f(x)$ 有两个不同的零点, 求 $a$ 的取值范围；

(3) 当 $a=e$ 时, 证明：对任意 $b>e^4$, 函数 $f(x)$ 有两个不同的零点 $x_1,x_2$, 满足 $x_2 > \frac{b\ln b}{2e^2} x_1 + \frac{e^2}{b}$.

(注：$e=2.71828\cdots$ 是自然对数的底数)
\begin{solution}
\textbf{答案：}(1) 见解析; (2) $(1,e^2]$; (3) 见解析

\textbf{分析：}(1) 求导, 讨论 $b$ 的正负; (2) 转化为方程有解, 利用导数研究函数; (3) 利用 (2) 的结论, 结合放缩法证明.

\textbf{详解：}(1) $f'(x)=a^x\ln a - b$. 若 $b\leq0$, 则 $f'(x)>0$, $f(x)$ 在 $\mathbb{R}$ 上单调递增；若 $b>0$, 则当 $x<\log_a\frac{b}{\ln a}$ 时 $f'(x)<0$, $f(x)$ 单调递减；当 $x>\log_a\frac{b}{\ln a}$ 时 $f'(x)>0$, $f(x)$ 单调递增.
(2) $f(x)=0$ 有两个不同零点等价于 $e^{x\ln a} - bx + e^2=0$ 有两个不同解. 令 $t=x\ln a$, 则 $e^t - \frac{b}{\ln a}t + e^2=0$, 即 $\frac{b}{\ln a} = \frac{e^t+e^2}{t}$. 设 $g(t)=\frac{e^t+e^2}{t}$, $g'(t)=\frac{e^t(t-1)-e^2}{t^2}$, 记 $h(t)=e^t(t-1)-e^2$, $h'(t)=e^t t>0$, 且 $h(2)=0$, 故 $g(t)$ 在 $(0,2)$ 递减, $(2,+\infty)$ 递增, $g(2)=e^2$. 所以 $\frac{b}{\ln a} > e^2$, 即 $\ln a < \frac{b}{e^2}$. 对任意 $b>2e^2$, 需 $\frac{b}{e^2}>2$, 则 $\ln a \leq 2$, 故 $1<a\leq e^2$.
(3) $a=e$ 时, $f(x)=e^x - bx + e^2$. 由 (2) 知有两正根 $x_1<x_2$, 且 $b=\frac{e^{x_1}+e^2}{x_1} = \frac{e^{x_2}+e^2}{x_2} > e^4$. 由 $b>e^4$ 得 $x_2>5$. 由 $b=\frac{e^{x_1}+e^2}{x_1} < \frac{2e^2}{x_1}$ (因为 $e^{x_1}<e^2$ 故不严谨, 实际上 $x_1<2$, 所以 $e^{x_1}<e^2$ 不成立, 需另证). 由 $b=\frac{e^{x_1}+e^2}{x_1} > \frac{e^2}{x_1}$ 得 $x_1 > \frac{e^2}{b}$, 但我们需要 $x_1<\frac{2e^2}{b}$? 事实上, 由 $b=\frac{e^{x_2}+e^2}{x_2} > \frac{e^{x_2}}{x_2}$ 可得 $\ln b > x_2 - \ln x_2$, 即 $x_2 < \ln b + \ln x_2$. 要证 $x_2 > \frac{b\ln b}{2e^2}x_1 + \frac{e^2}{b}$, 先由 $x_1 < \frac{2e^2}{b}$? 实际上由 $b=\frac{e^{x_1}+e^2}{x_1}$ 且 $x_1<2$, $e^{x_1}<e^2$, 得 $b>\frac{2e^2}{x_1}$, 故 $x_1<\frac{2e^2}{b}$. 则右边 $< \frac{b\ln b}{2e^2}\cdot\frac{2e^2}{b}+\frac{e^2}{b} = \ln b + \frac{e^2}{b}$. 而 $b = \frac{e^{x_2}+e^2}{x_2} < \frac{2e^{x_2}}{x_2}$ (因为 $e^2<e^{x_2}$ 当 $x_2>2$), 故 $\ln b < x_2 + \ln2 - \ln x_2$, 只需证 $x_2 > x_2 + \ln2 - \ln x_2 + \frac{e^2}{b}$, 即 $\ln x_2 > \ln2 + \frac{e^2}{b}$, 当 $b>e^4$ 时 $\frac{e^2}{b}<\frac{1}{e^2}$, $\ln x_2 > \ln5 > \ln2 + \frac{1}{e^2}$ 成立. 故原不等式成立.
\end{solution}
\end{question}

\section*{2022 年}

\begin{question}[points=15]
\textbf{来源：}2022 年 北京卷，第 20 题\par
已知函数 $f(x) = e^x \ln(1+x)$。

（1）求曲线 $y=f(x)$ 在点 $(0,f(0))$ 处的切线方程；

（2）设 $g(x)=f'(x)$，讨论函数 $g(x)$ 在 $[0,+\infty)$ 上的单调性；

（3）证明：对任意的 $s,t\in (0,+\infty)$，有 $f(s+t) > f(s)+f(t)$。
\begin{solution}
\textbf{答案：}（1）$y=x$；（2）$g(x)$ 在 $[0,+\infty)$ 上单调递增；（3）证明见解析

\textbf{分析：}（1）先求出切点坐标，再由导数求得切线斜率，即得切线方程；（2）在求一次导数无法判断的情况下，构造新的函数，再求一次导数，问题即得解；（3）令 $m(x)=f(x+t)-f(x)$，$(x,t>0)$，即证 $m(x)>m(0)$，由第二问结论可知 $m(x)$ 在 $[0,+\infty)$ 上单调递增，即得证。

\textbf{详解：}（1）解：因为 $f(x)=e^x\ln(1+x)$，所以 $f(0)=0$，即切点坐标为 $(0,0)$。又 $f'(x)=e^x\left(\ln(1+x)+\dfrac{1}{1+x}\right)$，切线斜率 $k=f'(0)=1$，所以切线方程为 $y=x$。
（2）解：$g(x)=f'(x)=e^x\left(\ln(1+x)+\dfrac{1}{1+x}\right)$，$g'(x)=e^x\left(\ln(1+x)+\dfrac{1}{1+x}\right)+e^x\left(\dfrac{1}{1+x}-\dfrac{1}{(1+x)^2}\right)=e^x\left(\ln(1+x)+\dfrac{2}{1+x}-\dfrac{1}{(1+x)^2}\right)$。
令 $h(x)=\ln(1+x)+\dfrac{2}{1+x}-\dfrac{1}{(1+x)^2}$，则 $h'(x)=\dfrac{1}{1+x}-\dfrac{2}{(1+x)^2}+\dfrac{2}{(1+x)^3}=\dfrac{(1+x)^2-2(1+x)+2}{(1+x)^3}=\dfrac{x^2+1}{(1+x)^3}>0$，所以 $h(x)$ 在 $[0,+\infty)$ 上单调递增，故 $h(x)\ge h(0)=1>0$，所以 $g'(x)>0$ 在 $[0,+\infty)$ 上恒成立，故 $g(x)$ 在 $[0,+\infty)$ 上单调递增。
（3）证明：原不等式等价于 $f(s+t)-f(s)>f(t)-f(0)$。令 $m(x)=f(x+t)-f(x)$，$(x,t>0)$，即证 $m(x)>m(0)$。因为 $m'(x)=f'(x+t)-f'(x)=g(x+t)-g(x)$，由（2）知 $g(x)$ 在 $[0,+\infty)$ 上单调递增，且 $x+t>x$，所以 $g(x+t)>g(x)$，故 $m'(x)>0$，所以 $m(x)$ 在 $(0,+\infty)$ 上单调递增。又因为 $x>0$，所以 $m(x)>m(0)$，即 $f(x+t)-f(x)>f(t)-f(0)=f(t)$，所以 $f(s+t)>f(s)+f(t)$。
\end{solution}
\end{question}

\begin{question}[points=5]
\textbf{来源：}2022 年 全国甲卷-理科，第 6 题\par
当 $x=1$ 时，函数 $f(x)=a\ln x + \frac{b}{x}$ 取得最大值 $-2$，则 $f'(2)=$（ ）
\begin{choices}
  \item $-1$
  \item $-\frac{1}{2}$
  \item $\frac{1}{2}$
  \item $1$
\end{choices}
\begin{solution}
\textbf{答案：}B

\textbf{分析：}根据题意可知 $f(1)=-2$, $f'(1)=0$ 即可解得 $a,b$，再根据 $f'(x)$ 即可解出.

\textbf{详解：}因为函数 $f(x)$ 定义域为 $(0,+\infty)$，所以依题可知 $f(1)=-2$, $f'(1)=0$，而 $f'(x)=\frac{a}{x}-\frac{b}{x^2}$，所以 $b=-2$, $a-b=0$，即 $a=-2$, $b=-2$，所以 $f'(x)=-\frac{2}{x}+\frac{2}{x^2}$，因此函数 $f(x)$ 在 $(0,1)$ 上递增，在 $(1,+\infty)$ 上递减，$x=1$ 时取最大值，满足题意，即有 $f'(2)=-1+\frac{1}{2}=-\frac{1}{2}$. 故选B.
\end{solution}
\end{question}

\begin{question}[points=12]
\textbf{来源：}2022 年 全国甲卷-理科，第 21 题\par
已知函数 $f(x)=\frac{e^x}{x}-\ln x+x-a$.

(1) 若 $f(x)\geq 0$，求a的取值范围;

(2) 证明：若 $f(x)$ 有两个零点 $x_1, x_2$，则 $x_1x_2<1$.
\begin{solution}
\textbf{答案：}(1) $(-\infty, e+1]$；(2) 证明见解析

\textbf{分析：}(1) 由导数确定函数单调性及最值，即可得解. (2) 利用分析法，转化要证明条件为 $\frac{e^x}{x}-xe^{\frac{1}{x}}-2\left[\ln x-\frac{1}{2}\left(x-\frac{1}{x}\right)\right]>0$，再利用导数即可得证.

\textbf{详解：}(1) $f(x)$ 的定义域为 $(0,+\infty)$，则 $f'(x)=\left(\frac{1}{x}-\frac{1}{x^2}\right)e^x-\frac{1}{x}+1=\frac{x-1}{x}\left(\frac{e^x}{x}+1\right)$，令 $f'(x)=0$ 得 $x=1$，当 $x\in(0,1)$ 时 $f'(x)<0$, $f(x)$ 单调递减，当 $x\in(1,+\infty)$ 时 $f'(x)>0$, $f(x)$ 单调递增，$f(x)\geq f(1)=e+1-a$，若 $f(x)\geq0$，则 $e+1-a\geq0$，即 $a\leq e+1$，所以a的取值范围为 $(-\infty, e+1]$.
(2) 由题知，$f(x)$ 一个零点小于1，一个零点大于1，不妨设 $x_1<1<x_2$，要证 $x_1x_2<1$，即证 $x_1<\frac{1}{x_2}$，因为 $x_1,\frac{1}{x_2}\in(0,1)$，即证 $f(x_1)>f\left(\frac{1}{x_2}\right)$，又因为 $f(x_1)=f(x_2)$，故只需证 $f(x_2)>f\left(\frac{1}{x_2}\right)$，即证 $\frac{e^{x_2}}{x_2}-\ln x_2+x_2-\left(\frac{e^{\frac{1}{x_2}}}{\frac{1}{x_2}}-\ln\frac{1}{x_2}+\frac{1}{x_2}\right)>0$，即证 $\frac{e^{x_2}}{x_2}-x_2 e^{\frac{1}{x_2}}-2\left[\ln x_2-\frac{1}{2}\left(x_2-\frac{1}{x_2}\right)\right]>0$, $x_2>1$. 下面证明 $x>1$ 时，$\frac{e^x}{x}-xe^{\frac{1}{x}}>0$ 且 $\ln x-\frac{1}{2}\left(x-\frac{1}{x}\right)<0$. 设 $g(x)=\frac{e^x}{x}-xe^{\frac{1}{x}}$, $x>1$，则 $g'(x)=\left(\frac{1}{x}-\frac{1}{x^2}\right)e^x-\left(e^{\frac{1}{x}}+xe^{\frac{1}{x}}\cdot(-\frac{1}{x^2})\right)=\frac{x-1}{x}\left(\frac{e^x}{x}-e^{\frac{1}{x}}\right)$，设 $\varphi(x)=\frac{e^x}{x}\ (x>1)$，$\varphi'(x)=\frac{x-1}{x^2}e^x>0$，所以 $\varphi(x)>\varphi(1)=e$，而 $e^{\frac{1}{x}}<e$，所以 $\frac{e^x}{x}-e^{\frac{1}{x}}>0$，故 $g'(x)>0$，$g(x)$ 在 $(1,+\infty)$ 单调递增，$g(x)>g(1)=0$，所以 $\frac{e^x}{x}-xe^{\frac{1}{x}}>0$. 令 $h(x)=\ln x-\frac{1}{2}\left(x-\frac{1}{x}\right)$, $x>1$，则 $h'(x)=\frac{1}{x}-\frac{1}{2}\left(1+\frac{1}{x^2}\right)=\frac{2x-x^2-1}{2x^2}=\frac{-(x-1)^2}{2x^2}<0$，所以 $h(x)$ 在 $(1,+\infty)$ 单调递减，$h(x)<h(1)=0$，所以 $\ln x-\frac{1}{2}\left(x-\frac{1}{x}\right)<0$. 综上，$\frac{e^x}{x}-xe^{\frac{1}{x}}-2\left[\ln x-\frac{1}{2}\left(x-\frac{1}{x}\right)\right]>0$，所以 $x_1x_2<1$.
\end{solution}
\end{question}

\begin{question}[points=5]
\textbf{来源：}2022 年 全国甲卷-文科，第 8 题\par
当 $x=1$ 时，函数 $f(x)=a\ln x+\frac{b}{x}$ 取得最大值 $-2$，则 $f'(2)=$
\begin{choices}
  \item $-1$
  \item $-\frac{1}{2}$
  \item $\frac{1}{2}$
  \item $1$
\end{choices}
\begin{solution}
\textbf{答案：}B
\end{solution}
\end{question}

\begin{question}[points=12]
\textbf{来源：}2022 年 全国甲卷-文科，第 20 题\par
已知函数 $f(x)=x^3-x, g(x)=x^2+a$，曲线 $y=f(x)$ 在点 $(x_1, f(x_1))$ 处的切线也是曲线 $y=g(x)$ 的切线．

（1）若 $x_1=-1$，求 $a$；

（2）求 $a$ 的取值范围．
\begin{solution}
\textbf{答案：}（1）$a=3$；（2）$[-1,+\infty)$
\end{solution}
\end{question}

\begin{question}[points=5]
\textbf{来源：}2022 年 全国乙卷-理科，第 16 题\par
已知$x=x_1$和$x=x_2$分别是函数$f(x)=2a^x - ex^2$（$a>0$且$a\neq1$）的极小值点和极大值点．若$x_1<x_2$，则$a$的取值范围是．
\fillin[{$\left(\frac{1}{e},1\right)$}]
\begin{solution}
\textbf{答案：}$\left(\frac{1}{e},1\right)$

\textbf{分析：}由$x_1,x_2$分别是函数$f(x)=2a^x-ex^2$的极小值点和极大值点，可得$x\in(-\infty,x_1)\cup(x_2,+\infty)$时，$f'(x)<0$，$x\in(x_1,x_2)$时，$f'(x)>0$，再分$a>1$和$0<a<1$两种情况讨论，方程$2\ln a\cdot a^x - 2ex = 0$的两个根为$x_1,x_2$，即函数$y=\ln a\cdot a^x$与函数$y=ex$的图象有两个不同的交点，构造函数$g(x)=\ln a\cdot a^x$，根据导数的结合意义结合图象即可得出答案.

\textbf{详解：}解：$f'(x)=2\ln a\cdot a^x - 2ex$，因为$x_1,x_2$分别是函数$f(x)=2a^x-ex^2$的极小值点和极大值点，所以函数$f(x)$在$(-\infty,x_1)$和$(x_2,+\infty)$上递减，在$(x_1,x_2)$上递增，所以当$x\in(-\infty,x_1)\cup(x_2,+\infty)$时，$f'(x)<0$，当$x\in(x_1,x_2)$时，$f'(x)>0$，若$a>1$时，当$x<0$时，$2\ln a\cdot a^x>0, 2ex<0$，则此时$f'(x)>0$，与前面矛盾，故$a>1$不符合题意，若$0<a<1$时，则方程$2\ln a\cdot a^x - 2ex = 0$的两个根为$x_1,x_2$，即方程$\ln a\cdot a^x = ex$的两个根为$x_1,x_2$，即函数$y=\ln a\cdot a^x$与函数$y=ex$的图象有两个不同的交点，令$g(x)=\ln a\cdot a^x$，则$g'(x)=\ln^2 a\cdot a^x, 0<a<1$，设过原点且与函数$y=g(x)$的图象相切的直线的切点为$(x_0,\ln a\cdot a^{x_0})$，则切线的斜率为$g'(x_0)=\ln^2 a\cdot a^{x_0}$，故切线方程为$y-\ln a\cdot a^{x_0}=\ln^2 a\cdot a^{x_0}(x-x_0)$，则有$-\ln a\cdot a^{x_0} = -x_0\ln^2 a\cdot a^{x_0}$，解得$x_0=\frac{1}{\ln a}$，则切线的斜率为$\ln^2 a\cdot a^{\frac{1}{\ln a}} = e\ln^2 a$，因为函数$y=\ln a\cdot a^x$与函数$y=ex$的图象有两个不同的交点，所以$e\ln^2 a < e$，解得$\frac{1}{e}<a<e$，又$0<a<1$，所以$\frac{1}{e}<a<1$，综上所述，$a$的范围为$\left(\frac{1}{e},1\right)$.
\end{solution}
\end{question}

\begin{question}[points=12]
\textbf{来源：}2022 年 全国乙卷-理科，第 21 题\par
已知函数$f(x)=\ln(1+x)+axe^{-x}$．

(1)当$a=1$时，求曲线$y=f(x)$在点$(0,f(0))$处的切线方程；

(2)若$f(x)$在区间$(-1,0),(0,+\infty)$各恰有一个零点，求$a$的取值范围．
\begin{solution}
\textbf{答案：}(1) $y=2x$
(2) $(-\infty,-1)$

\textbf{分析：}(1)先算出切点,再求导算出斜率即可；(2)求导,对$a$分类讨论,对$x$分$(-1,0),(0,+\infty)$两部分研究

\textbf{详解：}【小问1详解】\\$f(x)$的定义域为$(-1,+\infty)$当$a=1$时,$f(x)=\ln(1+x)+\frac{x}{e^x}, f(0)=0$,所以切点为$(0,0)$\\$f'(x)=\frac{1}{1+x}+\frac{1-x}{e^x}, f'(0)=2$,所以切线斜率为2所以曲线$y=f(x)$在点$(0,f(0))$处的切线方程为$y=2x$\\【小问2详解】\\$f(x)=\ln(1+x)+\frac{ax}{e^x}$\\$f'(x)=\frac{1}{1+x}+\frac{a(1-x)}{e^x}=\frac{e^x+a(1-x^2)}{(1+x)e^x}$设$g(x)=e^x+a(1-x^2)$\\$1^\circ$若$a>0$,当$x\in(-1,0), g(x)=e^x+a(1-x^2)>0$,即$f'(x)>0$所以$f(x)$在$(-1,0)$上单调递增,$f(x)<f(0)=0$故$f(x)$在$(-1,0)$上没有零点,不合题意\\$2^\circ$若$-1\leq a\leq 0$,当$x\in(0,+\infty)$,则$g'(x)=e^x-2ax>0$所以$g(x)$在$(0,+\infty)$上单调递增所以$g(x)>g(0)=1+a\geq0$,即$f'(x)>0$所以$f(x)$在$(0,+\infty)$上单调递增,$f(x)>f(0)=0$故$f(x)$在$(0,+\infty)$上没有零点,不合题意\\$3^\circ$若$a<-1$\\(1)当$x\in(0,+\infty)$,则$g'(x)=e^x-2ax>0$,所以$g(x)$在$(0,+\infty)$上单调递增\\$g(0)=1+a<0, g(1)=e>0$所以存在$m\in(0,1)$,使得$g(m)=0$,即$f'(m)=0$当$x\in(0,m), f'(x)<0, f(x)$单调递减当$x\in(m,+\infty), f'(x)>0, f(x)$单调递增所以当$x\in(0,m), f(x)<f(0)=0$当$x\to+\infty, f(x)\to+\infty$所以$f(x)$在$(m,+\infty)$上有唯一零点又$(0,m)$没有零点,即$f(x)$在$(0,+\infty)$上有唯一零点\\(2)当$x\in(-1,0), g(x)=e^x+a(1-x^2)$设$h(x)=g'(x)=e^x-2ax$\\$h'(x)=e^x-2a>0$所以$g'(x)$在$(-1,0)$单调递增\\$g'(-1)=\frac{1}{e}+2a<0, g'(0)=1>0$所以存在$n\in(-1,0)$,使得$g'(n)=0$当$x\in(-1,n), g'(x)<0, g(x)$单调递减当$x\in(n,0), g'(x)>0, g(x)$单调递增$g(x)<g(0)=1+a<0$又$g(-1)=\frac{1}{e}>0$所以存在$t\in(-1,n)$,使得$g(t)=0$,即$f'(t)=0$当$x\in(-1,t), f(x)$单调递增,当$x\in(t,0), f(x)$单调递减有$x\to-1, f(x)\to-\infty$而$f(0)=0$,所以当$x\in(t,0), f(x)>0$所以$f(x)$在$(-1,t)$上有唯一零点,$(t,0)$上无零点即$f(x)$在$(-1,0)$上有唯一零点所以$a<-1$,符合题意所以若$f(x)$在区间$(-1,0),(0,+\infty)$各恰有一个零点,求$a$的取值范围为$(-\infty,-1)$
\end{solution}
\end{question}

\begin{question}[points=5]
\textbf{来源：}2022 年 全国乙卷-文科，第 11 题\par
函数 $f(x) = \cos x + (x+1)\sin x + 1$ 在区间 $[0, 2\pi]$ 的最小值、最大值分别为
\begin{choices}
  \item $-\frac{\pi}{2}, \frac{\pi}{2}$
  \item $-\frac{3\pi}{2}, \frac{\pi}{2}$
  \item $-\frac{\pi}{2}, \frac{\pi}{2}+2$
  \item $-\frac{3\pi}{2}, \frac{\pi}{2}+2$
\end{choices}
\begin{solution}
\textbf{答案：}D

\textbf{分析：}利用导数求得单调区间，从而判断最值.

\textbf{详解：}$f'(x) = -\sin x + \sin x + (x+1)\cos x = (x+1)\cos x$, 在 $(0,\frac{\pi}{2})$ 和 $(\frac{3\pi}{2},2\pi)$ 上 $f'(x)>0$, $f(x)$ 单增；在 $(\frac{\pi}{2},\frac{3\pi}{2})$ 上 $f'(x)<0$, $f(x)$ 单减. 又 $f(0)=f(2\pi)=2$, $f(\frac{\pi}{2})=\frac{\pi}{2}+2$, $f(\frac{3\pi}{2}) = -\frac{3\pi}{2}$, 所以最小值为 $-\frac{3\pi}{2}$, 最大值为 $\frac{\pi}{2}+2$.
\end{solution}
\end{question}

\begin{question}[points=12]
\textbf{来源：}2022 年 全国乙卷-文科，第 20 题\par
已知函数 $f(x) = ax - \frac{1}{x} - (a+1)\ln x$.

(1) 当 $a=0$ 时，求 $f(x)$ 的最大值；

(2) 若 $f(x)$ 恰有一个零点，求 $a$ 的取值范围.
\begin{solution}
\textbf{答案：}(1) $-1$; (2) $(0, +\infty)$.

\textbf{分析：}(1) 求导得单调性，即可得最大值.
(2) 求导分类讨论 $a\le0$, $0<a<1$, $a=1$, $a>1$ 时函数的单调性和零点情况.

\textbf{详解：}(1) $a=0$ 时 $f(x)=-\frac1x-\ln x$, $f'(x)=\frac{1-x}{x^2}$, 当 $x\in(0,1)$ 时 $f'(x)>0$, $f(x)$ 单增；当 $x\in(1,+\infty)$ 时 $f'(x)<0$, $f(x)$ 单减，所以 $f(x)_{\max}=f(1)=-1$.
(2) $f'(x)=\frac{(ax-1)(x-1)}{x^2}$. 当 $a\le0$ 时，$ax-1<0$, $f(x)$ 在 $(0,1)$ 单增，在 $(1,+\infty)$ 单减，$f(1)=a-1<0$, 无零点. 当 $0<a<1$ 时，$f(x)$ 在 $(0,1), (\frac1a,+\infty)$ 单增，在 $(1,\frac1a)$ 单减，$f(1)=a-1<0$, 取 $m=(\frac3a+2)^2$, 则 $f(m)>0$, 所以仅有一个零点. 当 $a=1$ 时，$f'(x)\ge0$, $f(x)$ 单增，$f(1)=0$, 唯一零点. 当 $a>1$ 时，$f(x)$ 在 $(0,\frac1a), (1,+\infty)$ 单增，在 $(\frac1a,1)$ 单减，$f(1)=a-1>0$, 取 $n=\frac1{4(a+1)}$, 则 $f(n)<0$, 所以有一个零点. 综上 $a>0$.
\end{solution}
\end{question}

\begin{question}[points=15]
\textbf{来源：}2022 年 天津卷，第 20 题\par
已知 $a,b\in\mathbb{R}$，函数 $f(x)=e^x-a\sin x$，$g(x)=b\sqrt{x}$．

（1）求函数 $y=f(x)$ 在 $(0,f(0))$ 处的切线方程；

（2）若 $y=f(x)$ 和 $y=g(x)$ 有公共点，

（i）当 $a=0$ 时，求 $b$ 的取值范围；

（ii）求证：$a^2+b^2>e$．
\begin{solution}
\textbf{答案：}（1）$y=(1-a)x+1$；（2）（i）$b\in[\sqrt{2e},+\infty)$；（ii）证明见解析

\textbf{分析：}（1）求出 $f'(0)$ 可求切线方程；（2）（i）当 $a=0$ 时，曲线 $y=f(x)$ 和 $y=g(x)$ 有公共点即 $e^x=b\sqrt{x}$ 有解，设 $t=\sqrt{x}$，则 $e^{t^2}=bt$ 在 $[0,+\infty)$ 上有解，设 $s(t)=e^{t^2}-bt$，求导后分类讨论结合零点存在定理可求 $b\in[\sqrt{2e},+\infty)$；（ii）曲线 $y=f(x)$ 和 $y=g(x)$ 有公共点即 $a\sin x_0+b\sqrt{x_0}-e^{x_0}=0$，利用点到直线的距离得到 $\sqrt{a^2+b^2}\ge\frac{e^{x_0}}{\sqrt{\sin^2 x_0+x_0}}$，利用导数证明 $\frac{e^{2x_0}}{\sin^2 x_0+x_0}>e$，从而可得不等式成立.

\textbf{详解：}（1）$f'(x)=e^x-a\cos x$，$f'(0)=1-a$，$f(0)=1$，切线方程为 $y=(1-a)x+1$.
（2）（i）$a=0$ 时，$e^x=b\sqrt{x}$ 有解，令 $t=\sqrt{x}\ge0$，则 $e^{t^2}=bt$，即 $e^{t^2}-bt=0$。设 $s(t)=e^{t^2}-bt$，$s'(t)=2te^{t^2}-b$。若 $b\le0$，$s(t)>0$ 无解；若 $b>0$，令 $u(t)=2te^{t^2}-b$，$u'(t)=2(1+2t^2)e^{t^2}>0$，$u(0)=-b<0$，$u(\sqrt{b})=2\sqrt{b}e^{b}-b>0$，故存在唯一 $t_0>0$ 使 $u(t_0)=0$，$s(t)$ 在 $(0,t_0)$ 递减，$(t_0,+\infty)$ 递增，$s(t)_{\min}=s(t_0)=e^{t_0^2}-bt_0\le0$ 时有解，结合 $b=2t_0e^{t_0^2}$ 得 $e^{t_0^2}-2t_0^2e^{t_0^2}\le0$，即 $t_0^2\ge\frac{1}{2}$，所以 $b=2t_0e^{t_0^2}\ge2\cdot\frac{1}{\sqrt{2}}e^{\frac{1}{2}}=\sqrt{2e}$，故 $b\in[\sqrt{2e},+\infty)$.
（ii）设公共点为 $(x_0,y_0)$，则 $e^{x_0}-a\sin x_0=b\sqrt{x_0}$，即 $a\sin x_0+b\sqrt{x_0}-e^{x_0}=0$，原点到直线 $a\sin x_0+b\sqrt{x_0}-e^{x_0}=0$ 的距离为 $\frac{|e^{x_0}|}{\sqrt{\sin^2 x_0+x_0}}$，所以 $\sqrt{a^2+b^2}\ge\frac{e^{x_0}}{\sqrt{\sin^2 x_0+x_0}}$，即 $a^2+b^2\ge\frac{e^{2x_0}}{\sin^2 x_0+x_0}$。下证 $\frac{e^{2x}}{\sin^2 x+x}>e$ 对 $x>0$ 恒成立。等价于 $e^{2x-1}>\sin^2 x+x$，即 $2x-1>\ln(\sin^2 x+x)$。令 $h(x)=2x-1-\ln(\sin^2 x+x)$，$h'(x)=2-\frac{2\sin x\cos x+1}{\sin^2 x+x}=2-\frac{\sin2x+1}{\sin^2 x+x}$，由于 $\sin2x\le1$，$\sin^2 x+x\ge x>0$，$\frac{\sin2x+1}{\sin^2 x+x}\le\frac{2}{x}$，当 $x\ge2$ 时，$h'(x)>0$；当 $0<x<2$ 时，可证明 $h(x)>0$。或者直接证 $e^{2x-1}>x+\sin^2 x$，利用 $e^x\ge x+1$ 得 $e^{2x-1}\ge2x$，而 $2x\ge x+\sin^2 x$ 因为 $x\ge\sin^2 x$。所以 $\frac{e^{2x_0}}{\sin^2 x_0+x_0}\ge e$，等号不成立，故 $a^2+b^2>e$.
\end{solution}
\end{question}

\begin{question}[points=5]
\textbf{来源：}2022 年 新高考II卷，第 14 题\par
曲线$y=\ln|x|$过坐标原点的两条切线的方程为，.
\fillin[{$y=\dfrac{1}{e}x$，$y=-\dfrac{1}{e}x$}]
\begin{solution}
\textbf{答案：}$y=\dfrac{1}{e}x$，$y=-\dfrac{1}{e}x$

\textbf{分析：}分$x>0$和$x<0$两种情况，当$x>0$时设切点为$(x_0,\ln x_0)$，求出函数的导函数，即可求出切线的斜率，从而表示出切线方程，再根据切线过坐标原点求出$x_0$，即可求出切线方程，当$x<0$时同理可得.

\textbf{详解：}解：因为$y=\ln|x|$，当$x>0$时$y=\ln x$，设切点为$(x_0,\ln x_0)$，由$y'=\dfrac{1}{x}$，所以$y'|_{x=x_0}=\dfrac{1}{x_0}$，所以切线方程为$y-\ln x_0=\dfrac{1}{x_0}(x-x_0)$，又切线过坐标原点，所以$-\ln x_0=\dfrac{1}{x_0}(-x_0)$，解得$x_0=e$，所以切线方程为$y-1=\dfrac{1}{e}(x-e)$，即$y=\dfrac{1}{e}x$；当$x<0$时$y=\ln(-x)$，设切点为$(x_1,\ln(-x_1))$，由$y'=\dfrac{1}{x}$（因为$(-x)'=-1$，复合函数求导），得$y'=\dfrac{1}{x}$（注意此处$x<0$，$\dfrac{1}{x}$也为负），所以$y'|_{x=x_1}=\dfrac{1}{x_1}$，所以切线方程为$y-\ln(-x_1)=\dfrac{1}{x_1}(x-x_1)$，又切线过坐标原点，所以$-\ln(-x_1)=\dfrac{1}{x_1}(-x_1)$，解得$x_1=-e$，所以切线方程为$y-1=\dfrac{1}{-e}(x+e)$，即$y=-\dfrac{1}{e}x$；故答案为$y=\dfrac{1}{e}x$；$y=-\dfrac{1}{e}x$．
\end{solution}
\end{question}

\begin{question}[points=12]
\textbf{来源：}2022 年 新高考II卷，第 22 题\par
已知函数$f(x)=xe^{ax}-e^x$．

(1)当$a=1$时，讨论$f(x)$的单调性；

(2)当$x>0$时，$f(x)<-1$，求$a$的取值范围；

(3)设$n\in\mathbf{N}^*$，证明：$\dfrac{1}{1^2+1}+\dfrac{1}{2^2+2}+\cdots+\dfrac{1}{n^2+n}>\ln(n+1)$．
\begin{solution}
\textbf{答案：}（1）$f(x)$的减区间为$(-\infty,0)$，增区间为$(0,+\infty)$；（2）$a\leq\dfrac{1}{2}$；（3）见解析

\textbf{分析：}（1）求出$f'(x)$，讨论其符号后可得$f(x)$的单调性.（2）设$h(x)=xe^{ax}-e^x+1$，求出$h''(x)$，先讨论$a>\dfrac{1}{2}$时题设中的不等式不成立，再就$0<a\leq\dfrac{1}{2}$结合放缩法讨论$h'(x)$符号，最后就$a\leq 0$结合放缩法讨论$h(x)$的范围后可得参数的取值范围.（3）由（2）可得$2\ln t<t-\dfrac{1}{t}$对任意的$t>1$恒成立，从而可得$\ln(n+1)-\ln n<\dfrac{1}{n^2+n}$对任意的$n\in\mathbf{N}^*$恒成立，结合裂项相消法可证题设中的不等式.

\textbf{详解：}（1）当$a=1$时，$f(x)=(x-1)e^x$，则$f'(x)=xe^x$，当$x<0$时，$f'(x)<0$，当$x>0$时，$f'(x)>0$，故$f(x)$的减区间为$(-\infty,0)$，增区间为$(0,+\infty)$.
（2）设$h(x)=xe^{ax}-e^x+1$，则$h(0)=0$，又$h'(x)=(1+ax)e^{ax}-e^x$，设$g(x)=(1+ax)e^{ax}-e^x$，则$g'(x)=(2a+a^2x)e^{ax}-e^x$，若$a>\dfrac{1}{2}$，则$g'(0)=2a-1>0$，因为$g'(x)$为连续不间断函数，故存在$x_0>0$，使得$\forall x\in(0,x_0)$，总有$g'(x)>0$，故$g(x)$在$(0,x_0)$为增函数，故$g(x)>g(0)=0$，故$h(x)$在$(0,x_0)$为增函数，故$h(x)>h(0)=0$，即$f(x)>-1$，与题设矛盾. 若$0<a\leq\dfrac{1}{2}$，则$h'(x)=(1+ax)e^{ax}-e^x=e^{ax+\ln(1+ax)}-e^x$，下证：对任意$x>0$，总有$\ln(1+x)<x$成立，证明：设$S(x)=\ln(1+x)-x$，故$S'(x)=\dfrac{1}{1+x}-1=\dfrac{-x}{1+x}<0$，故$S(x)$在$(0,+\infty)$上为减函数，故$S(x)<S(0)=0$即$\ln(1+x)<x$成立. 由上述不等式有$e^{ax+\ln(1+ax)}-e^x<e^{ax+ax}-e^x=e^{2ax}-e^x\leq 0$（因为$2a\leq 1$，且$x>0$时$e^{2ax}\leq e^x$），故$h'(x)\leq 0$总成立，即$h(x)$在$(0,+\infty)$上为减函数，所以$h(x)<h(0)=0$，即$f(x)<-1$. 当$a\leq 0$时，有$h'(x)=e^{ax}-e^x+axe^{ax}<1-1+0=0$（因为$e^{ax}\leq 1$，$e^x>1$，$axe^{ax}\leq 0$），所以$h(x)$在$(0,+\infty)$上为减函数，所以$h(x)<h(0)=0$，即$f(x)<-1$. 综上，$a\leq\dfrac{1}{2}$.
（3）取$a=\dfrac{1}{2}$，则$\forall x>0$，总有$xe^{\frac{x}{2}}-e^x+1<0$成立，令$t=e^{\frac{x}{2}}$，则$t>1$，$t^2=e^x$，$x=2\ln t$，故$2t\ln t<t^2-1$即$2\ln t<t-\dfrac{1}{t}$对任意的$t>1$恒成立. 所以对任意的$n\in\mathbf{N}^*$，有$2\ln\dfrac{n+1}{n}<\dfrac{n+1}{n}-\dfrac{n}{n+1}$，整理得到$\ln(n+1)-\ln n<\dfrac{1}{n^2+n}$，故$\dfrac{1}{1^2+1}+\dfrac{1}{2^2+2}+\cdots+\dfrac{1}{n^2+n}>\ln2-\ln1+\ln3-\ln2+\cdots+\ln(n+1)-\ln n=\ln(n+1)$，故不等式成立.
\end{solution}
\end{question}

\begin{question}[points=5]
\textbf{来源：}2022 年 新高考I卷，第 15 题\par
若曲线 $y = (x + a)e^x$ 有两条过坐标原点的切线, 则 $a$ 的取值范围是 .
\fillin[{$(-\infty, -4) \cup (0, +\infty)$}]
\begin{solution}
\textbf{答案：}$(-\infty, -4) \cup (0, +\infty)$

\textbf{分析：}设出切点横坐标 $x_0$, 利用导数的几何意义求得切线方程, 根据切线经过原点得到关于 $x_0$ 的方程, 根据此方程应有两个不同的实数根, 求得 $a$ 的取值范围.

\textbf{详解：}$\because y = (x + a)e^x$, $\therefore y' = (x + 1 + a)e^x$. 设切点为 $(x_0, y_0)$, 则 $y_0 = (x_0 + a)e^{x_0}$, 切线斜率 $k = (x_0 + 1 + a)e^{x_0}$. 切线方程为 $y - (x_0 + a)e^{x_0} = (x_0 + 1 + a)e^{x_0}(x - x_0)$. $\because$ 切线过原点, $\therefore -(x_0 + a)e^{x_0} = (x_0 + 1 + a)e^{x_0}(-x_0)$, 整理得 $x_0^2 + a x_0 - a = 0$. $\because$ 切线有两条, $\therefore \Delta = a^2 + 4a > 0$, 解得 $a < -4$ 或 $a > 0$. $\therefore a$ 的取值范围是 $(-\infty, -4) \cup (0, +\infty)$.
\end{solution}
\end{question}

\begin{question}[points=12]
\textbf{来源：}2022 年 新高考I卷，第 22 题\par
已知函数 $f(x) = e^x - ax$ 和 $g(x) = ax - \ln x$ 有相同的最小值.

(1) 求 $a$;

(2) 证明: 存在直线 $y = b$, 其与两条曲线 $y = f(x)$ 和 $y = g(x)$ 共有三个不同的交点, 并且从左到右的三个交点的横坐标成等差数列.
\begin{solution}
\textbf{答案：}(1) $a = 1$; (2) 见解析

\textbf{分析：}(1) 根据导数可得函数的单调性, 从而可得相应的最小值, 根据最小值相等可求 $a$. 注意分类讨论. (2) 根据 (1) 可得当 $b > 1$ 时, $e^x - x = b$ 的解的个数、 $x - \ln x = b$ 的解的个数均为 2, 构建新函数 $h(x) = e^x + \ln x - 2x$, 利用导数可得该函数只有一个零点且可得 $f(x)$, $g(x)$ 的大小关系, 根据存在直线 $y = b$ 与曲线 $y = f(x)$, $y = g(x)$ 有三个不同的交点可得 $b$ 的取值, 再根据两类方程的根的关系可证明三根成等差数列.

\textbf{详解：}(1) $f(x) = e^x - ax$ 的定义域为 $\mathbb{R}$, $f'(x) = e^x - a$. 若 $a \le 0$, 则 $f'(x) > 0$, 此时 $f(x)$ 无最小值, 故 $a > 0$. 当 $x < \ln a$ 时, $f'(x) < 0$, $f(x)$ 在 $(-\infty, \ln a)$ 上为减函数; 当 $x > \ln a$ 时, $f'(x) > 0$, $f(x)$ 在 $(\ln a, +\infty)$ 上为增函数. 故 $f(x)_{\min} = f(\ln a) = a - a\ln a$. $g(x) = ax - \ln x$ 的定义域为 $(0, +\infty)$, $g'(x) = a - \frac{1}{x} = \frac{ax-1}{x}$. 当 $0 < x < \frac{1}{a}$ 时, $g'(x) < 0$, $g(x)$ 在 $(0, \frac{1}{a})$ 上为减函数; 当 $x > \frac{1}{a}$ 时, $g'(x) > 0$, $g(x)$ 在 $(\frac{1}{a}, +\infty)$ 上为增函数. 故 $g(x)_{\min} = g\left(\frac{1}{a}\right) = 1 - \ln\frac{1}{a} = 1 + \ln a$. 因为 $f(x)$ 和 $g(x)$ 有相同的最小值, 所以 $a - a\ln a = 1 + \ln a$, 整理得 $\frac{a-1}{a+1} = \ln a$, 其中 $a > 0$. 设 $\varphi(a) = \frac{a-1}{a+1} - \ln a$, $a > 0$, 则 $\varphi'(a) = \frac{2}{(a+1)^2} - \frac{1}{a} = \frac{2a - (a+1)^2}{a(a+1)^2} = \frac{-(a^2+1)}{a(a+1)^2} < 0$, 故 $\varphi(a)$ 为 $(0, +\infty)$ 上的减函数, 而 $\varphi(1) = 0$, 故 $\frac{a-1}{a+1} = \ln a$ 的唯一解为 $a = 1$. 综上, $a = 1$.
(2) 由 (1) 可得 $f(x) = e^x - x$, $g(x) = x - \ln x$ 的最小值为 $1$. 当 $b > 1$ 时, 考虑 $e^x - x = b$ 的解的个数、 $x - \ln x = b$ 的解的个数. 设 $S(x) = e^x - x - b$, $S'(x) = e^x - 1$. 当 $x < 0$ 时, $S'(x) < 0$; 当 $x > 0$ 时, $S'(x) > 0$. 故 $S(x)$ 在 $(-\infty, 0)$ 上为减函数, 在 $(0, +\infty)$ 上为增函数. $S(x)_{\min} = S(0) = 1 - b < 0$. $S(-b) = e^{-b} > 0$, $S(b) = e^b - 2b$. 设 $u(b) = e^b - 2b$, $b > 1$, 则 $u'(b) = e^b - 2 > 0$, 故 $u(b)$ 在 $(1, +\infty)$ 上为增函数, $u(b) > u(1) = e - 2 > 0$, 故 $S(b) > 0$. 故 $S(x) = 0$ 有两个不同的零点, 即 $e^x - x = b$ 的解的个数为 2. 设 $T(x) = x - \ln x - b$, $T'(x) = 1 - \frac{1}{x} = \frac{x-1}{x}$. 当 $0 < x < 1$ 时, $T'(x) < 0$; 当 $x > 1$ 时, $T'(x) > 0$. 故 $T(x)$ 在 $(0,1)$ 上为减函数, 在 $(1, +\infty)$ 上为增函数. $T(x)_{\min} = T(1) = 1 - b < 0$. $T(e^{-b}) = e^{-b} + b - b = e^{-b} > 0$, $T(e^b) = e^b - b - b = e^b - 2b > 0$. 故 $T(x) = 0$ 有两个不同的零点, 即 $x - \ln x = b$ 的解的个数为 2. 当 $b = 1$ 时, 由 (1) 讨论可得 $e^x - x = b$, $x - \ln x = b$ 仅有一个零点; 当 $b < 1$ 时, 均无零点. 故若存在直线 $y = b$ 与曲线 $y = f(x)$, $y = g(x)$ 有三个不同的交点, 则 $b > 1$. 设 $h(x) = e^x + \ln x - 2x$, $x > 0$, 则 $h'(x) = e^x + \frac{1}{x} - 2$. 设 $s(x) = e^x - x - 1$, $x > 0$, 则 $s'(x) = e^x - 1 > 0$, 故 $s(x)$ 在 $(0, +\infty)$ 上为增函数, $s(x) > s(0) = 0$, 即 $e^x > x + 1$. 所以 $h'(x) > (x+1) + \frac{1}{x} - 2 = x + \frac{1}{x} - 1 \ge 2 - 1 = 1 > 0$, 故 $h(x)$ 在 $(0, +\infty)$ 上为增函数. $h(1) = e - 2 > 0$, $h(\frac{1}{e^3}) = e^{\frac{1}{e^3}} + \ln\frac{1}{e^3} - \frac{2}{e^3} = e^{\frac{1}{e^3}} - 3 - \frac{2}{e^3} < e^{\frac{1}{e^3}} - 3 < 0$, 故 $h(x)$ 在 $(0, +\infty)$ 上有且只有一个零点 $x_0$, 且 $\frac{1}{e^3} < x_0 < 1$. 当 $0 < x < x_0$ 时, $h(x) < 0$, 即 $e^x - x < x - \ln x$, 即 $f(x) < g(x)$; 当 $x > x_0$ 时, $h(x) > 0$, 即 $f(x) > g(x)$. 因此若存在直线 $y = b$ 与曲线 $y = f(x)$, $y = g(x)$ 有三个不同的交点, 则 $b = f(x_0) = g(x_0) > 1$. 此时 $e^x - x = b$ 有两个不同的零点 $x_1$, $x_0$ ($x_1 < 0 < x_0$); $x - \ln x = b$ 有两个不同的零点 $x_0$, $x_4$ ($0 < x_0 < 1 < x_4$). 所以 $e^{x_1} - x_1 = b$, $e^{x_0} - x_0 = b$, $x_4 - \ln x_4 = b$, $x_0 - \ln x_0 = b$. 由 $x_4 - b = \ln x_4$ 得 $e^{x_4 - b} = x_4$, 即 $e^{x_4 - b} - (x_4 - b) - b = 0$, 故 $x_4 - b$ 为方程 $e^x - x = b$ 的解. 同理, $x_0 - b$ 也为方程 $e^x - x = b$ 的解. 又 $e^{x_1} - x_1 = b$ 可化为 $e^{x_1} = x_1 + b$, 即 $x_1 - \ln(x_1 + b) = 0$, 即 $(x_1 + b) - \ln(x_1 + b) - b = 0$, 故 $x_1 + b$ 为方程 $x - \ln x = b$ 的解. 同理, $x_0 + b$ 也为方程 $x - \ln x = b$ 的解. 所以 $\{x_1, x_0\} = \{x_0 - b, x_4 - b\}$. 由于 $b > 1$, 故 $x_0 < x_4 - b$? 实际上, 因为 $x_0 < 1$, $b > 1$, 所以 $x_0 - b < 0$, 而 $x_1 < 0$, 所以 $x_1 = x_0 - b$, $x_0 = x_4 - b$. 从而 $x_1 + x_4 = 2x_0$. 故从左到右的三个交点的横坐标成等差数列.
\end{solution}
\end{question}

\begin{question}[points=15]
\textbf{来源：}2022 年 浙江卷，第 22 题\par
设函数 $f(x) = \dfrac{e}{2x} + \ln x$ $(x>0)$．

（1）求 $f(x)$ 的单调区间；

（2）已知 $a,b \in \mathbb{R}$，曲线 $y=f(x)$ 上不同的三点 $(x_1, f(x_1)), (x_2, f(x_2)), (x_3, f(x_3))$ 处的切线都经过点 $(a,b)$．证明：

（ⅰ）若 $a>e$，则 $0 < b - f(a) < \dfrac{1}{2}\left(\dfrac{a}{e} - 1\right)$；

（ⅱ）若 $0<a<e, x_1<x_2<x_3$，则 $\dfrac{2\sqrt{e-a}}{e} + \dfrac{1}{6e^2} < \dfrac{1}{x_1} + \dfrac{1}{x_3} < \dfrac{2\sqrt{e-a}}{a} - \dfrac{1}{6e}$．

（注：$e=2.71828\cdots$ 是自然对数的底数）
\begin{solution}
\textbf{答案：}见解析

\textbf{分析：}（1）求出导数，讨论其符号可得单调区间；（2）（ⅰ）由题设构造关于切点横坐标的方程，根据方程有3个不同的解可证明不等式成立；（ⅱ）通过比值代换等证明．

\textbf{详解：}（1）$f'(x) = -\frac{e}{2x^2} + \frac{1}{x} = \frac{2x-e}{2x^2}$，当 $0<x<\frac{e}{2}$ 时 $f'(x)<0$，当 $x>\frac{e}{2}$ 时 $f'(x)>0$，所以 $f(x)$ 的减区间为 $(0,\frac{e}{2})$，增区间为 $(\frac{e}{2},+\infty)$．
（2）（ⅰ）设切点 $(x,f(x))$，切线方程为 $y-f(x)=f'(x)(x-a)$，即 $g(x)=f(x)-b-f'(x)(x-a)=0$，化简得 $g(x)=1-\frac{a+e}{x}+\frac{ae}{2x^2}-\ln x+b$，$g'(x)=\frac{1}{x^3}(x-e)(x-a)$．当 $a>e$ 时，$g(x)$ 在 $(0,e)$、$(a,+\infty)$ 递减，在 $(e,a)$ 递增．由 $g(x)$ 有三个零点得 $g(e)<0$ 且 $g(a)>0$，代入得 $b<\frac{a}{2e}+1$ 且 $b>\frac{e}{2a}+\ln a=f(a)$，从而 $0<b-f(a)<\frac{1}{2}\left(\frac{a}{e}-1\right)$．
（ⅱ）类似（ⅰ）可得 $\frac{a}{2e}+1<b<\frac{e}{2a}+\ln a$．设 $t=\frac{e}{x}$，$m=\frac{a}{e}\in(0,1)$，则 $t_1,t_2,t_3$ 满足 $-\frac{m}{2}t^2+(m+1)t-\ln t-b=0$．由 $t_1t_3=1$ 及韦达定理可证．具体证明过程略．
\end{solution}
\end{question}

\section*{2023 年}

\begin{question}[points=15]
\textbf{来源：}2023 年 北京卷，第 20 题\par
设函数 $f(x) = x - x^3 e^{ax+b}$，曲线 $y=f(x)$ 在点 $(1,f(1))$ 处的切线方程为 $y = -x+1$．

（1）求 $a,b$ 的值；

（2）设函数 $g(x) = f'(x)$，求 $g(x)$ 的单调区间；

（3）求 $f(x)$ 的极值点个数．
\begin{solution}
\textbf{答案：}（1）$a=-1,b=1$ （2）$g(x)$ 的单调递减区间为 $(0,3-\sqrt{3})$ 和 $(3+\sqrt{3},+\infty)$，单调递增区间为 $(-\infty,0)$ 和 $(3-\sqrt{3},3+\sqrt{3})$ （3）3个

\textbf{分析：}（1）先对 $f(x)$ 求导，利用导数的几何意义得到 $f(1)=0$，$f'(1)=-1$，从而得到关于 $a,b$ 的方程组，解之即可；（2）由（1）得 $g(x)$ 的解析式，从而求得 $g'(x)$，利用数轴穿根法求得 $g'(x)<0$ 与 $g'(x)>0$ 的解，由此求得 $g(x)$ 的单调区间；（3）结合（2）中结论，利用零点存在定理，依次分类讨论区间 $(-\infty,0)$，$(0,x_1)$，$(x_1,x_2)$ 与 $(x_2,+\infty)$ 上 $f'(x)$ 的零点的情况，从而利用导数与函数的极值点的关系求得 $f(x)$ 的极值点个数.

\textbf{详解：}（1）因为 $f(x)=x-x^3e^{ax+b}$，$x\in\mathbb{R}$，所以 $f'(x)=1-(3x^2+ax^3)e^{ax+b}$，因为 $f(x)$ 在 $(1,f(1))$ 处的切线方程为 $y=-x+1$，所以 $f(1)=-1+1=0$，$f'(1)=-1$，则 $\begin{cases} 1-1^3\times e^{a+b}=0 \\ 1-(3+a)e^{a+b}=-1 \end{cases}$，解得 $a=-1,b=1$。
（2）由（1）得 $g(x)=f'(x)=1-(3x^2-x^3)e^{-x+1}$，$x\in\mathbb{R}$，则 $g'(x)=-x(x^2-6x+6)e^{-x+1}$，令 $x^2-6x+6=0$，解得 $x=3\pm\sqrt{3}$，设 $x_1=3-\sqrt{3}$，$x_2=3+\sqrt{3}$，则 $0<x_1<x_2$，易知 $e^{-x+1}>0$ 恒成立。所以令 $g'(x)<0$，解得 $0<x<x_1$ 或 $x>x_2$；令 $g'(x)>0$，解得 $x<0$ 或 $x_1<x<x_2$。所以 $g(x)$ 在 $(0,x_1)$，$(x_2,+\infty)$ 上单调递减，在 $(-\infty,0)$，$(x_1,x_2)$ 上单调递增。
（3）由（1）得 $f(x)=x-x^3e^{-x+1}$，$f'(x)=1-(3x^2-x^3)e^{-x+1}$。由（2）知 $f'(x)$ 在 $(0,x_1)$，$(x_2,+\infty)$ 上单调递减，在 $(-\infty,0)$，$(x_1,x_2)$ 上单调递增。当 $x<0$ 时，$f'(-1)=1-4e^2<0$，$f'(0)=1>0$，所以 $f'(x)$ 在 $(-\infty,0)$ 上存在唯一零点，设为 $x_3$，则 $-1<x_3<0$，此时当 $x<x_3$ 时 $f'(x)<0$，$f(x)$ 单调递减；当 $x_3<x<0$ 时 $f'(x)>0$，$f(x)$ 单调递增，所以 $f(x)$ 在 $(-\infty,0)$ 上有一个极小值点。当 $x\in(0,x_1)$ 时，$f'(x)$ 在 $(0,x_1)$ 上单调递减，$f'(x_1)=f'(3-\sqrt{3})<f'(1)=1-2<0$，故 $f'(0)f'(x_1)<0$，所以 $f'(x)$ 在 $(0,x_1)$ 上存在唯一零点，设为 $x_4$，则 $0<x_4<x_1$，此时当 $0<x<x_4$ 时 $f'(x)>0$，$f(x)$ 单调递增；当 $x_4<x<x_1$ 时 $f'(x)<0$，$f(x)$ 单调递减，所以 $f(x)$ 在 $(0,x_1)$ 上有一个极大值点。当 $x\in(x_1,x_2)$ 时，$f'(x)$ 在 $(x_1,x_2)$ 上单调递增，$f'(x_2)=f'(3+\sqrt{3})>f'(3)=1>0$，故 $f'(x_1)f'(x_2)<0$，所以 $f'(x)$ 在 $(x_1,x_2)$ 上存在唯一零点，设为 $x_5$，则 $x_1<x_5<x_2$，此时当 $x_1<x<x_5$ 时 $f'(x)<0$，$f(x)$ 单调递减；当 $x_5<x<x_2$ 时 $f'(x)>0$，$f(x)$ 单调递增，所以 $f(x)$ 在 $(x_1,x_2)$ 上有一个极小值点。当 $x>x_2=3+\sqrt{3}>3$ 时，$3x^2-x^3=x^2(3-x)<0$，所以 $f'(x)=1-(3x^2-x^3)e^{-x+1}>0$，$f(x)$ 单调递增，无极值点。综上，$f(x)$ 在 $(-\infty,0)$ 和 $(x_1,x_2)$ 上各有一个极小值点，在 $(0,x_1)$ 上有一个极大值点，共有3个极值点。
\end{solution}
\end{question}

\begin{question}[points=12]
\textbf{来源：}2023 年 全国甲卷-理科，第 21 题\par
已知函数 $f(x) = ax - \frac{\sin x}{\cos^3 x}$，$x \in \left(0, \frac{\pi}{2}\right)$．

（1）当 $a=8$ 时，讨论 $f(x)$ 的单调性；

（2）若 $f(x) < \sin 2x$ 恒成立，求 $a$ 的取值范围．
\begin{solution}
\textbf{答案：}（1）$f(x)$ 在 $(0,\frac{\pi}{4})$ 上单调递增，在 $(\frac{\pi}{4},\frac{\pi}{2})$ 上单调递减；（2）$(-\infty,3]$

\textbf{分析：}（1）求导，然后令 $t=\cos^2 x$，讨论导数的符号即可；（2）构造 $g(x)=f(x)-\sin 2x$，计算 $g'(x)$ 的最大值，然后与 0 比较大小，得出 $a$ 的分界点，再对 $a$ 讨论即可．

\textbf{详解：}（1）$f'(x) = a - \frac{\cos x \cos^3 x + 3\sin x \cos^2 x \sin x}{\cos^6 x} = a - \frac{\cos^2 x + 3\sin^2 x}{\cos^4 x} = a - \frac{3-2\cos^2 x}{\cos^4 x}$，令 $\cos^2 x = t$，则 $t\in(0,1)$，则 $f'(x) = g(t) = a - \frac{3-2t}{t^2} = \frac{at^2+2t-3}{t^2}$，当 $a=8$ 时，$f'(x) = g(t) = \frac{8t^2+2t-3}{t^2} = \frac{(2t-1)(4t+3)}{t^2}$，当 $t\in(0,\frac{1}{2})$ 时，$f'(x)<0$；当 $t\in(\frac{1}{2},1)$ 时，$f'(x)>0$．所以 $f(x)$ 在 $(0,\frac{\pi}{4})$ 上单调递增，在 $(\frac{\pi}{4},\frac{\pi}{2})$ 上单调递减．
（2）设 $g(x)=f(x)-\sin 2x$，$g'(x)=f'(x)-2\cos 2x = g(t) - 2(2\cos^2 x -1) = \frac{at^2+2t-3}{t^2} - 2(2t-1) = a+2 - 4t + \frac{2}{t} - \frac{3}{t^2}$，设 $\varphi(t)=a+2-4t+\frac{2}{t}-\frac{3}{t^2}$，$\varphi'(t)=-4-\frac{2}{t^2}+\frac{6}{t^3} = \frac{-4t^3-2t+6}{t^3} = -\frac{2(t-1)(2t^2+2t+3)}{t^3}>0$，所以 $\varphi(t)<\varphi(1)=a-3$．①若 $a\in(-\infty,3]$，$g'(x)=\varphi(t)<a-3\le0$，即 $g(x)$ 在 $(0,\frac{\pi}{2})$ 上单调递减，所以 $g(x)<g(0)=0$，所以当 $a\in(-\infty,3]$ 时，$f(x)<\sin 2x$，符合题意．②若 $a\in(3,+\infty)$，当 $t\to0$ 时，$\frac{2}{t}-\frac{3}{t^2} = -3\left(\frac{1}{t}-\frac{1}{3}\right)^2+\frac{1}{3} \to -\infty$，所以 $\varphi(t)\to -\infty$，$\varphi(1)=a-3>0$，所以 $\exists t_0\in(0,1)$，使得 $\varphi(t_0)=0$，即 $\exists x_0\in(0,\frac{\pi}{2})$，使得 $g'(x_0)=0$．当 $t\in(t_0,1)$ 时，$\varphi(t)>0$，即当 $x\in(0,x_0)$ 时，$g'(x)>0$，$g(x)$ 单调递增，所以当 $x\in(0,x_0)$ 时，$g(x)>g(0)=0$，不合题意．综上，$a$ 的取值范围为 $(-\infty,3]$．
\end{solution}
\end{question}

\begin{question}[points=5]
\textbf{来源：}2023 年 全国甲卷-文科，第 8 题\par
曲线 $y = \frac{e^x}{x+1}$ 在点 $\left(1, \frac{e}{2}\right)$ 处的切线方程为（    ）
\begin{choices}
  \item $y = \frac{e}{4}x$
  \item $y = \frac{e}{2}x$
  \item $y = \frac{e}{4}x + \frac{e}{4}$
  \item $y = \frac{e}{2}x + \frac{3e}{4}$
\end{choices}
\begin{solution}
\textbf{答案：}C

\textbf{分析：}先由切点设切线方程，再求函数的导数，把切点的横坐标代入导数得到切线的斜率，代入所设方程即可求解.

\textbf{详解：}设曲线 $y = \frac{e^x}{x+1}$ 在点 $\left(1, \frac{e}{2}\right)$ 处的切线方程为 $y - \frac{e}{2} = k(x-1)$，因为 $y = \frac{e^x}{x+1}$，所以 $y' = \frac{e^x(x+1)-e^x}{(x+1)^2} = \frac{xe^x}{(x+1)^2}$，所以 $k = y'|_{x=1}= \frac{e}{4}$，所以 $y - \frac{e}{2} = \frac{e}{4}(x-1)$，即 $y = \frac{e}{4}x + \frac{e}{4}$，故选C.
\end{solution}
\end{question}

\begin{question}
\textbf{来源：}2023 年 全国甲卷-文科，第 20 题\par
已知函数 $f(x)=ax-\frac{\sin x}{\cos^2 x}, x\in(0,\frac{\pi}{2})$．

（1）当 $a=1$ 时，讨论 $f(x)$ 的单调性；

（2）若 $f(x)+\sin x<0$，求 $a$ 的取值范围．
\begin{solution}
\textbf{答案：}（1）$f(x)$ 在 $(0,\frac{\pi}{2})$ 上单调递减；（2）$a\le0$

\textbf{分析：}（1）代入 $a=1$ 后，再对 $f(x)$ 求导，同时利用三角函数的平方关系化简 $f'(x)$，再利用换元法判断得其分子与分母的正负情况，从而得解；（2）法一：构造函数 $g(x)=f(x)+\sin x$，从而得到 $g(x)<0$，注意到 $g(0)=0$，从而得到 $g'(0)\le0$，进而得到 $a\le0$，再分类讨论 $a=0$ 与 $a<0$ 两种情况即可得解；法二：先化简并判断得 $\sin x-\frac{\sin x}{\cos^2 x}<0$ 恒成立，再分类讨论 $a=0$，$a<0$ 与 $a>0$ 三种情况，利用零点存在定理与隐零点的知识判断得 $a>0$ 时不满足题意，从而得解.

\textbf{详解：}（1）因为 $a=1$，所以 $f(x)=x-\frac{\sin x}{\cos^2 x}, x\in(0,\frac{\pi}{2})$，则 $f'(x)=1-\frac{\cos x\cdot\cos^2 x-\cos x\cdot(-\sin x)\sin x}{\cos^4 x}=1-\frac{\cos^2 x+2\sin^2 x}{\cos^3 x}=\frac{\cos^3 x-\cos^2 x-2(1-\cos^2 x)}{\cos^3 x}=\frac{\cos^3 x+\cos^2 x-2}{\cos^3 x}$，令 $t=\cos x$，由于 $x\in(0,\frac{\pi}{2})$，所以 $t=\cos x\in(0,1)$，所以 $\cos^3 x+\cos^2 x-2=t^3+t^2-2=(t-1)(t^2+2t+2)$，因为 $t^2+2t+2>0$，$t-1<0$，$\cos^3 x=t^3>0$，所以 $f'(x)<0$ 在 $(0,\frac{\pi}{2})$ 上恒成立，所以 $f(x)$ 在 $(0,\frac{\pi}{2})$ 上单调递减.
（2）法一：构建 $g(x)=f(x)+\sin x=ax-\frac{\sin x}{\cos^2 x}+\sin x(0<x<\frac{\pi}{2})$，则 $g'(x)=a-\frac{1+\sin^2 x}{\cos^3 x}+\cos x$，若 $g(x)=f(x)+\sin x<0$，且 $g(0)=f(0)+\sin0=0$，则 $g'(0)=a-1+1=a\le0$，解得 $a\le0$，当 $a=0$ 时，因为 $\sin x-\frac{\sin x}{\cos^2 x}=\sin x\left(1-\frac{1}{\cos^2 x}\right)$，又 $x\in(0,\frac{\pi}{2})$，所以 $0<\sin x<1$，$0<\cos x<1$，则 $\frac{1}{\cos^2 x}>1$，所以 $f(x)+\sin x=\sin x-\frac{\sin x}{\cos^2 x}<0$，满足题意；当 $a<0$ 时，由于 $0<x<\frac{\pi}{2}$，显然 $ax<0$，所以 $f(x)+\sin x=ax-\frac{\sin x}{\cos^2 x}+\sin x<\sin x-\frac{\sin x}{\cos^2 x}<0$，满足题意；综上所述：若 $f(x)+\sin x<0$，等价于 $a\le0$，所以 $a$ 的取值范围为 $(-\infty,0]$．
法二：因为 $\sin x-\frac{\sin x}{\cos^2 x}=\frac{\sin x\cos^2 x-\sin x}{\cos^2 x}=\frac{-\sin^3 x}{\cos^2 x}$，因为 $x\in(0,\frac{\pi}{2})$，所以 $0<\sin x<1$，$0<\cos x<1$，故 $\sin x-\frac{\sin x}{\cos^2 x}<0$ 在 $(0,\frac{\pi}{2})$ 上恒成立，所以当 $a=0$ 时，$f(x)+\sin x=\sin x-\frac{\sin x}{\cos^2 x}<0$，满足题意；当 $a<0$ 时，由于 $0<x<\frac{\pi}{2}$，显然 $ax<0$，所以 $f(x)+\sin x=ax-\frac{\sin x}{\cos^2 x}+\sin x<\sin x-\frac{\sin x}{\cos^2 x}<0$，满足题意；当 $a>0$ 时，因为 $f(x)+\sin x=ax-\frac{\sin x}{\cos^2 x}+\sin x=ax-\frac{\sin^3 x}{\cos^2 x}$，令 $g(x)=ax-\frac{\sin^3 x}{\cos^2 x}(0<x<\frac{\pi}{2})$，则 $g'(x)=a-\frac{3\sin^2 x\cos^2 x+2\sin^4 x}{\cos^3 x}$，注意到 $g'(0)=a-0=a>0$，若 $\forall0<x<\frac{\pi}{2}$，$g'(x)>0$，则 $g(x)$ 在 $(0,\frac{\pi}{2})$ 上单调递增，注意到 $g(0)=0$，所以 $g(x)>g(0)=0$，即 $f(x)+\sin x>0$，不满足题意；若 $\exists0<x_0<\frac{\pi}{2}$，$g'(x_0)<0$，则 $g'(0)g'(x_0)<0$，所以在 $(0,\frac{\pi}{2})$ 上最靠近 $x=0$ 处必存在零点 $x_1\in(0,\frac{\pi}{2})$，使得 $g'(x_1)=0$，此时 $g'(x)$ 在 $(0,x_1)$ 上有 $g'(x)>0$，所以 $g(x)$ 在 $(0,x_1)$ 上单调递增，则在 $(0,x_1)$ 上有 $g(x)>g(0)=0$，即 $f(x)+\sin x>0$，不满足题意；综上：$a\le0$.
\end{solution}
\end{question}

\begin{question}[points=12]
\textbf{来源：}2023 年 全国乙卷-理科，第 21 题\par
已知函数 $f(x)=(\frac{1}{x}+a)\ln(1+x)$.



(1) 当 $a=-1$ 时，求曲线 $y=f(x)$ 在点 $(1,f(1))$ 处的切线方程；



(2) 是否存在 $a$，$b$，使得曲线 $y=f(\frac{1}{x})$ 关于直线 $x=b$ 对称，若存在，求 $a$，$b$ 的值，若不存在，说明理由；



(3) 若 $f(x)$ 在 $(0,+\infty)$ 存在极值，求 $a$ 的取值范围.
\begin{solution}

\textbf{详解：}(1) 当 $a=-1$ 时，$f(x)=(\frac{1}{x}-1)\ln(x+1)$，则 $f'(x)=-\frac{1}{x^2}\ln(x+1)+(\frac{1}{x}-1)\times\frac{1}{x+1}$，$f(1)=0$，$f'(1)=-\ln2$，函数在 $(1,f(1))$ 处的切线方程为 $y-0=-\ln2(x-1)$，即 $(\ln2)x+y-\ln2=0$.

(2) 由函数的解析式可得 $f(\frac{1}{x})=(x+a)\ln(\frac{1}{x}+1)$，函数的定义域满足 $\frac{1}{x}+1=\frac{x+1}{x}>0$，即函数的定义域为 $(-\infty,-1)\cup(0,+\infty)$，定义域关于直线 $x=-\frac{1}{2}$ 对称，由题意可得 $b=-\frac{1}{2}$，由对称性可知 $f(-\frac{1}{2}+m)=f(-\frac{1}{2}-m)$（$m>\frac{1}{2}$），取 $m=\frac{3}{2}$ 可得 $f(1)=f(-2)$，即 $(a+1)\ln2=(a-2)\ln\frac{1}{2}$，则 $a+1=2-a$，解得 $a=\frac{1}{2}$，经检验 $a=\frac{1}{2},b=-\frac{1}{2}$ 满足题意，故存在 $a=\frac{1}{2},b=-\frac{1}{2}$.

(3) 由函数的解析式可得 $f'(x)=(-\frac{1}{x^2})\ln(x+1)+(\frac{1}{x}+a)\cdot\frac{1}{x+1}$，由 $f(x)$ 在区间 $(0,+\infty)$ 存在极值点，则 $f'(x)$ 在区间 $(0,+\infty)$ 上存在变号零点；令 $(-\frac{1}{x^2})\ln(x+1)+(\frac{1}{x}+a)\cdot\frac{1}{x+1}=0$，则 $-(x+1)\ln(x+1)+x+ax^2=0$，令 $g(x)=ax^2+x-(x+1)\ln(x+1)$，$f(x)$ 在区间 $(0,+\infty)$ 存在极值点，等价于 $g(x)$ 在区间 $(0,+\infty)$ 上存在变号零点，$g'(x)=2ax-\ln(x+1)$，$g''(x)=2a-\frac{1}{x+1}$. 当 $a\le0$ 时，$g'(x)<0$，$g(x)$ 在区间 $(0,+\infty)$ 上单调递减，此时 $g(x)<g(0)=0$，$g(x)$ 在区间 $(0,+\infty)$ 上无零点，不合题意；当 $a\ge\frac{1}{2}$，$2a\ge1$ 时，由于 $\frac{1}{x+1}<1$，所以 $g''(x)>0$，$g'(x)$ 在区间 $(0,+\infty)$ 上单调递增，所以 $g'(x)>g'(0)=0$，$g(x)$ 在区间 $(0,+\infty)$ 上单调递增，$g(x)>g(0)=0$，所以 $g(x)$ 在区间 $(0,+\infty)$ 上无零点，不符合题意；当 $0<a<\frac{1}{2}$ 时，由 $g''(x)=2a-\frac{1}{x+1}=0$ 可得 $x=\frac{1}{2a}-1$，当 $x\in(0,\frac{1}{2a}-1)$ 时，$g''(x)<0$，$g'(x)$ 单调递减，当 $x\in(\frac{1}{2a}-1,+\infty)$ 时，$g''(x)>0$，$g'(x)$ 单调递增，故 $g'(x)$ 的最小值为 $g'(\frac{1}{2a}-1)=1-2a+\ln2a$，令 $m(x)=1-x+\ln x$（$0<x<1$），则 $m'(x)=\frac{-x+1}{x}>0$，函数 $m(x)$ 在定义域内单调递增，$m(x)<m(1)=0$，据此可得 $1-x+\ln x<0$ 恒成立，则 $g'(\frac{1}{2a}-1)=1-2a+\ln2a<0$，令 $h(x)=\ln x-x^2+x$（$x>0$），则 $h'(x)=\frac{-2x^2+x+1}{x}$，当 $x\in(0,1)$ 时，$h'(x)>0$，$h(x)$ 单调递增，当 $x\in(1,+\infty)$ 时，$h'(x)<0$，$h(x)$ 单调递减，故 $h(x)\le h(1)=0$，即 $\ln x\le x^2-x$（取等条件为 $x=1$），所以 $g'(x)=2ax-\ln(x+1)>2ax-[(x+1)^2-(x+1)]=2ax-x^2+x$，$g'(2a-1)>2a(2a-1)-[(2a-1)^2+(2a-1)]=0$，且注意到 $g'(0)=0$，根据零点存在性定理可知 $g'(x)$ 在区间 $(0,+\infty)$ 上存在唯一零点 $x_0$. 当 $x\in(0,x_0)$ 时，$g'(x)<0$，$g(x)$ 单调减，当 $x\in(x_0,+\infty)$ 时，$g'(x)>0$，$g(x)$ 单调递增，所以 $g(x_0)<g(0)=0$. 令 $n(x)=\ln x-\frac{1}{2}(x-\frac{1}{x})$，则 $n'(x)=\frac{1}{x}-\frac{1}{2}(1+\frac{1}{x^2})=\frac{-(x-1)^2}{2x^2}\le0$，则 $n(x)$ 单调递减，注意到 $n(1)=0$，故当 $x\in(1,+\infty)$ 时，$\ln x-\frac{1}{2}(x-\frac{1}{x})<0$，从而有 $\ln x<\frac{1}{2}(x-\frac{1}{x})$，所以 $g(x)=ax^2+x-(x+1)\ln(x+1)>ax^2+x-(x+1)\times\frac{1}{2}[(x+1)-\frac{1}{x+1}]=(a-\frac{1}{2})x^2+\frac{1}{2}$，令 $(a-\frac{1}{2})x^2+\frac{1}{2}=0$ 得 $x^2=\frac{1}{1-2a}$，所以 $g(\frac{1}{\sqrt{1-2a}})>0$，所以函数 $g(x)$ 在区间 $(0,+\infty)$ 上存在变号零点，符合题意. 综合上面可知实数 $a$ 的取值范围是 $(0,\frac{1}{2})$.
\end{solution}
\end{question}

\begin{question}[points=12]
\textbf{来源：}2023 年 全国乙卷-文科，第 20 题\par
已知函数$f(x)=\left(\frac{1}{x}+a\right)\ln(1+x)$．

（1）当$a=-1$时，求曲线$y=f(x)$在点$(1,f(1))$处的切线方程．

（2）若函数$f(x)$在$(0,+\infty)$单调递增，求$a$的取值范围．
\begin{solution}

\textbf{分析：}（1）由题意首先求得导函数的解析式，然后由导数的几何意义确定切线的斜率和切点坐标，最后求解切线方程即可；
（2）原问题即$f'(x)\ge0$在区间$(0,+\infty)$上恒成立，整理变形可得$g(x)=ax^2+x-(x+1)\ln(x+1)\ge0$在区间$(0,+\infty)$上恒成立，然后分类讨论$a\le0, a\ge\frac{1}{2}, 0<a<\frac{1}{2}$三种情况即可求得实数$a$的取值范围.

\textbf{详解：}（1）当$a=-1$时，$f(x)=\left(\frac{1}{x}-1\right)\ln(x+1)(x>-1)$，则$f'(x)=-\frac{1}{x^2}\ln(x+1)+\left(\frac{1}{x}-1\right)\cdot\frac{1}{x+1}$，据此可得$f(1)=0, f'(1)=-\ln2$，所以函数在$(1,f(1))$处的切线方程为$y-0=-\ln2(x-1)$，即$(\ln2)x+y-\ln2=0$.
（2）由函数的解析式可得$f'(x)=\left(-\frac{1}{x^2}\right)\ln(x+1)+\left(\frac{1}{x}+a\right)\cdot\frac{1}{x+1}(x>-1)$，满足题意时$f'(x)\ge0$在区间$(0,+\infty)$上恒成立. 令$\left(-\frac{1}{x^2}\right)\ln(x+1)+\left(\frac{1}{x}+a\right)\cdot\frac{1}{x+1}\ge0$，则$-(x+1)\ln(x+1)+(x+ax^2)\ge0$，令$g(x)=ax^2+x-(x+1)\ln(x+1)$，原问题等价于$g(x)\ge0$在区间$(0,+\infty)$上恒成立，则$g'(x)=2ax-\ln(x+1)$，当$a\le0$时，由于$2ax\le0, \ln(x+1)>0$，故$g'(x)<0$，$g(x)$在区间$(0,+\infty)$上单调递减，此时$g(x)<g(0)=0$，不合题意；令$h(x)=g'(x)=2ax-\ln(x+1)$，则$h'(x)=2a-\frac{1}{x+1}$，当$a\ge\frac{1}{2}$，$2a\ge1$时，由于$\frac{1}{x+1}<1$，所以$h'(x)>0, h(x)$在区间$(0,+\infty)$上单调递增，即$g'(x)$在区间$(0,+\infty)$上单调递增，所以$g'(x)>g'(0)=0$，$g(x)$在区间$(0,+\infty)$上单调递增，$g(x)>g(0)=0$，满足题意. 当$0<a<\frac{1}{2}$时，由$h'(x)=2a-\frac{1}{x+1}=0$可得$x=\frac{1}{2a}-1$，当$x\in(0,\frac{1}{2a}-1)$时，$h'(x)<0, h(x)$在区间$(0,\frac{1}{2a}-1)$上单调递减，即$g'(x)$单调递减，注意到$g'(0)=0$，故当$x\in(0,\frac{1}{2a}-1)$时，$g'(x)<g'(0)=0$，$g(x)$单调递减，由于$g(0)=0$，故当$x\in(0,\frac{1}{2a}-1)$时，$g(x)<g(0)=0$，不合题意. 综上可知：实数$a$的取值范围是$\left\{a\mid a\ge\frac{1}{2}\right\}$.
\end{solution}
\end{question}

\begin{question}[points=18]
\textbf{来源：}2023 年 上海卷，第 21 题；次专题：数列\par
令 $f(x)=\ln x$，取点 $P_1(x_1,y_1)$ 过其曲线 $y=f(x)$ 做切线交 $x$ 轴于 $(x_2,0)$，取点 $P_2(x_2,f(x_2))$ 过其做切线交 $x$ 轴于 $(x_3,0)$，若 $|x_{n+1}-x_n|\le 1$ 则停止，以此类推，得到数列 $\{x_n\}$。

(1)若正整数 $k$，证明 $x_{n+1}=x_n-\frac{f(x_n)}{f'(x_n)}$；

(2)若正整数 $n$，试比较 $x_n$ 与 $e$ 大小；

(3)若正整数 $k$，是否存在 $k$ 使得 $x_1,x_2,x_3$ 依次成等差数列？若存在，求出 $k$ 的所有取值，若不存在，试说明理由。
\begin{solution}

\textbf{详解：}(1) $f(x)=\ln x$，则 $f'(x)=\frac{1}{x}$。曲线在 $P_n(x_n,\ln x_n)$ 处的切线方程为 $y-\ln x_n=\frac{1}{x_n}(x-x_n)$，令 $y=0$，得 $0-\ln x_n=\frac{1}{x_n}(x_{n+1}-x_n)$，所以 $-\ln x_n = \frac{x_{n+1}-x_n}{x_n}$，即 $x_{n+1}=x_n-x_n\ln x_n$，而 $\frac{f(x_n)}{f'(x_n)}=\frac{\ln x_n}{1/x_n}=x_n\ln x_n$，故 $x_{n+1}=x_n-\frac{f(x_n)}{f'(x_n)}$。
(2) $x_{n+1}=x_n(1-\ln x_n)$。考虑 $g(x)=x(1-\ln x)$，$g'(x)=1-\ln x-1=-\ln x$。当 $x>1$ 时 $g'(x)<0$，$g(x)$ 递减；$0<x<1$ 时 $g'(x)>0$，$g(x)$ 递增。且 $g(e)=e(1-1)=0$。当 $x_1>e$ 时，$x_2=x_1(1-\ln x_1)<0$，但 $x_n$ 应取正，故需 $x_1\le e$。若 $x_1=e$，则 $x_2=0$，停止。若 $x_1<e$，则 $x_2>0$，且 $x_2<x_1$？具体比较：当 $x_n<e$ 时，$\ln x_n<1$，$1-\ln x_n>0$，$x_{n+1}=x_n(1-\ln x_n)$，由于 $1-\ln x_n<1$，所以 $x_{n+1}<x_n$。且 $x_{n+1}>0$。又 $x_{n+1}=x_n(1-\ln x_n) >? e$？当 $x_n>1$ 时，$1-\ln x_n<1$，$x_{n+1}<x_n$，但可能仍大于1。综合得 $x_n<e$。
(3) 假设存在 $k$ 使 $x_1,x_2,x_3$ 成等差，则 $2x_2=x_1+x_3$。由递推 $x_2=x_1(1-\ln x_1)$，$x_3=x_2(1-\ln x_2)$，代入得 $2x_1(1-\ln x_1)=x_1+x_1(1-\ln x_1)(1-\ln(x_1(1-\ln x_1)))$。化简得方程，然后通过构造函数 $h(x)=...$ 研究零点。原解析：存在 $k=1$。
\end{solution}
\end{question}

\begin{question}[points=16]
\textbf{来源：}2023 年 天津卷，第 20 题；次专题：数列\par
已知函数 $f(x) = \left(\dfrac{1}{x} + \dfrac{1}{2}\right)\ln(x+1)$.

(1) 求曲线 $y = f(x)$ 在 $x = 2$ 处切线的斜率;

(2) 当 $x > 0$ 时, 证明: $f(x) > 1$;

(3) 证明: $\dfrac{5}{6} < \ln(n!) - \left(n+\dfrac{1}{2}\right)\ln n + n \le 1$.
\begin{solution}
\textbf{答案：}(1) $\dfrac{1}{3} - \dfrac{\ln 3}{4}$; (2) 证明见解析; (3) 证明见解析

\textbf{详解：}(1) $f(x) = \dfrac{\ln(x+1)}{x} + \dfrac{\ln(x+1)}{2}$, $f'(x) = \dfrac{1}{x(x+1)} - \dfrac{\ln(x+1)}{x^2} + \dfrac{1}{2(x+1)}$, 所以 $f'(2) = \dfrac{1}{6} - \dfrac{\ln3}{4} + \dfrac{1}{6} = \dfrac{1}{3} - \dfrac{\ln3}{4}$.
(2) 要证 $x>0$ 时 $f(x) > 1$, 即 $\ln(x+1) > \dfrac{2x}{x+2}$. 令 $g(x) = \ln(x+1) - \dfrac{2x}{x+2}$, $g'(x) = \dfrac{1}{x+1} - \dfrac{4}{(x+2)^2} = \dfrac{x^2}{(x+1)(x+2)^2} > 0$, 所以 $g(x)$ 在 $(0,+\infty)$ 上递增, $g(x) > g(0) = 0$, 故 $\ln(x+1) > \dfrac{2x}{x+2}$, 即 $f(x) > 1$.
(3) 设 $h(n) = \ln(n!) - \left(n+\frac{1}{2}\right)\ln n + n$, $n \in \mathbb{N}^*$. 则 $h(n+1)-h(n) = 1 - \left(n+\frac{1}{2}\right)\ln\left(1+\frac{1}{n}\right)$. 由 (2) 知 $\left(n+\frac{1}{2}\right)\ln\left(1+\frac{1}{n}\right) > 1$, 所以 $h(n+1)-h(n) < 0$, $\{h(n)\}$ 递减, 故 $h(n) \le h(1) = 1$. 下证 $h(n) > \dfrac{5}{6}$. 令 $\varphi(x) = \ln x - \dfrac{(x+5)(x-1)}{4x+2}$, $x>0$, 则 $\varphi'(x) = \dfrac{(x-1)^2(1-x)}{x(2x+1)^2}$, 当 $0<x<1$ 时 $\varphi'(x)>0$, 当 $x>1$ 时 $\varphi'(x)<0$, 所以 $\varphi(x) \le \varphi(1)=0$, 即 $\ln x \le \dfrac{(x+5)(x-1)}{4x+2}$ 恒成立. 则 $h(n)-h(n+1) = \left(n+\frac{1}{2}\right)\ln\left(1+\frac{1}{n}\right) - 1 \le \left(n+\frac{1}{2}\right) \cdot \dfrac{\left(\frac{6+\frac{1}{n}}{\frac{2}{n}}\right)}{\frac{4}{n}+2} = \dfrac{1}{12}\left(\frac{1}{n-1} - \frac{1}{n}\right)$. 累加得 $h(2)-h(n) < \frac{1}{12}\left(1-\frac{1}{n}\right)$, 而 $h(2) = 2-\ln2$, 所以 $h(1)-h(n) < \ln2 - 1 + \frac{3}{2} + \frac{1}{12}\left(1-\frac{1}{n}\right) < \ln2 + \frac{1}{2} + \frac{1}{12} = \ln2 + \frac{7}{12} < \frac{5}{6}$, 故 $h(n) > \frac{5}{6}$. 综上, $\frac{5}{6} < h(n) \le 1$.
\end{solution}
\end{question}

\begin{question}[points=5]
\textbf{来源：}2023 年 新高考II卷，第 6 题\par
已知函数$f(x)=ae^{x}-\ln x$在区间$(1,2)$上单调递增，则$a$的最小值为（ \quad ）
\begin{choices}
  \item A. $e^{2}$
  \item B. $e$
  \item C. $e^{-1}$
  \item D. $e^{-2}$
\end{choices}
\begin{solution}
\textbf{答案：}C

\textbf{分析：}根据$f'(x)=ae^{x}-\dfrac{1}{x}\ge0$在$(1,2)$上恒成立，再根据分参求最值即可求出.

\textbf{详解：}依题可知，$f'(x)=ae^{x}-\dfrac{1}{x}\ge0$在$(1,2)$上恒成立，显然$a>0$，所以$xe^{x}\ge\dfrac{1}{a}$，设$g(x)=xe^{x},x\in(1,2)$，则$g'(x)=(x+1)e^{x}>0$，所以$g(x)$在$(1,2)$上单调递增，$g(x)>g(1)=e$，故$e\ge\dfrac{1}{a}$，即$a\ge e^{-1}$，即$a$的最小值为$e^{-1}$.
\end{solution}
\end{question}

\begin{question}[points=5]
\textbf{来源：}2023 年 新高考II卷，第 11 题\par
若函数$f(x)=a\ln x+\dfrac{b}{x}+\dfrac{c}{x^{2}}(a\neq0)$既有极大值也有极小值，则（ \quad ）
\begin{choices}
  \item A. $bc>0$
  \item B. $ab>0$
  \item C. $b^{2}+8ac>0$
  \item D. $ac<0$
\end{choices}
\begin{solution}
\textbf{答案：}BCD

\textbf{分析：}求出函数$f(x)$的导数$f'(x)$，由已知可得$f'(x)$在$(0,+\infty)$上有两个变号零点，转化为一元二次方程有两个不等的正根判断作答.

\textbf{详解：}函数$f(x)=a\ln x+\dfrac{b}{x}+\dfrac{c}{x^{2}}$的定义域为$(0,+\infty)$，求导得$f'(x)=\dfrac{a}{x}-\dfrac{b}{x^{2}}-\dfrac{2c}{x^{3}}=\dfrac{ax^{2}-bx-2c}{x^{3}}$，因为函数$f(x)$既有极大值也有极小值，则函数$f'(x)$在$(0,+\infty)$上有两个变号零点，而$a\neq0$，因此方程$ax^{2}-bx-2c=0$有两个不等的正根$x_{1},x_{2}$，于是$\begin{cases}\Delta=b^{2}+8ac>0\\ x_{1}+x_{2}=\dfrac{b}{a}>0\\ x_{1}x_{2}=-\dfrac{2c}{a}>0\end{cases}$，即有$b^{2}+8ac>0$，$ab>0$，$ac<0$，显然$a^{2}bc<0$，即$bc<0$，A错误，BCD正确.
\end{solution}
\end{question}

\begin{question}[points=10]
\textbf{来源：}2023 年 新高考II卷，第 22 题；次专题：三角函数\par
（1）证明：当$0<x<1$时，$x-x^{2}<\sin x<x$；

（2）已知函数$f(x)=\cos ax-\ln(1-x^{2})$，若$x=0$是$f(x)$的极大值点，求$a$的取值范围．
\begin{solution}
\textbf{答案：}（1）证明见详解；（2）$(-\infty,-\sqrt{2})\cup(\sqrt{2},+\infty)$.

\textbf{分析：}（1）分别构建$F(x)=x-\sin x$，$x\in(0,1)$，$G(x)=x-x^{2}+\sin x$，$x\in(0,1)$，求导，利用导数判断原函数的单调性，进而可得结果；
（2）根据题意结合偶函数的性质可知只需要研究$f(x)$在$(0,1)$上的单调性，求导，分类讨论$0<a^{2}\le 2$和$a^{2}\ge 2$，结合（1）中的结论放缩，根据极大值的定义分析求解.

\textbf{详解：}（1）构建$F(x)=x-\sin x$，$x\in(0,1)$，则$F'(x)=1-\cos x>0$对$\forall x\in(0,1)$恒成立，则$F(x)$在$(0,1)$上单调递增，可得$F(x)>F(0)=0$，所以$x>\sin x$，$x\in(0,1)$；构建$G(x)=\sin x-(x-x^{2})=x-x^{2}+\sin x$，$x\in(0,1)$，则$G'(x)=2x-1+\cos x$，$x\in(0,1)$，构建$g(x)=G'(x)$，$x\in(0,1)$，则$g'(x)=2-\sin x>0$对$\forall x\in(0,1)$恒成立，则$g(x)$在$(0,1)$上单调递增，可得$g(x)>g(0)=0$，即$G'(x)>0$对$\forall x\in(0,1)$恒成立，则$G(x)$在$(0,1)$上单调递增，可得$G(x)>G(0)=0$，所以$\sin x>x-x^{2}$，$x\in(0,1)$；综上所述：$x-x^{2}<\sin x<x$.
（2）令$1-x^{2}>0$，解得$-1<x<1$，即函数$f(x)$的定义域为$(-1,1)$，若$a=0$，则$f(x)=-\ln(1-x^{2})$，$x\in(-1,1)$，因为$y=-\ln u$在定义域内单调递减，$y=1-x^{2}$在$(-1,0)$上单调递增，在$(0,1)$上单调递减，则$f(x)=-\ln(1-x^{2})$在$(-1,0)$上单调递减，在$(0,1)$上单调递增，故$x=0$是$f(x)$的极小值点，不合题意，所以$a\neq0$.
当$a\neq0$时，令$b=|a|>0$，因为$f(x)=\cos ax-\ln(1-x^{2})=\cos(|a|x)-\ln(1-x^{2})=\cos bx-\ln(1-x^{2})$，且$f(-x)=\cos(-bx)-\ln[1-(-x)^{2}]=\cos bx-\ln(1-x^{2})=f(x)$，所以函数$f(x)$在定义域内为偶函数，由题意可得：$f'(x)=-b\sin bx-\dfrac{2x}{x^{2}-1}$，$x\in(-1,1)$.
（i）当$0<b^{2}\le 2$时，取$m=\min\{\dfrac{1}{b},1\}$，$x\in(0,m)$，则$bx\in(0,1)$，由（1）可得$f'(x)=-b\sin(bx)-\dfrac{2x}{x^{2}-1}>-b^{2}x-\dfrac{2x}{x^{2}-1}=\dfrac{x(b^{2}x^{2}+2-b^{2})}{1-x^{2}}$，且$b^{2}x^{2}>0$，$2-b^{2}\ge0$，$1-x^{2}>0$，所以$f'(x)>\dfrac{x(b^{2}x^{2}+2-b^{2})}{1-x^{2}}>0$，即当$x\in(0,m)\subset(0,1)$时，$f'(x)>0$，则$f(x)$在$(0,m)$上单调递增，结合偶函数的对称性可知：$f(x)$在$(-m,0)$上单调递减，所以$x=0$是$f(x)$的极小值点，不合题意；
（ii）当$b^{2}>2$时，取$x\in(0,\dfrac{1}{b})\subset(0,1)$，则$bx\in(0,1)$，由（1）可得$f'(x)=-b\sin(bx)-\dfrac{2x}{x^{2}-1}<-b(bx-b^{2}x^{2})-\dfrac{2x}{x^{2}-1}=\dfrac{x(-b^{3}x^{3}+b^{2}x^{2}+b^{3}x+2-b^{2})}{1-x^{2}}$，构建$h(x)=-b^{3}x^{3}+b^{2}x^{2}+b^{3}x+2-b^{2}$，$x\in(0,\dfrac{1}{b})$，则$h'(x)=-3b^{3}x^{2}+2b^{2}x+b^{3}$，$x\in(0,\dfrac{1}{b})$，且$h'(0)=b^{3}>0$，$h'(\dfrac{1}{b})=b^{3}-b>0$，则$h'(x)>0$对$\forall x\in(0,\dfrac{1}{b})$恒成立，可知$h(x)$在$(0,\dfrac{1}{b})$上单调递增，且$h(0)=2-b^{2}<0$，$h(\dfrac{1}{b})=2>0$，所以$h(x)$在$(0,\dfrac{1}{b})$内存在唯一的零点$n\in(0,\dfrac{1}{b})$，当$x\in(0,n)$时，则$h(x)<0$，且$x>0$，$1-x^{2}>0$，则$f'(x)<\dfrac{x}{1-x^{2}}(-b^{3}x^{3}+b^{2}x^{2}+b^{3}x+2-b^{2})<0$，即当$x\in(0,n)\subset(0,1)$时，$f'(x)<0$，则$f(x)$在$(0,n)$上单调递减，结合偶函数的对称性可知：$f(x)$在$(-n,0)$上单调递增，所以$x=0$是$f(x)$的极大值点，符合题意；
综上所述：$b^{2}>2$，即$a^{2}>2$，解得$a>\sqrt{2}$或$a<-\sqrt{2}$，故$a$的取值范围为$(-\infty,-\sqrt{2})\cup(\sqrt{2},+\infty)$.
\end{solution}
\end{question}

\begin{question}[points=12]
\textbf{来源：}2023 年 新高考I卷，第 19 题\par
已知函数 $f(x)=a(e^x+a)-x$．

（1）讨论 $f(x)$ 的单调性；

（2）证明：当 $a>0$ 时，$f(x)>2\ln a+\frac{3}{2}$．
\begin{solution}
\textbf{答案：}（1）答案见解析；（2）证明见解析

\textbf{分析：}（1）先求导，再分类讨论 $a\leq 0$ 与 $a>0$ 两种情况，结合导数与函数单调性的关系即可得解；（2）方法一：结合（1）中结论，将问题转化为 $a^2-\frac{1}{2}-\ln a>0$ 的恒成立问题，构造函数 $g(a)=a^2-\frac{1}{2}-\ln a(a>0)$，利用导数证得 $g(a)>0$ 即可。方法二：构造函数 $h(x)=e^x-x-1$，证得 $e^x\geq x+1$，从而得到 $f(x)\geq x+\ln a+1+a^2-x$，进而将问题转化为 $a^2-\frac{1}{2}-\ln a>0$ 的恒成立问题，由此得证。

\textbf{详解：}（1）因为 $f(x)=a(e^x+a)-x$，定义域为 $\mathbb{R}$，所以 $f'(x)=ae^x-1$。当 $a\leq 0$ 时，由于 $e^x>0$，则 $ae^x\leq 0$，故 $f'(x)=ae^x-1<0$ 恒成立，所以 $f(x)$ 在 $\mathbb{R}$ 上单调递减；当 $a>0$ 时，令 $f'(x)=ae^x-1=0$，解得 $x=-\ln a$，当 $x<-\ln a$ 时，$f'(x)<0$，则 $f(x)$ 在 $(-\infty,-\ln a)$ 上单调递减；当 $x>-\ln a$ 时，$f'(x)>0$，则 $f(x)$ 在 $(-\ln a,+\infty)$ 上单调递增。综上：当 $a\leq 0$ 时，$f(x)$ 在 $\mathbb{R}$ 上单调递减；当 $a>0$ 时，$f(x)$ 在 $(-\infty,-\ln a)$ 上单调递减，$f(x)$ 在 $(-\ln a,+\infty)$ 上单调递增。（2）方法一：由（1）得，$f(x)_{\min}=f(-\ln a)=a(e^{-\ln a}+a)+\ln a=1+a^2+\ln a$，要证 $f(x)>2\ln a+\frac{3}{2}$，即证 $1+a^2+\ln a>2\ln a+\frac{3}{2}$，即证 $a^2-\frac{1}{2}-\ln a>0$ 恒成立，令 $g(a)=a^2-\frac{1}{2}-\ln a(a>0)$，则 $g'(a)=2a-\frac{1}{a}=\frac{2a^2-1}{a}$，令 $g'(a)<0$，则 $0<a<\frac{\sqrt{2}}{2}$；令 $g'(a)>0$，则 $a>\frac{\sqrt{2}}{2}$；所以 $g(a)$ 在 $(0,\frac{\sqrt{2}}{2})$ 上单调递减，在 $(\frac{\sqrt{2}}{2},+\infty)$ 上单调递增，所以 $g(a)_{\min}=g(\frac{\sqrt{2}}{2})=(\frac{\sqrt{2}}{2})^2-\frac{1}{2}-\ln\frac{\sqrt{2}}{2}=\ln2>0$，则 $g(a)>0$ 恒成立，所以当 $a>0$ 时，$f(x)>2\ln a+\frac{3}{2}$ 恒成立。方法二：令 $h(x)=e^x-x-1$，则 $h'(x)=e^x-1$，由于 $y=e^x$ 在 $\mathbb{R}$ 上单调递增，所以 $h'(x)=e^x-1$ 在 $\mathbb{R}$ 上单调递增，又 $h'(0)=e^0-1=0$，所以当 $x<0$ 时，$h'(x)<0$；当 $x>0$ 时，$h'(x)>0$；所以 $h(x)$ 在 $(-\infty,0)$ 上单调递减，在 $(0,+\infty)$ 上单调递增，故 $h(x)\geq h(0)=0$，则 $e^x\geq x+1$，当且仅当 $x=0$ 时，等号成立，因为 $f(x)=a(e^x+a)-x=ae^x+a^2-x=e^{x+\ln a}+a^2-x\geq x+\ln a+1+a^2-x$，当且仅当 $x+\ln a=0$，即 $x=-\ln a$ 时，等号成立，所以要证 $f(x)>2\ln a+\frac{3}{2}$，即证 $x+\ln a+1+a^2-x>2\ln a+\frac{3}{2}$，即证 $a^2-\frac{1}{2}-\ln a>0$，令 $g(a)=a^2-\frac{1}{2}-\ln a(a>0)$，则 $g'(a)=2a-\frac{1}{a}=\frac{2a^2-1}{a}$，令 $g'(a)<0$，则 $0<a<\frac{\sqrt{2}}{2}$；令 $g'(a)>0$，则 $a>\frac{\sqrt{2}}{2}$；所以 $g(a)$ 在 $(0,\frac{\sqrt{2}}{2})$ 上单调递减，在 $(\frac{\sqrt{2}}{2},+\infty)$ 上单调递增，所以 $g(a)_{\min}=g(\frac{\sqrt{2}}{2})=(\frac{\sqrt{2}}{2})^2-\frac{1}{2}-\ln\frac{\sqrt{2}}{2}=\ln2>0$，则 $g(a)>0$ 恒成立，所以当 $a>0$ 时，$f(x)>2\ln a+\frac{3}{2}$ 恒成立。
\end{solution}
\end{question}

\section*{2024 年}

\begin{question}[points=15]
\textbf{来源：}2024 年 北京卷，第 20 题\par
已知 $f(x)=x+k\ln(1+x)$ 在 $(t,f(t))(t>0)$ 处切线为 $l$。

（1）若切线 $l$ 的斜率 $k=-1$，求 $f(x)$ 单调区间；

（2）证明：切线 $l$ 不经过 $(0,0)$；

（3）已知 $k=1$，$A(t,f(t))$，$C(0,f(t))$，$O(0,0)$，其中 $t>0$，切线 $l$ 与 $y$ 轴交于点 $B$ 时。当 $2S_{\triangle ACO}=15S_{\triangle ABO}$，符合条件的 $A$ 的个数为？

（参考数据：$1.09<\ln3<1.10$，$1.60<\ln5<1.61$，$1.94<\ln7<1.95$）
\begin{solution}
\textbf{答案：}（1）单调递减区间为 $(-1,0)$，单调递增区间为 $(0,+\infty)$；（2）证明见解析；（3）2

\textbf{分析：}（1）直接代入 $k=-1$，再利用导数研究其单调性即可；（2）写出切线方程 $y-f(t)=\left(1+\dfrac{k}{1+t}\right)(x-t)(t>0)$，将 $(0,0)$ 代入再设新函数 $F(t)=\ln(1+t)-\dfrac{t}{1+t}$，利用导数研究其零点即可；（3）分别写出面积表达式，代入 $2S_{\triangle ACO}=15S_{\triangle ABO}$ 得到 $13\ln(1+t)-2t-15\dfrac{t}{1+t}=0$，再设新函数 $h(t)=13\ln(1+t)-2t-\dfrac{15t}{1+t}(t>0)$ 研究其零点即可。

\textbf{详解：}（1）当 $k=-1$ 时，$f(x)=x-\ln(1+x)$，$f'(x)=1-\dfrac{1}{1+x}=\dfrac{x}{1+x}(x>-1)$。当 $x\in(-1,0)$ 时，$f'(x)<0$；当 $x\in(0,+\infty)$ 时，$f'(x)>0$。所以 $f(x)$ 的单调递减区间为 $(-1,0)$，单调递增区间为 $(0,+\infty)$。
（2）$f'(x)=1+\dfrac{k}{1+x}$，切线 $l$ 的斜率为 $1+\dfrac{k}{1+t}$，则切线方程为 $y-f(t)=\left(1+\dfrac{k}{1+t}\right)(x-t)(t>0)$。将 $(0,0)$ 代入则 $-f(t)=-t\left(1+\dfrac{k}{1+t}\right)$，$f(t)=t\left(1+\dfrac{k}{1+t}\right)$，即 $t+k\ln(1+t)=t+\dfrac{kt}{1+t}$，则 $\ln(1+t)=\dfrac{t}{1+t}$，$\ln(1+t)-\dfrac{t}{1+t}=0$。令 $F(t)=\ln(1+t)-\dfrac{t}{1+t}$，假设 $l$ 过 $(0,0)$，则 $F(t)$ 在 $t>0$ 存在零点。$F'(t)=\dfrac{1}{1+t}-\dfrac{(1+t)-t}{(1+t)^2}=\dfrac{t}{(1+t)^2}>0$，所以 $F(t)$ 在 $(0,+\infty)$ 上单调递增，$F(t)>F(0)=0$，$F(t)$ 在 $(0,+\infty)$ 无零点，与假设矛盾，故直线 $l$ 不过 $(0,0)$。
（3）$k=1$ 时，$f(x)=x+\ln(1+x)$，$f'(x)=1+\dfrac{1}{1+x}=\dfrac{x+2}{1+x}>0$。$S_{\triangle ACO}=\dfrac{1}{2}tf(t)$，设 $l$ 与 $y$ 轴交点 $B$ 为 $(0,q)$，$t>0$ 时，若 $q<0$，则此时 $l$ 与 $f(x)$ 必有交点，与切线定义矛盾。由（2）知 $q\neq0$，所以 $q>0$。则切线 $l$ 的方程为 $y-t-\ln(t+1)=\left(1+\dfrac{1}{1+t}\right)(x-t)$，令 $x=0$，则 $y=q=\ln(1+t)-\dfrac{t}{t+1}$。因为 $2S_{\triangle ACO}=15S_{\triangle ABO}$，则 $2tf(t)=15t\left[\ln(1+t)-\dfrac{t}{t+1}\right]$，即 $13\ln(1+t)-2t-15\dfrac{t}{1+t}=0$。记 $h(t)=13\ln(1+t)-2t-\dfrac{15t}{1+t}(t>0)$，则满足条件的 $A$ 有几个即 $h(t)$ 有几个零点。$h'(t)=\dfrac{13}{1+t}-2-\dfrac{15}{(1+t)^2}=\dfrac{13(1+t)-2(1+t)^2-15}{(1+t)^2}=\dfrac{-2t^2+9t-4}{(1+t)^2}=\dfrac{(-2t+1)(t-4)}{(1+t)^2}$。当 $t\in(0,\frac{1}{2})$ 时，$h'(t)<0$，$h(t)$ 单调递减；当 $t\in(\frac{1}{2},4)$ 时，$h'(t)>0$，$h(t)$ 单调递增；当 $t\in(4,+\infty)$ 时，$h'(t)<0$，$h(t)$ 单调递减。因为 $h(0)=0$，$h(\frac{1}{2})<0$，$h(4)=13\ln5-20>13\times1.6-20=0.8>0$，$h(24)=13\ln25-48-\dfrac{15\times24}{25}=26\ln5-48-\dfrac{72}{5}<26\times1.61-48-14.4=41.86-62.4=-20.54<0$，所以由零点存在性定理及 $h(t)$ 的单调性，$h(t)$ 在 $(\frac{1}{2},4)$ 上必有一个零点，在 $(4,24)$ 上必有一个零点。综上所述，$h(t)$ 有两个零点，即满足 $2S_{\triangle ACO}=15S_{\triangle ABO}$ 的 $A$ 有两个。
\end{solution}
\end{question}

\begin{question}[points=5]
\textbf{来源：}2024 年 全国甲卷-理科，第 6 题\par
设函数 $f(x)=\dfrac{e^x+2\sin x}{1+x^2}$，则曲线 $y=f(x)$ 在 $(0,1)$ 处的切线与两坐标轴围成的三角形的面积为（   ）
\begin{choices}
  \item $\dfrac{1}{6}$
  \item $\dfrac{1}{3}$
  \item $\dfrac{1}{2}$
  \item $\dfrac{2}{3}$
\end{choices}
\begin{solution}
\textbf{答案：}A

\textbf{分析：}借助导数的几何意义计算可得其在点 $(0,1)$ 处的切线方程，即可得其与坐标轴交点坐标，即可得其面积。

\textbf{详解：}$f'(x)=\frac{(e^x+2\cos x)(1+x^2)-(e^x+2\sin x)\cdot2x}{(1+x^2)^2}$，则 $f'(0)=\frac{(e^0+2\cos0)(1+0)-(e^0+2\sin0)\times0}{(1+0)^2}=3$，即该切线方程为 $y-1=3x$，即 $y=3x+1$，令 $x=0$，则 $y=1$，令 $y=0$，则 $x=-\frac{1}{3}$，故该切线与两坐标轴所围成的三角形面积 $S=\frac{1}{2}\times1\times\left|-\frac{1}{3}\right|=\frac{1}{6}$。
\end{solution}
\end{question}

\begin{question}[points=12]
\textbf{来源：}2024 年 全国甲卷-理科，第 21 题\par
已知函数 $f(x)=(1-ax)\ln(1+x)-x$。

（1）当 $a=-2$ 时，求 $f(x)$ 的极值；

（2）当 $x\ge0$ 时，$f(x)\ge0$ 恒成立，求 $a$ 的取值范围。
\begin{solution}
\textbf{答案：}（1）极小值为0，无极大值；（2）$a\le -\dfrac{1}{2}$

\textbf{分析：}（1）求出函数的导数，根据导数的单调性和零点可求函数的极值；（2）求出函数的二阶导数，就 $a\le -\frac{1}{2}$、$-\frac{1}{2}<a<0$、$a\ge0$ 分类讨论后可得参数的取值范围。

\textbf{详解：}（1）当 $a=-2$ 时，$f(x)=(1+2x)\ln(1+x)-x$，故 $f'(x)=2\ln(1+x)+\frac{1+2x}{1+x}-1=2\ln(1+x)-\frac{1}{1+x}+1$，因为 $y=2\ln(1+x),y=-\frac{1}{1+x}+1$ 在 $(-1,+\infty)$ 上为增函数，故 $f'(x)$ 在 $(-1,+\infty)$ 上为增函数，而 $f'(0)=0$，故当 $-1<x<0$ 时，$f'(x)<0$，当 $x>0$ 时，$f'(x)>0$，故 $f(x)$ 在 $x=0$ 处取极小值且极小值为 $f(0)=0$，无极大值。
（2）$f'(x)=-a\ln(1+x)+\frac{1-ax}{1+x}-1=-a\ln(1+x)-\frac{(a+1)x}{1+x},x>0$，设 $s(x)=-a\ln(1+x)-\frac{(a+1)x}{1+x},x>0$，则 $s'(x)=-\frac{a}{x+1}-\frac{a+1}{(1+x)^2}=-\frac{ax+2a+1}{(1+x)^2}$，当 $a\le -\frac{1}{2}$ 时，$s'(x)>0$，故 $s(x)$ 在 $(0,+\infty)$ 上为增函数，故 $s(x)>s(0)=0$，即 $f'(x)>0$，所以 $f(x)$ 在 $[0,+\infty)$ 上为增函数，故 $f(x)\ge f(0)=0$。当 $-\frac{1}{2}<a<0$ 时，当 $0<x<-\frac{2a+1}{a}$ 时，$s'(x)<0$，故 $s(x)$ 在 $(0,-\frac{2a+1}{a})$ 上为减函数，故在 $(0,-\frac{2a+1}{a})$ 上 $s(x)<s(0)$，即在 $(0,-\frac{2a+1}{a})$ 上 $f'(x)<0$ 即 $f(x)$ 为减函数，故在 $(0,-\frac{2a+1}{a})$ 上 $f(x)<f(0)=0$，不合题意，舍。当 $a\ge0$，此时 $s'(x)<0$ 在 $(0,+\infty)$ 上恒成立，同理可得在 $(0,+\infty)$ 上 $f(x)<f(0)=0$ 恒成立，不合题意，舍；综上，$a\le -\frac{1}{2}$。
\end{solution}
\end{question}

\begin{question}[points=5]
\textbf{来源：}2024 年 全国甲卷-文科，第 7 题\par
曲线 $f(x) = x^6 + 3x - 1$ 在 $(0,-1)$ 处的切线与坐标轴围成的面积为（ ）
\begin{choices}
  \item $\frac{1}{6}$
  \item $\frac{3}{2}$
  \item $\frac{1}{2}$
  \item $-\frac{3}{2}$
\end{choices}
\begin{solution}
\textbf{答案：}A

\textbf{分析：}先求出切线方程，再求出切线的截距，从而可求面积.

\textbf{详解：}$f'(x) = 6x^5 + 3$，所以 $f'(0) = 3$，故切线方程为 $y = 3(x-0) - 1 = 3x - 1$，故切线的横截距为 $\frac{1}{3}$，纵截距为 $-1$，故切线与坐标轴围成的面积为 $\frac{1}{2} \times 1 \times \frac{1}{3} = \frac{1}{6}$.
\end{solution}
\end{question}

\begin{question}[points=12]
\textbf{来源：}2024 年 全国甲卷-文科，第 17 题\par
已知函数 $f(x) = a(x-1) - \ln x + 1$．

（1）求 $f(x)$ 的单调区间；

（2）若 $a \le 2$ 时，证明：当 $x > 1$ 时，$f(x) < e^{x-1}$ 恒成立．
\begin{solution}
\textbf{答案：}（1）见解析；（2）见解析

\textbf{分析：}（1）求导，含参分类讨论得出导函数的符号，从而得出原函数的单调性；（2）先根据题设条件将问题转化成证明当 $x > 1$ 时，$e^{x-1} - 2x + 1 + \ln x > 0$ 即可.

\textbf{详解：}（1）$f(x)$ 定义域为 $(0, +\infty)$，$f'(x) = a - \frac{1}{x} = \frac{ax - 1}{x}$．当 $a \le 0$ 时，$f'(x) = \frac{ax - 1}{x} < 0$，故 $f(x)$ 在 $(0, +\infty)$ 上单调递减；当 $a > 0$ 时，$x \in \left(\frac{1}{a}, +\infty\right)$ 时，$f'(x) > 0$，$f(x)$ 单调递增，当 $x \in \left(0, \frac{1}{a}\right)$ 时，$f'(x) < 0$，$f(x)$ 单调递减．综上所述，当 $a \le 0$ 时，$f(x)$ 在 $(0, +\infty)$ 上单调递减；$a > 0$ 时，$f(x)$ 在 $\left(\frac{1}{a}, +\infty\right)$ 上单调递增，在 $\left(0, \frac{1}{a}\right)$ 上单调递减．
（2）$a \le 2$，且 $x > 1$ 时，$e^{x-1} - f(x) = e^{x-1} - a(x-1) + \ln x - 1 \ge e^{x-1} - 2x + 1 + \ln x$，令 $g(x) = e^{x-1} - 2x + 1 + \ln x \ (x > 1)$，下证 $g(x) > 0$ 即可．$g'(x) = e^{x-1} - 2 + \frac{1}{x}$，再令 $h(x) = g'(x)$，则 $h'(x) = e^{x-1} - \frac{1}{x^2}$，显然 $h'(x)$ 在 $(1, +\infty)$ 上递增，则 $h'(x) > h'(1) = e^0 - 1 = 0$，即 $g'(x) = h(x)$ 在 $(1, +\infty)$ 上递增，故 $g'(x) > g'(1) = e^0 - 2 + 1 = 0$，即 $g(x)$ 在 $(1, +\infty)$ 上单调递增，故 $g(x) > g(1) = e^0 - 2 + 1 + \ln 1 = 0$，问题得证.
\end{solution}
\end{question}

\begin{question}[points=18]
\textbf{来源：}2024 年 上海卷，第 21 题\par
对于一个函数$f(x)$和一个点$M(a,b)$，令$s(x)=(x-a)^2+(f(x)-b)^2$，若$P(x_0,f(x_0))$是$s(x)$取到最小值的点，则称$P$是$M$在$f(x)$的“最近点”．（1）对于$f(x)=\frac1x$（$x>0$），求证：对于点$M(0,0)$，存在点$P$，使得点$P$是$M$在$f(x)$的“最近点”；（2）对于$f(x)=e^x$，$M(1,0)$，请判断是否存在一个点$P$，它是$M$在$f(x)$的“最近点”，且直线$MP$与$y=f(x)$在点$P$处的切线垂直；（3）已知$y=f(x)$在定义域$\mathbb{R}$上存在导函数$f'(x)$，且函数$g(x)$在定义域$\mathbb{R}$上恒正，设点$M_1(t-1,f(t)-g(t))$，$M_2(t+1,f(t)+g(t))$．若对任意的$t\in\mathbb{R}$，存在点$P$同时是$M_1$，$M_2$在$f(x)$的“最近点”，试判断$f(x)$的单调性．
\begin{solution}
\textbf{答案：}（1）证明见解析；（2）存在，$P(0,1)$；（3）严格单调递减

\textbf{分析：}（1）代入M(0,0)，利用基本不等式即可；（2）由题得$s(x)=(x-1)^2+e^{2x}$，利用导函数得到其最小值，则得到P，再证明直线MP与切线垂直即可；（3）根据题意得到$s_1'(x_0)=s_2'(x_0)=0$，对两等式化简得$f'(x_0)=-\frac1{g(t)}$，再利用“最近点”的定义得到不等式组，即可证明$x_0=t$，最后得到函数单调性.

\textbf{详解：}（1）当M(0,0)时，$s(x)=(x-0)^2+(\frac1x-0)^2=x^2+\frac1{x^2}\ge2\sqrt{x^2\cdot\frac1{x^2}}=2$，当且仅当$x^2=\frac1{x^2}$即$x=1$时取等号，故对于点M(0,0)，存在点P(1,1)，使得该点是M(0,0)在$f(x)$的“最近点”．（2）由题设可得$s(x)=(x-1)^2+(e^x-0)^2=(x-1)^2+e^{2x}$，则$s'(x)=2(x-1)+2e^{2x}$，因为$y=2(x-1)$，$y=2e^{2x}$均为R上单调递增函数，则$s'(x)=2(x-1)+2e^{2x}$在R上为严格增函数，而$s'(0)=0$，故当$x<0$时，$s'(x)<0$，当$x>0$时，$s'(x)>0$，故$s(x)_{\min}=s(0)=2$，此时P(0,1)，而$f'(x)=e^x$，$k=f'(0)=1$，故$f(x)$在点P处的切线方程为$y=x+1$．而$k_{MP}=\frac{0-1}{1-0}=-1$，故$k_{MP}\cdot k=-1$，故直线MP与$y=f(x)$在点P处的切线垂直．（3）设$s_1(x)=(x-t+1)^2+(f(x)-f(t)+g(t))^2$，$s_2(x)=(x-t-1)^2+(f(x)-f(t)-g(t))^2$，而$s_1'(x)=2(x-t+1)+2(f(x)-f(t)+g(t))f'(x)$，$s_2'(x)=2(x-t-1)+2(f(x)-f(t)-g(t))f'(x)$，若对任意的$t\in\mathbb{R}$，存在点P同时是$M_1$，$M_2$在$f(x)$的“最近点”，设$P(x_0,y_0)$，则$x_0$既是$s_1(x)$的最小值点，也是$s_2(x)$的最小值点，因为两函数的定义域均为R，则$x_0$也是两函数的极小值点，则存在$x_0$，使得$s_1'(x_0)=s_2'(x_0)=0$，即$s_1'(x_0)=2(x_0-t+1)+2f'(x_0)[f(x_0)-f(t)+g(t)]=0$①，$s_2'(x_0)=2(x_0-t-1)+2f'(x_0)[f(x_0)-f(t)-g(t)]=0$②，由①②相等得$4+4g(t)f'(x_0)=0$，即$1+f'(x_0)g(t)=0$，即$f'(x_0)=-\frac1{g(t)}$，又因为函数$g(x)$在定义域R上恒正，则$f'(x_0)=-\frac1{g(t)}<0$恒成立，接下来证明$x_0=t$，因为$x_0$既是$s_1(x)$的最小值点，也是$s_2(x)$的最小值点，则$s_1(x_0)\le s_1(t)$，$s_2(x_0)\le s_2(t)$，即$(x_0-t+1)^2+(f(x_0)-f(t)+g(t))^2\le1+(g(t))^2$③，$(x_0-t-1)^2+(f(x_0)-f(t)-g(t))^2\le1+(g(t))^2$④，③+④得$2(x_0-t)^2+2+2[f(x_0)-f(t)]^2+2g^2(t)\le2+2g^2(t)$，即$(x_0-t)^2+[f(x_0)-f(t)]^2\le0$，因为$(x_0-t)^2\ge0$，$[f(x_0)-f(t)]^2\ge0$，则$\begin{cases} x_0-t=0 \\ f(x_0)-f(t)=0 \end{cases}$，解得$x_0=t$，则$f'(t)=-\frac1{g(t)}<0$恒成立，因为t的任意性，则$f(x)$严格单调递减.
\end{solution}
\end{question}

\begin{question}[points=16]
\textbf{来源：}2024 年 天津卷，第 20 题\par
设函数 $f(x)=x\ln x$.

(1) 求 $f(x)$ 图象上点 $(1,f(1))$ 处的切线方程;

(2) 若 $f(x)\geq a(x-\sqrt{x})$ 在 $x\in(0,+\infty)$ 时恒成立, 求 $a$ 的取值范围;

(3) 若 $x_1,x_2\in(0,1)$, 证明 $|f(x_1)-f(x_2)|\leq \sqrt{|x_1-x_2|}$.
\begin{solution}
\textbf{答案：}(1) $y=x-1$; (2) $\{2\}$; (3) 证明见解析

\textbf{分析：}(1) 直接使用导数的几何意义. (2) 先由题设条件得到 $a=2$, 再证明 $a=2$ 时条件满足. (3) 先确定 $f(x)$ 的单调性, 再对 $x_1,x_2$ 分类讨论.

\textbf{详解：}(1) $f'(x)=\ln x+1$, $f(1)=0$, $f'(1)=1$, 切线方程为 $y=x-1$.
(2) 设 $g(t)=a(t-1)-2\ln t$, 则 $f(x)-a(x-\sqrt{x})=x\left[a\left(\frac{1}{\sqrt{x}}-1\right)-2\ln\frac{1}{\sqrt{x}}\right]=x\cdot g\left(\frac{1}{\sqrt{x}}\right)$. 故恒成立等价于 $g(t)\geq 0$ 对 $t>0$ 恒成立. 取 $t=2$ 得 $a\geq 1$, 取 $t=\frac{2}{a}$ 得 $a=2$. 反之, 当 $a=2$ 时, $g(t)=2(t-1)-2\ln t=2(t-1-\ln t)\geq 0$ (因为 $t-1\geq\ln t$). 故 $a$ 的取值范围是 $\{2\}$.
(3) 先证引理: 对 $0<a<b$, $\ln a+1<\frac{f(b)-f(a)}{b-a}<\ln b+1$. 由 $f'(x)=\ln x+1$, $f(x)$ 在 $(0,\frac{1}{e})$ 递减, 在 $(\frac{1}{e},+\infty)$ 递增. 不妨设 $x_1\leq x_2$. 分三种情况: (i) $\frac{1}{e}\leq x_1\leq x_2<1$: $|f(x_1)-f(x_2)|=f(x_2)-f(x_1)<(\ln x_2+1)(x_2-x_1)<x_2-x_1\leq \sqrt{x_2-x_1}$; (ii) $0<x_1\leq x_2\leq \frac{1}{e}$: 利用构造的 $\varphi(x)=x\ln x - c\ln c - \sqrt{c-x}$, 可证 $x_1\ln x_1 - x_2\ln x_2 \leq \sqrt{x_2-x_1}$; (iii) $0<x_1\leq \frac{1}{e}\leq x_2<1$: 利用中间值 $\frac{1}{e}$ 和 (i)(ii) 的结论可得 $|f(x_1)-f(x_2)|\leq \max\{|f(x_1)-f(\frac{1}{e})|, |f(\frac{1}{e})-f(x_2)|\}\leq \sqrt{\frac{1}{e}-x_1}\leq \sqrt{x_2-x_1}$ 或 $\sqrt{x_2-\frac{1}{e}}\leq \sqrt{x_2-x_1}$. 故结论成立.
\end{solution}
\end{question}

\begin{question}[points=15]
\textbf{来源：}2024 年 新高考II卷，第 16 题\par
已知函数$f(x)=e^x-ax-a^3$.

(1)当$a=1$时，求曲线$y=f(x)$在点$(1,f(1))$处的切线方程；

(2)若$f(x)$有极小值，且极小值小于0，求$a$的取值范围.
\begin{solution}
\textbf{答案：}(1)$(e-1)x-y-1=0$
(2)$(1,+\infty)$

\textbf{分析：}(1)求导，结合导数的几何意义求切线方程；(2)解法一：求导，分析$a\le0$和$a>0$两种情况，利用导数判断单调性和极值，分析可得$a^2+\ln a-1>0$，构建函数解不等式即可；解法二：求导，可知$f'(x)=e^x-a$有零点，可得$a>0$，进而利用导数求$f(x)$的单调性和极值，分析可得$a^2+\ln a-1>0$，构建函数解不等式即可.

\textbf{详解：}（1）当$a=1$时，则$f(x)=e^x-x-1$，$f'(x)=e^x-1$，可得$f(1)=e-2$，$f'(1)=e-1$，即切点坐标为$(1,e-2)$，切线斜率$k=e-1$，所以切线方程为$y-(e-2)=(e-1)(x-1)$，即$(e-1)x-y-1=0$。
（2）解法一：因为$f(x)$的定义域为$\mathbb{R}$，且$f'(x)=e^x-a$。若$a\le0$，则$f'(x)\ge0$对任意$x\in\mathbb{R}$恒成立，可知$f(x)$在$\mathbb{R}$上单调递增，无极值，不合题意；若$a>0$，令$f'(x)>0$，解得$x>\ln a$；令$f'(x)<0$，解得$x<\ln a$；可知$f(x)$在$(-\infty,\ln a)$内单调递减，在$(\ln a,+\infty)$内单调递增，则$f(x)$有极小值$f(\ln a)=a-a\ln a-a^3$，无极大值。由题意可得：$f(\ln a)=a-a\ln a-a^3<0$，即$a^2+\ln a-1>0$。构建$g(a)=a^2+\ln a-1,a>0$，则$g'(a)=2a+\dfrac{1}{a}>0$，可知$g(a)$在$(0,+\infty)$内单调递增，且$g(1)=0$，不等式$a^2+\ln a-1>0$等价于$g(a)>g(1)$，解得$a>1$，所以$a$的取值范围为$(1,+\infty)$。
解法二：因为$f(x)$的定义域为$\mathbb{R}$，且$f'(x)=e^x-a$。若$f(x)$有极小值，则$f'(x)=e^x-a$有零点，令$f'(x)=e^x-a=0$，可得$e^x=a$，可知$y=e^x$与$y=a$有交点，则$a>0$。若$a>0$，令$f'(x)>0$，解得$x>\ln a$；令$f'(x)<0$，解得$x<\ln a$；可知$f(x)$在$(-\infty,\ln a)$内单调递减，在$(\ln a,+\infty)$内单调递增，则$f(x)$有极小值$f(\ln a)=a-a\ln a-a^3$，无极大值，符合题意。由题意可得：$f(\ln a)=a-a\ln a-a^3<0$，即$a^2+\ln a-1>0$。构建$g(a)=a^2+\ln a-1,a>0$，因为$y=a^2$和$y=\ln a-1$在$(0,+\infty)$内单调递增，可知$g(a)$在$(0,+\infty)$内单调递增，且$g(1)=0$，不等式$a^2+\ln a-1>0$等价于$g(a)>g(1)$，解得$a>1$，所以$a$的取值范围为$(1,+\infty)$.
\end{solution}
\end{question}

\begin{question}[points=5]
\textbf{来源：}2024 年 新高考I卷，第 13 题\par
若曲线 $y=e^x+x$ 在点 $(0,1)$ 处的切线也是曲线 $y=\ln(x+1)+a$ 的切线，则 $a=$ .
\fillin[{$\ln 2$}]
\begin{solution}
\textbf{答案：}$\ln 2$

\textbf{分析：}先求出曲线 $y=e^x+x$ 在 $(0,1)$ 的切线方程，再设曲线 $y=\ln(x+1)+a$ 的切点为 $(x_0,\ln(x_0+1)+a)$，求出 $y'$，利用公切线斜率相等求出 $x_0$，表示出切线方程，结合两切线方程相同即可求解.

\textbf{详解：}由 $y=e^x+x$ 得 $y'=e^x+1$，$y'|_{x=0}=e^0+1=2$，故曲线 $y=e^x+x$ 在 $(0,1)$ 处的切线方程为 $y=2x+1$；由 $y=\ln(x+1)+a$ 得 $y'=\dfrac{1}{x+1}$，设切线与曲线 $y=\ln(x+1)+a$ 相切的切点为 $(x_0,\ln(x_0+1)+a)$，由两曲线有公切线得 $\dfrac{1}{x_0+1}=2$，解得 $x_0=-\dfrac{1}{2}$，则切点为 $\left(-\dfrac{1}{2},a+\ln\dfrac{1}{2}\right)$，切线方程为 $y=2\left(x+\dfrac{1}{2}\right)+a+\ln\dfrac{1}{2}=2x+1+a-\ln 2$，根据两切线重合，所以 $a-\ln 2=0$，解得 $a=\ln 2$.
\end{solution}
\end{question}

\begin{question}[points=17]
\textbf{来源：}2024 年 新高考I卷，第 18 题\par
已知函数 $f(x)=\ln\dfrac{x}{2-x}+ax+b(x-1)^3$.

(1) 若 $b=0$，且 $f'(x)\ge 0$，求 $a$ 的最小值；

(2) 证明：曲线 $y=f(x)$ 是中心对称图形；

(3) 若 $f(x)>-2$ 当且仅当 $1<x<2$，求 $b$ 的取值范围.
\begin{solution}
\textbf{答案：}(1) $a_{\min}=-2$ (2) 证明见解析 (3) $b\ge -\dfrac{2}{3}$

\textbf{分析：}(1) 求出 $f'(x)$ 的最小值后根据 $f'(x)\ge 0$ 可求 $a$ 的最小值；(2) 设 $P(m,n)$ 为 $y=f(x)$ 图象上任意一点，可证 $P(m,n)$ 关于 $(1,a)$ 的对称点 $Q(2-m,2a-n)$ 也在函数的图像上，从而可证对称性；(3) 根据题设可判断 $f(1)=-2$ 即 $a=-2$，再根据 $f(x)>-2$ 在 $(1,2)$ 上恒成立可求得 $b\ge -\dfrac{2}{3}$.

\textbf{详解：}(1) $b=0$ 时，$f(x)=\ln\dfrac{x}{2-x}+ax$，其中 $x\in(0,2)$，则 $f'(x)=\dfrac{1}{x}+\dfrac{1}{2-x}+a=\dfrac{2}{x(2-x)}+a$，因为 $x(2-x)\le\left(\dfrac{2-x+x}{2}\right)^2=1$，当且仅当 $x=1$ 时等号成立，故 $f'(x)_{\min}=2+a$，而 $f'(x)\ge 0$ 成立，故 $a+2\ge 0$ 即 $a\ge -2$，所以 $a$ 的最小值为 $-2$.
(2) $f(x)$ 的定义域为 $(0,2)$，设 $P(m,n)$ 为 $y=f(x)$ 图象上任意一点，$P(m,n)$ 关于 $(1,a)$ 的对称点为 $Q(2-m,2a-n)$，因为 $P(m,n)$ 在 $y=f(x)$ 图象上，故 $n=\ln\dfrac{m}{2-m}+am+b(m-1)^3$，而 $f(2-m)=\ln\dfrac{2-m}{m}+a(2-m)+b(2-m-1)^3=-\left[\ln\dfrac{m}{2-m}+am+b(m-1)^3\right]+2a=-n+2a$，所以 $Q(2-m,2a-n)$ 也在 $y=f(x)$ 图象上，由 $P$ 的任意性可得 $y=f(x)$ 图象为中心对称图形，且对称中心为 $(1,a)$.
(3) 因为 $f(x)>-2$ 当且仅当 $1<x<2$，故 $x=1$ 为 $f(x)=-2$ 的一个解，所以 $f(1)=-2$ 即 $a=-2$，先考虑 $1<x<2$ 时，$f(x)>-2$ 恒成立. 此时 $f(x)>-2$ 即为 $\ln\dfrac{x}{2-x}+2(1-x)+b(x-1)^3>0$ 在 $(1,2)$ 上恒成立，设 $t=x-1\in(0,1)$，则 $\ln\dfrac{1+t}{1-t}-2t+bt^3>0$ 在 $(0,1)$ 上恒成立，设 $g(t)=\ln\dfrac{1+t}{1-t}-2t+bt^3, t\in(0,1)$，则 $g'(t)=\dfrac{2}{1-t^2}-2+3bt^2=\dfrac{2t^2+3bt^2(1-t^2)}{1-t^2}=\dfrac{t^2(2+3b-3bt^2)}{1-t^2}$，当 $b\ge 0$，$-3bt^2+2+3b\ge -3b+2+3b=2>0$，故 $g'(t)>0$ 恒成立，故 $g(t)$ 在 $(0,1)$ 上为增函数，故 $g(t)>g(0)=0$ 即 $f(x)>-2$ 在 $(1,2)$ 上恒成立. 当 $-\dfrac{2}{3}\le b<0$ 时，$-3bt^2+2+3b\ge 2+3b\ge 0$，故 $g'(t)\ge 0$ 恒成立，故 $g(t)$ 在 $(0,1)$ 上为增函数，故 $g(t)>g(0)=0$ 即 $f(x)>-2$ 在 $(1,2)$ 上恒成立. 当 $b<-\dfrac{2}{3}$，则当 $0<t<\sqrt{1+\dfrac{2}{3b}}$ 时，$g'(t)<0$，故在 $\left(0,\sqrt{1+\dfrac{2}{3b}}\right)$ 上 $g(t)$ 为减函数，故 $g(t)<g(0)=0$，不合题意. 综上，$f(x)>-2$ 在 $(1,2)$ 上恒成立时 $b\ge -\dfrac{2}{3}$.
\end{solution}
\end{question}

\section*{2025 年}

\begin{question}[points=17]
\textbf{来源：}2025 年 全国II卷，第 18 题\par
已知函数 $f(x)=\ln(1+x)-x+\frac{1}{2}x^2-kx^3$，其中 $0<k<\frac{1}{3}$.

（1）证明：$f(x)$ 在区间 $(0,+\infty)$ 存在唯一的极值点和唯一的零点；

（2）设 $x_1$，$x_2$ 分别为 $f(x)$ 在区间 $(0,+\infty)$ 的极值点和零点.

（i）设函数 $g(t)=f(x_1+t)-f(x_1-t)$，证明：$g(t)$ 在区间 $(0,x_1)$ 单调递减；

（ii）比较 $2x_1$ 与 $x_2$ 的大小，并证明你的结论.
\begin{solution}
\textbf{答案：}（1）证明见解析
（2）（i）证明见解析；（ii）$2x_1>x_2$，证明见解析

\textbf{分析：}（1）先由题意求得 $f'(x)=x^2\left(\frac{1}{1+x}-3k\right)$，接着构造函数 $h(x)=\frac{1}{1+x}-3k,x>0$，利用导数工具研究函数 $h(x)$ 的单调性和函数值情况，从而得到函数的单调性，进而得证函数 $f(x)$ 在区间 $(0,+\infty)$ 上存在唯一极值点；再结合 $f(0)=0$ 和 $x\to+\infty$ 时 $f(x)$ 的正负情况即可得证 $f(x)$ 在区间 $(0,+\infty)$ 上存在唯一零点；（2）（i）由（1）$x_1+1=\frac{1}{3k}$ 和 $g'(t)=f'(x_1+t)-f'(x_1-t)$ 结合（1）中所得导函数 $f'(x_1)$ 计算得到 $g'(t)=-t\left[\frac{(x_1+t)^2}{(x_1+t+1)(x_1+1)}+\frac{(x_1-t)^2}{(x_1-t+1)(x_1+1)}\right]$，再结合 $t\in(0,x_1)$ 得 $g'(t)<0$ 即可得证；（ii）由函数 $g(t)$ 在区间 $(0,x_1)$ 上单调递减得到 $0>f(2x_1)$，再结合 $f(x_2)=0$ 和函数 $f(x)$ 的单调性以及函数值的情况即可得证.

\textbf{详解：}（1）由题得 $f'(x)=\frac{1}{1+x}-1+x-3kx^2=\frac{x^2}{1+x}-3kx^2=x^2\left(\frac{1}{1+x}-3k\right)$，因为 $x\in(0,+\infty)$，所以 $x^2>0$，设 $h(x)=\frac{1}{1+x}-3k,x>0$，则 $h'(x)=-\frac{1}{(1+x)^2}<0$ 在 $(0,+\infty)$ 上恒成立，所以 $h(x)$ 在 $(0,+\infty)$ 上单调递减，$h(0)=1-3k>0$，令 $h(x_1)=0\Rightarrow x_1=\frac{1}{3k}-1$，所以当 $x\in(0,x_1)$ 时，$h(x)>0$，则 $f'(x)>0$；当 $x\in(x_1,+\infty)$ 时，$h(x)<0$，则 $f'(x)<0$，所以 $f(x)$ 在 $(0,x_1)$ 上单调递增，在 $(x_1,+\infty)$ 上单调递减，所以 $f(x)$ 在 $(0,+\infty)$ 上存在唯一极值点. 对函数 $y=\ln(1+x)-x$ 有 $y'=\frac{1}{1+x}-1=-\frac{x}{1+x}<0$ 在 $(0,+\infty)$ 上恒成立，所以 $y=\ln(1+x)-x$ 在 $(0,+\infty)$ 上单调递减，所以 $y=\ln(1+x)-x< y|_{x=0}=0$ 在 $(0,+\infty)$ 上恒成立，又因为 $f(0)=0$，$x\to+\infty$ 时 $\frac{1}{2}x^2-kx^3=\frac{1}{2}x^2(1-2kx)<0$，所以 $x\to+\infty$ 时 $f(x)<0$，所以存在唯一 $x_2\in(0,+\infty)$ 使得 $f(x_2)=0$，即 $f(x)$ 在 $(0,+\infty)$ 上存在唯一零点.
（2）（i）由（1）知 $x_1=\frac{1}{3k}-1$，则 $x_1+1=\frac{1}{3k}$，$f'(x)=x^2\left(\frac{1}{1+x}-3k\right)$，则 $g'(t)=f'(x_1+t)-f'(x_1-t)=(x_1+t)^2\left(\frac{1}{x_1+t+1}-3k\right)-(x_1-t)^2\left(\frac{1}{x_1-t+1}-3k\right)=(x_1+t)^2\left(\frac{1}{x_1+t+1}-\frac{1}{x_1+1}\right)-(x_1-t)^2\left(\frac{1}{x_1-t+1}-\frac{1}{x_1+1}\right)=\frac{-t(x_1+t)^2}{(x_1+t+1)(x_1+1)}-\frac{t(x_1-t)^2}{(x_1-t+1)(x_1+1)}=-t\left[\frac{(x_1+t)^2}{(x_1+t+1)(x_1+1)}+\frac{(x_1-t)^2}{(x_1-t+1)(x_1+1)}\right]$，因为 $t\in(0,x_1)$，所以 $x_1-t+1>0$，所以 $\frac{(x_1+t)^2}{(x_1+t+1)(x_1+1)}+\frac{(x_1-t)^2}{(x_1-t+1)(x_1+1)}>0$，所以 $g'(t)<0$，所以函数 $g(t)$ 在区间 $(0,x_1)$ 上单调递减.
（ii）$2x_1>x_2$，证明如下：由（i）知函数 $g(t)$ 在区间 $(0,x_1)$ 上单调递减，所以 $g(0)>g(x_1)$ 即 $0>f(2x_1)$，又 $f(x_2)=0$，由（1）可知 $f(x)$ 在 $(x_1,+\infty)$ 上单调递减，$x_2\in(x_1,+\infty)$，且对任意 $x\in(0,x_2)$，$f(x)>0$，所以 $2x_1>x_2$.
\end{solution}
\end{question}

\begin{question}[points=5]
\textbf{来源：}2025 年 全国I卷，第 12 题\par
若直线 $y = 2x + 5$ 是曲线 $y = e^x + x + a$ 的切线，则 $a =$ .
\fillin[{$4$}]
\begin{solution}
\textbf{答案：}$4$

\textbf{分析：}利用导数的几何性质与导数的四则运算求得切点，进而代入曲线方程即可得解.

\textbf{详解：}法一：对于 $y = e^x + x + a$，$y' = e^x + 1$，令 $e^x+1=2$，得 $x=0$，代入切线方程得 $y=5$，切点 $(0,5)$ 在曲线上，则 $5 = e^0 + 0 + a$，解得 $a=4$. 法二：略.
\end{solution}
\end{question}

\begin{question}[points=14]
\textbf{来源：}2025 年 上海卷，第 19 题\par
已知 $f(x) = x^2 - (m+2)x + m\ln x, m \in \mathbf{R}$．

(1) 若 $f(1) = 0$，求不等式 $f(x) \leq x^2 - 1$ 的解集；

(2) 若函数 $y = f(x)$ 满足在 $(0, +\infty)$ 上存在极大值，求 $m$ 的取值范围；
\begin{solution}
\textbf{答案：}(1) $[1, +\infty)$； (2) $m > 0$ 且 $m \neq 2$

\textbf{分析：}(1) 先求出 $m$，从而原不等式即为 $x + \ln x \geq 1$，构建新函数 $s(x) = x + \ln x, x > 0$，由该函数为增函数可求不等式的解； (2) 求出函数的导数，就 $m \leq 0, 0 < m < 2, m = 2, m > 2$ 分类讨论后可得参数的取值范围.

\textbf{详解：}(1) 因为 $f(1) = 0$，故 $1 - m - 2 + 0 = 0$，故 $m = -1$，故 $f(x) = x^2 - x - \ln x$，故 $f(x) \leq x^2 - 1$ 即为 $x + \ln x \geq 1$，设 $s(x) = x + \ln x, x > 0$，则 $s'(x) = 1 + \frac{1}{x} > 0$，故 $s(x)$ 在 $(0, +\infty)$ 上为增函数，而 $x + \ln x \geq 1$ 即为 $s(x) \geq s(1)$，故 $x \geq 1$，故原不等式的解为 $[1, +\infty)$．
(2) $f(x)$ 在 $(0, +\infty)$ 有极大值即为有极大值点．$f'(x) = 2x - (m+2) + \frac{m}{x} = \frac{2x^2 - (m+2)x + m}{x} = \frac{(2x - m)(x - 1)}{x}$．若 $m \leq 0$，则 $x \in (0,1)$ 时，$f'(x) < 0$，$x \in (1, +\infty)$ 时，$f'(x) > 0$，故 $x = 1$ 为 $f(x)$ 的极小值点，无极大值点，故舍；若 $0 < \frac{m}{2} < 1$ 即 $0 < m < 2$，则 $x \in \left(\frac{m}{2}, 1\right)$ 时，$f'(x) < 0$，$x \in \left(0, \frac{m}{2}\right) \cup (1, +\infty)$ 时，$f'(x) > 0$，故 $x = \frac{m}{2}$ 为 $f(x)$ 的极大值点，符合题设要求；若 $m = 2$，则 $x \in (0, +\infty)$ 时，$f'(x) \geq 0$，$f(x)$ 无极值点，舍；若 $\frac{m}{2} > 1$ 即 $m > 2$，则 $x \in \left(1, \frac{m}{2}\right)$ 时，$f'(x) < 0$，$x \in (0,1) \cup \left(\frac{m}{2}, +\infty\right)$ 时，$f'(x) > 0$，故 $x = 1$ 为 $f(x)$ 的极大值点，符合题设要求；综上，$m > 0$ 且 $m \neq 2$.
\end{solution}
\end{question}

\begin{question}[points=15]
\textbf{来源：}2025 年 天津卷，第 20 题\par
已知函数$f(x) = ax - (\ln x)^2$.

(1) $a = 1$时, 求$f(x)$在点$(1, f(1))$处的切线方程;

(2) $f(x)$有3个零点$x_1, x_2, x_3$且$x_1 < x_2 < x_3$.

(i) 求$a$的取值范围;

(ii) 证明: $(\ln x_2 - \ln x_1) \cdot \ln x_3 < \dfrac{4e}{e-1}$.
\begin{solution}

\textbf{分析：}(1)利用导数的几何意义，求导数值得斜率，由点斜式方程可得；(2)(i)令$f(x)=0$, 分离参数得$a = \dfrac{(\ln x)^2}{x}$, 作出函数$g(x) = \dfrac{(\ln x)^2}{x}$的图象，数形结合可得$a$范围；(ii)由(2)结合图象，可得$x_1, x_2, x_3$范围，整体换元$\ln x_1 = t_1$, $\ln x_2 = t_2$, $\ln x_3 = t_3$, 转化为$\begin{cases} ae^{t_1} = t_1^2 \\ ae^{t_2} = t_2^2 \\ ae^{t_3} = t_3^2 \end{cases}$, 结合②③可得$\begin{cases} \ln a + t_2 = 2\ln t_2 \\ \ln a + t_3 = 2\ln t_3 \end{cases}$, 两式作差, 利用对数平均不等式可得$t_2 t_3 < 4$, 再由$t_1 = ae^{t_1} < a$得$-t_1 < a$, 结合$a = \dfrac{t_3^2}{e^{t_3}}$减元处理, 再构造函数求最值, 放缩法可证明不等式.

\textbf{详解：}（1）当$a=1$时, $f(x) = x - (\ln x)^2$, $x > 0$, 则$f'(x) = 1 - \dfrac{2\ln x}{x}$, 则$f'(1) = 1$, 且$f(1) = 1$, 则切点$(1,1)$, 且切线的斜率为1, 故函数$f(x)$在点$(1, f(1))$处的切线方程为$y = x$.
（2）(i)令$f(x) = ax - (\ln x)^2 = 0$, $x > 0$, 得$a = \dfrac{(\ln x)^2}{x}$. 设$g(x) = \dfrac{(\ln x)^2}{x}$, $x > 0$, 则$g'(x) = \dfrac{2\ln x \cdot x - (\ln x)^2}{x^2} = \dfrac{\ln x (2 - \ln x)}{x^2}$. 由$g'(x) = 0$解得$x=1$或$e^2$, 其中$g(1)=0$, $g(e^2) = \dfrac{4}{e^2}$. 当$0 < x < 1$时, $g'(x) < 0$, $g(x)$在$(0,1)$上单调递减；当$1 < x < e^2$时, $g'(x) > 0$, $g(x)$在$(1, e^2)$上单调递增；当$x > e^2$时, $g'(x) < 0$, $g(x)$在$(e^2, +\infty)$上单调递减；且当$x \to 0$时, $g(x) \to +\infty$；当$x \to +\infty$时, $g(x) \to 0$. 要使函数$f(x)$有3个零点, 则方程$a = g(x)$在$(0, +\infty)$内有3个根, 即直线$y = a$与函数$g(x)$的图象有3个交点. 结合图象可知, $0 < a < \dfrac{4}{e^2}$. 故$a$的取值范围为$(0, \dfrac{4}{e^2})$.
(ii)由图象可知, $0 < x_1 < 1 < x_2 < e^2 < x_3$, 设$\ln x_1 = t_1$, $\ln x_2 = t_2$, $\ln x_3 = t_3$, 则$t_1 < 0 < t_2 < 2 < t_3$, 满足$\begin{cases} ae^{t_1} = t_1^2 & \text{①} \\ ae^{t_2} = t_2^2 & \text{②} \\ ae^{t_3} = t_3^2 & \text{③} \end{cases}$, 由②③可得$\begin{cases} \ln a + t_2 = 2\ln t_2 \\ \ln a + t_3 = 2\ln t_3 \end{cases}$, 两式作差可得$t_3 - t_2 = 2(\ln t_3 - \ln t_2)$, 则由对数均值不等式可得$2 = \dfrac{t_3 - t_2}{\ln t_3 - \ln t_2} > \sqrt{t_2 t_3}$, 则$t_2 t_3 < 4$. 故要证$(\ln x_2 - \ln x_1) \cdot \ln x_3 < \dfrac{4e}{e-1}$, 即证$t_2 t_3 - t_1 t_3 < \dfrac{4e}{e-1}$, 只需证$4 - t_1 t_3 \le \dfrac{4e}{e-1}$, 即证$-t_1 t_3 \le \dfrac{4}{e-1}$. 又因为$t_1 < 0$, $t_1 = ae^{t_1} < a$, 则$-t_1 = |t_1| < a$, 所以$-t_1 t_3 < a t_3 = \dfrac{t_3^2}{e^{t_3}} \cdot t_3 = \dfrac{t_3^2}{e^{t_3}}$, 故只需证$\dfrac{t^2}{e^{t/2}} \le \dfrac{4}{e-1}$ (其中$t = t_3 > 2$). 设函数$\varphi(t) = \dfrac{t^2}{e^{t/2}}$, $t > 2$, 则$\varphi'(t) = \dfrac{2t e^{t/2} - t^2 \cdot \frac{1}{2} e^{t/2}}{e^t} = \dfrac{t(4 - t)}{2e^{t/2}}$. 当$2 < t < 4$时, $\varphi'(t) > 0$, 则$\varphi(t)$在$(2,4)$上单调递增；当$t > 4$时, $\varphi'(t) < 0$, 则$\varphi(t)$在$(4, +\infty)$上单调递减. 故$\varphi(t)_{\max} = \varphi(4) = \dfrac{16}{e^2}$, 即$\varphi(t) \le \dfrac{16}{e^2}$. 而由$4e^2 - 16e + 16 = 4(e-2)^2 > 0$, 可知$\dfrac{16}{e^2} < \dfrac{4}{e-1}$成立, 故命题得证.
\end{solution}
\end{question}

